%% GENERATED by generate_papers.py -- edit the bank_*.py files, not this file.
\documentclass[11pt,a4paper]{article}
\usepackage[a4paper, top=20mm, bottom=20mm, left=18mm, right=18mm, headheight=14pt]{geometry}
\usepackage[T1]{fontenc}
\usepackage{lmodern}
\usepackage{amsmath, amssymb, mathtools}
\usepackage{braket}
\usepackage{array, tabularx, booktabs, multicol}
\usepackage{enumitem}
\usepackage[dvipsnames]{xcolor}
\usepackage[most]{tcolorbox}
\usepackage{needspace}
\usepackage{fancyhdr, lastpage}
\usepackage{pifont}
\usepackage{tikz}
\usepackage{graphicx}
\usepackage{eso-pic}
% SILICO FELLER watermark: official logo, centred, ~8 % opacity, behind all content, on every page
\AddToShipoutPictureBG{\AtPageCenter{\makebox(0,0){\tikz\node[opacity=0.08, inner sep=0pt]{\includegraphics[width=0.72\paperwidth]{../../../assets/branding/silico_feller_logo.png}};}}}
\usepackage[hidelinks]{hyperref}
\newcolumntype{L}{>{\raggedright\arraybackslash}X}
\definecolor{brand}{HTML}{1F4E79}
\definecolor{okgreen}{HTML}{1B7F3B}
\setlength{\parindent}{0pt}
\setlength{\parskip}{3pt}
\newcommand{\optlabel}[1]{\textbf{(#1)}}
\newcommand{\correctlabel}[1]{\textcolor{okgreen}{\textbf{(#1)}\,\ding{51}}}
\newenvironment{qitem}[1]{\begin{list}{}{\setlength{\leftmargin}{3.3em}\setlength{\labelwidth}{3.1em}%
  \setlength{\labelsep}{0.2em}\setlength{\topsep}{3pt}\setlength{\parsep}{3pt}}\item[\textbf{Q.#1}\hfill]}{\end{list}}
\newcommand{\qsep}{\par\vspace{2pt}{\color{black!15}\hrule}\vspace{4pt}}
\newcommand{\qmeta}[2]{{\footnotesize\color{black!60}\textsc{#1} \;$\cdot$\; #2}\par}
\newtcolorbox{passagebox}[1]{enhanced, breakable, colback=brand!4, colframe=brand!70, fonttitle=\bfseries,
  title={#1}, left=6pt, right=6pt, top=4pt, bottom=4pt}
\newtcolorbox{sectionbanner}{enhanced, colback=brand, colframe=brand, coltext=white, fontupper=\large\bfseries,
  left=8pt, right=8pt, top=6pt, bottom=6pt, sharp corners}
\newtcolorbox{infobox}[1]{enhanced, breakable, colback=black!3, colframe=black!45, fonttitle=\bfseries,
  title={#1}, left=6pt, right=6pt, top=4pt, bottom=4pt}
\newtcolorbox{solbox}[1]{enhanced, breakable, colback=okgreen!3, colframe=okgreen!60!black, fonttitle=\bfseries,
  title={Answer: (#1) \quad Detailed solution}, left=6pt, right=6pt, top=3pt, bottom=4pt, before skip=4pt, after skip=6pt}
\allowdisplaybreaks
\setlength{\emergencystretch}{3em}

\begin{document}
\pagestyle{fancy}\fancyhf{}
\fancyhead[L]{\small\color{black!70}Quantum Noise Simulation --- Assessment}
\fancyhead[R]{\small\color{black!70}\textbf{Answer Book}}
\fancyfoot[C]{\small Page \thepage\ of \pageref*{LastPage}}
\renewcommand{\headrulewidth}{0.4pt}

\thispagestyle{empty}
\begin{center}
{\color{brand}\rule{\linewidth}{1.2pt}}\\[6pt]
{\Huge\bfseries\color{brand} Quantum Noise Simulation}\\[4pt]
{\Large Assessment of noise models, channels, simulation and mitigation}\\[4pt]
{\large\bfseries ANSWER BOOK}\\[2pt]
{\color{brand}\rule{\linewidth}{1.2pt}}
\end{center}
This book accompanies the separate \emph{Question Paper} (80 questions, 3 hours, 240 marks). It contains
\begin{enumerate}[leftmargin=1.4em, itemsep=0pt, topsep=2pt]
\item the \textbf{answer key}, section by section, with the topic tested by every question (page~\pageref{key}),
\item a \textbf{topic-wise map} grouping the questions by theme and a \textbf{quick-marking grid} (page~\pageref{topicmap}), and
\item the \textbf{Answer Bank}: every question restated with the correct option marked \textcolor{okgreen}{\ding{51}} and a
      complete solution with derivations (from page~\pageref{bank}).
\end{enumerate}
\begin{infobox}{Marking scheme (every question)}
Each question has \textbf{four} options (A), (B), (C) and (D). \textbf{ONLY ONE} of these four options is correct.\\[2pt]
\begin{tabular}{@{}ll@{}}
\emph{Full marks}: & $+3$ if ONLY the correct option is chosen;\\
\emph{Zero marks}: & $\phantom{+}0$ if none of the options is chosen (question unanswered);\\
\emph{Negative marks}: & $-1$ in all other cases (wrong option, or more than one option chosen).
\end{tabular}
\end{infobox}

\begin{infobox}{How to evaluate}
For each section count $N_{\mathrm c}$ = number of correct answers and $N_{\mathrm w}$ = number of questions answered wrongly (or with more than one option marked). Section marks $= 3N_{\mathrm c} - N_{\mathrm w}$; unanswered questions score $0$. The total is out of $240$.
Every numerical answer was re-computed independently (NumPy Kraus-operator and density-matrix calculations, exact arithmetic where possible, and CUDA-Q density-matrix and stabilizer simulations, including the noise model of \texttt{Noise\_Simulation.ipynb}) by
\texttt{verify\_answers.py} in the source folder.
\end{infobox}
\begin{infobox}{Conventions (identical to the Question Paper)}
\footnotesize
\begin{itemize}[leftmargin=1.2em, itemsep=1pt, topsep=2pt]
\item Qubit 0 is written first: $\ket{q_0q_1\dots}$; two-qubit matrices use the basis order $\ket{00}, \ket{01}, \ket{10}, \ket{11}$. $X, Y, Z$ are the Pauli matrices,
      $H$ the Hadamard gate, $R_P(\theta) = e^{-i\theta P/2}$; CNOT flips the target when the control is 1.
\item A qubit state is $\rho = \frac12(I + r_xX + r_yY + r_zZ)$ with Bloch vector $\vec r$, $|\vec r| \le 1$; purity $\operatorname{Tr}\rho^2$. State fidelity with a pure target:
      $F = \braket{\psi|\rho|\psi}$. Measurement probabilities are the diagonal of $\rho$; $\langle Z\rangle = P(0) - P(1)$.
\item A channel is $\mathcal E(\rho) = \sum_kK_k\rho K_k^\dagger$ with $\sum_kK_k^\dagger K_k = I$. Bit flip: $\sqrt{1-p}\,I, \sqrt p\,X$; phase flip: $\sqrt{1-p}\,I, \sqrt p\,Z$;
      bit-phase flip: $\sqrt{1-p}\,I, \sqrt p\,Y$; depolarizing (CUDA-Q): $\sqrt{1-p}\,I, \sqrt{p/3}\,X, \sqrt{p/3}\,Y, \sqrt{p/3}\,Z$ (two-qubit: $\sqrt{1-p}\,II$ and the 15
      other Paulis with $p/15$); amplitude damping: $\operatorname{diag}(1, \sqrt{1-\gamma})$, $\sqrt\gamma\ket0\!\bra1$; phase damping: $\operatorname{diag}(1, \sqrt{1-\lambda})$, $\operatorname{diag}(0, \sqrt\lambda)$.
\item $T_1$: $\rho_{11}(t) = \rho_{11}(0)e^{-t/T_1}$; $T_2$: $|\rho_{01}(t)| = |\rho_{01}(0)|e^{-t/T_2}$; $\frac1{T_2} = \frac1{2T_1} + \frac1{T_\varphi}$. $\ln$ is the natural logarithm.
\item \textbf{Notebook} (\texttt{Noise\_Simulation.ipynb}, defaults): 15-qubit GHZ circuit ($H$ on qubit 0, CNOT chain), the chosen channel with probability $p$ after
      every \texttt{h}, \texttt{x}, \texttt{ry}, \texttt{rz} gate and a two-qubit channel after every CNOT; no measurement noise; 10\,000 shots, seed 42.
      \textbf{Previously recorded run}: CUDA-Q 0.16.0, CPU target \texttt{stim}, no GPU, depolarizing $p = 0.01$: population $0.8936$, $X$-basis parity $0.6910$,
      $(1-p)^G = 0.8601$. No physical quantum processor was used.
\end{itemize}
\end{infobox}
\newpage
\section*{Answer key}\label{key}
{\small Answer distribution: (A) 20, (B) 20, (C) 20, (D) 20.}

\needspace{14\baselineskip}
\subsection*{Section A --- Mathematical and Quantum Foundations}
{\small\renewcommand{\arraystretch}{1.15}
\begin{tabularx}{\linewidth}{@{}r c L @{\hspace{1em}} r c L@{}}
\toprule
\textbf{Q.} & \textbf{Ans.} & \textbf{Topic tested} & \textbf{Q.} & \textbf{Ans.} & \textbf{Topic tested}\\ \midrule
1 & \textbf{C} & Density matrix of a pure state & 11 & \textbf{C} & Partial trace; entanglement and mixedness \\
2 & \textbf{C} & Purity $\operatorname{Tr}\rho^2$ & 12 & \textbf{B} & Eigenvalues of a qubit $\rho$ \\
3 & \textbf{B} & Bloch vector from $\rho$ & 13 & \textbf{A} & Independent qubits: tensor product \\
4 & \textbf{D} & Validity of density matrices & 14 & \textbf{C} & Conjugation $X\rho X$ on the Bloch sphere \\
5 & \textbf{A} & Trace preservation of Kraus sets & 15 & \textbf{D} & Bloch vectors of states \\
6 & \textbf{C} & Random phases $\to$ dephasing & 16 & \textbf{D} & Hermitian, unitary, projector, Kraus \\
7 & \textbf{A} & Error-free probability $(1-p)^n$ & 17 & \textbf{A} & Mixed states inside the Bloch ball \\
8 & \textbf{B} & Shot-noise standard error & 18 & \textbf{D} & Density-matrix conditions; Born rule \\
9 & \textbf{A} & Coherent angle error & 19 & \textbf{D} & Tomography: Bloch vector, purity \\
10 & \textbf{B} & Fidelity with a mixed state & 20 & \textbf{B} & Tomography: fidelity, probabilities \\
\bottomrule
\end{tabularx}}

\needspace{14\baselineskip}
\subsection*{Section B --- Quantum Noise Fundamentals}
{\small\renewcommand{\arraystretch}{1.15}
\begin{tabularx}{\linewidth}{@{}r c L @{\hspace{1em}} r c L@{}}
\toprule
\textbf{Q.} & \textbf{Ans.} & \textbf{Topic tested} & \textbf{Q.} & \textbf{Ans.} & \textbf{Topic tested}\\ \midrule
21 & \textbf{D} & Fixed points of bit flip & 31 & \textbf{A} & Readout errors on GHZ \\
22 & \textbf{D} & Phase flip on $\ket+$ & 32 & \textbf{B} & Coherent vs incoherent errors \\
23 & \textbf{C} & Depolarizing flip probability & 33 & \textbf{B} & Two-qubit depolarizing marginal \\
24 & \textbf{A} & Amplitude damping on $\ket+$ & 34 & \textbf{A} & Phase errors and the $X$ basis \\
25 & \textbf{C} & Phase damping $\equiv$ phase flip & 35 & \textbf{B} & Kraus operators of channels \\
26 & \textbf{D} & $T_1$ half-life & 36 & \textbf{D} & Channels on the Bloch ball \\
27 & \textbf{D} & $T_2$ from $T_1$ and $T_\varphi$ & 37 & \textbf{B} & Non-unital amplitude damping \\
28 & \textbf{B} & Constraint $T_2 \le 2T_1$ & 38 & \textbf{C} & Complete dephasing vs unitary $Z$ \\
29 & \textbf{A} & Coherence-limited depth & 39 & \textbf{A} & $T_1$/$T_2$ decay of $\ket+$ \\
30 & \textbf{C} & Readout error probabilities & 40 & \textbf{C} & Fidelity and $T_\varphi$ \\
\bottomrule
\end{tabularx}}

\needspace{14\baselineskip}
\subsection*{Section C --- Noise Models and Simulation}
{\small\renewcommand{\arraystretch}{1.15}
\begin{tabularx}{\linewidth}{@{}r c L @{\hspace{1em}} r c L@{}}
\toprule
\textbf{Q.} & \textbf{Ans.} & \textbf{Topic tested} & \textbf{Q.} & \textbf{Ans.} & \textbf{Topic tested}\\ \midrule
41 & \textbf{B} & CUDA-Q depolarizing convention & 51 & \textbf{D} & Hellinger fidelity and TVD \\
42 & \textbf{A} & Constant identity gates and noise & 52 & \textbf{C} & Parity from $X$-basis counts \\
43 & \textbf{D} & CPTP check & 53 & \textbf{D} & Trajectories vs density matrix \\
44 & \textbf{B} & Applying the bit-flip channel & 54 & \textbf{D} & \texttt{add\_all\_qubit\_channel} \\
45 & \textbf{C} & Composing depolarizing channels & 55 & \textbf{A} & Noise-capable CUDA-Q targets \\
46 & \textbf{C} & Notebook: depolarizing $(1-p)^G$ & 56 & \textbf{C} & Figures of merit \\
47 & \textbf{A} & Notebook: bit-flip $(1-p)^{29}$ & 57 & \textbf{D} & $(1-p)^G$: bound, not exact value \\
48 & \textbf{A} & Density-matrix memory & 58 & \textbf{A} & Density-matrix path in CUDA-Q 0.16 \\
49 & \textbf{C} & What \texttt{stim} can simulate & 59 & \textbf{B} & Recorded run: leaked probability \\
50 & \textbf{B} & \texttt{as\_counts\_from\_rho} & 60 & \textbf{B} & Recorded run: population vs parity \\
\bottomrule
\end{tabularx}}

\needspace{14\baselineskip}
\subsection*{Section D --- Advanced Noise Analysis and Implementation}
{\small\renewcommand{\arraystretch}{1.15}
\begin{tabularx}{\linewidth}{@{}r c L @{\hspace{1em}} r c L@{}}
\toprule
\textbf{Q.} & \textbf{Ans.} & \textbf{Topic tested} & \textbf{Q.} & \textbf{Ans.} & \textbf{Topic tested}\\ \midrule
61 & \textbf{A} & Process vs average gate fidelity & 71 & \textbf{B} & Surface-code scaling \\
62 & \textbf{D} & Randomized benchmarking & 72 & \textbf{A} & Coherent vs stochastic accumulation \\
63 & \textbf{D} & Linear ZNE & 73 & \textbf{B} & Pauli twirling \\
64 & \textbf{C} & Richardson ZNE & 74 & \textbf{D} & Notebook: \texttt{auto} method choice \\
65 & \textbf{C} & Gate folding scale factor & 75 & \textbf{A} & Mitigation, DD, QEC \\
66 & \textbf{B} & Readout mitigation & 76 & \textbf{B} & Error taxonomy \\
67 & \textbf{D} & Limits of readout mitigation & 77 & \textbf{C} & Cost and nature of ZNE \\
68 & \textbf{D} & Ramsey vs Hahn echo & 78 & \textbf{B} & Readout and hardware in the notebook \\
69 & \textbf{C} & DD and Markovian noise & 79 & \textbf{A} & ZNE extrapolation values \\
70 & \textbf{A} & 3-qubit bit-flip code & 80 & \textbf{C} & Interpreting ZNE \\
\bottomrule
\end{tabularx}}

\needspace{18\baselineskip}
\subsection*{Topic-wise map of the questions}\label{topicmap}
{\small Every question appears under at least one theme; several appear under two.\par}
{\small\renewcommand{\arraystretch}{1.2}
\begin{tabularx}{\linewidth}{@{}>{\raggedright\arraybackslash}p{0.42\linewidth}L@{}}
\toprule
\textbf{Theme} & \textbf{Question numbers}\\ \midrule
Probability, statistics and sampling & 7, 8, 13, 30, 31, 51, 52, 59 \\
Density matrices, purity and the Bloch sphere & 1, 2, 3, 4, 12, 15, 17, 18, 19 \\
Linear operators, partial trace and fidelity & 10, 11, 14, 16, 20, 40, 61 \\
Phases, decoherence and dephasing & 6, 22, 25, 34, 38, 68, 69 \\
Pauli channels: bit, phase, bit-phase flip, depolarizing & 21, 23, 33, 35, 36, 41, 44, 45 \\
Amplitude and phase damping, $T_1$, $T_2$ & 5, 24, 25, 26, 27, 28, 29, 37, 39, 40 \\
Kraus operators and CPTP maps & 5, 16, 35, 43, 44, 45, 57 \\
Measurement and readout errors & 30, 31, 66, 67, 78 \\
Coherent vs stochastic errors & 9, 32, 72, 73, 76 \\
Noise models and simulator targets in CUDA-Q & 41, 42, 48, 49, 53, 54, 55, 58, 74 \\
\texttt{Noise\_Simulation.ipynb}: code and recorded results & 46, 47, 50, 56, 59, 60, 74, 78 \\
Fidelity, error rates and benchmarking & 51, 56, 59, 61, 62 \\
Error mitigation: ZNE, readout mitigation & 63, 64, 65, 66, 67, 77, 79, 80 \\
Dynamical decoupling and quantum error correction & 68, 69, 70, 71, 75 \\
\bottomrule
\end{tabularx}}

\subsection*{Quick-marking grid}
\begin{center}\renewcommand{\arraystretch}{1.3}
\begin{tabular}{|c|c|c|c|c|c|c|c|c|c|}\hline
1\,\textbf{C} & 2\,\textbf{C} & 3\,\textbf{B} & 4\,\textbf{D} & 5\,\textbf{A} & 6\,\textbf{C} & 7\,\textbf{A} & 8\,\textbf{B} & 9\,\textbf{A} & 10\,\textbf{B} \\
\hline
11\,\textbf{C} & 12\,\textbf{B} & 13\,\textbf{A} & 14\,\textbf{C} & 15\,\textbf{D} & 16\,\textbf{D} & 17\,\textbf{A} & 18\,\textbf{D} & 19\,\textbf{D} & 20\,\textbf{B} \\
\hline
21\,\textbf{D} & 22\,\textbf{D} & 23\,\textbf{C} & 24\,\textbf{A} & 25\,\textbf{C} & 26\,\textbf{D} & 27\,\textbf{D} & 28\,\textbf{B} & 29\,\textbf{A} & 30\,\textbf{C} \\
\hline
31\,\textbf{A} & 32\,\textbf{B} & 33\,\textbf{B} & 34\,\textbf{A} & 35\,\textbf{B} & 36\,\textbf{D} & 37\,\textbf{B} & 38\,\textbf{C} & 39\,\textbf{A} & 40\,\textbf{C} \\
\hline
41\,\textbf{B} & 42\,\textbf{A} & 43\,\textbf{D} & 44\,\textbf{B} & 45\,\textbf{C} & 46\,\textbf{C} & 47\,\textbf{A} & 48\,\textbf{A} & 49\,\textbf{C} & 50\,\textbf{B} \\
\hline
51\,\textbf{D} & 52\,\textbf{C} & 53\,\textbf{D} & 54\,\textbf{D} & 55\,\textbf{A} & 56\,\textbf{C} & 57\,\textbf{D} & 58\,\textbf{A} & 59\,\textbf{B} & 60\,\textbf{B} \\
\hline
61\,\textbf{A} & 62\,\textbf{D} & 63\,\textbf{D} & 64\,\textbf{C} & 65\,\textbf{C} & 66\,\textbf{B} & 67\,\textbf{D} & 68\,\textbf{D} & 69\,\textbf{C} & 70\,\textbf{A} \\
\hline
71\,\textbf{B} & 72\,\textbf{A} & 73\,\textbf{B} & 74\,\textbf{D} & 75\,\textbf{A} & 76\,\textbf{B} & 77\,\textbf{C} & 78\,\textbf{B} & 79\,\textbf{A} & 80\,\textbf{C} \\
\hline\end{tabular}\end{center}
\newpage
\section*{Answer Bank --- questions with detailed solutions}\label{bank}

\begin{sectionbanner}SECTION A \enspace$\cdot$\enspace Mathematical and Quantum Foundations \hfill Q.1--20\end{sectionbanner}

\needspace{13\baselineskip}
\begin{qitem}{1}
\qmeta{Single-correct MCQ}{Density matrix of a pure state}
The density matrix of the pure state $\ket\psi = \frac{1}{\sqrt2}(\ket0 + i\ket1)$ is
\par\smallskip\noindent\begin{tabularx}{\linewidth}{@{}*{4}{L}@{}}
\optlabel{A}\,$\frac12\begin{pmatrix}1 & 0\\ 0 & 1\end{pmatrix}$ & \optlabel{B}\,$\frac12\begin{pmatrix}1 & 1\\ 1 & 1\end{pmatrix}$ & \correctlabel{C}\,$\frac12\begin{pmatrix}1 & -i\\ i & 1\end{pmatrix}$ & \optlabel{D}\,$\frac12\begin{pmatrix}1 & i\\ -i & 1\end{pmatrix}$
\end{tabularx}
\end{qitem}
\begin{solbox}{C}
\textbf{\color{brand}Correct answer.} \textcolor{okgreen}{\textbf{(C)}} $\frac12\begin{pmatrix}1 & -i\\ i & 1\end{pmatrix}$

\textbf{\color{brand}Concepts required.} A pure state's density matrix is the outer product $\ket\psi\!\bra\psi$; the bra is the conjugate transpose of the ket.

\textbf{\color{brand}Relevant formulae.} $\rho_{ij} = \psi_i\psi_j^*$; for $\alpha\ket0 + \beta\ket1$: $\rho = \begin{psmallmatrix}|\alpha|^2 & \alpha\beta^*\\ \alpha^*\beta & |\beta|^2\end{psmallmatrix}$.

\textbf{\color{brand}Step-by-step explanation.} $\rho = \ket\psi\!\bra\psi$ with $\ket\psi = \frac{1}{\sqrt2}\begin{pmatrix}1\\ i\end{pmatrix}$ and $\bra\psi = \frac{1}{\sqrt2}(1,\ -i)$ (complex conjugate
transpose). Hence $\rho = \frac12\begin{pmatrix}1\cdot1 & 1\cdot(-i)\\ i\cdot1 & i\cdot(-i)\end{pmatrix} = \frac12\begin{pmatrix}1 & -i\\ i & 1\end{pmatrix}$.
\emph{Check}: Hermitian, trace 1, $\rho^2 = \rho$ (pure). Its Bloch vector is $(0, 1, 0)$.

\textbf{\color{brand}Why the other options are incorrect.}
\begin{itemize}[leftmargin=2.4em, itemsep=1pt, topsep=2pt]
\item[\textbf{(A)}] $I/2$ is the maximally mixed state; a pure state has $\rho^2 = \rho$ and non-zero coherence here.
\item[\textbf{(B)}] This is $\ket+\!\bra+$; the relative phase $i$ must appear in the off-diagonal elements.
\item[\textbf{(D)}] This forgets the complex conjugate in the bra: it is the density matrix of $\frac{1}{\sqrt2}(\ket0 - i\ket1)$.
\end{itemize}
\textbf{\color{brand}Key takeaway.} $\rho_{01} = \alpha\beta^*$: phases live in the off-diagonal elements.
\end{solbox}
\qsep

\needspace{13\baselineskip}
\begin{qitem}{2}
\qmeta{Single-correct MCQ}{Purity $\operatorname{Tr}\rho^2$}
The purity $\operatorname{Tr}\rho^2$ of $\rho = \begin{pmatrix}0.7 & 0.3\\ 0.3 & 0.3\end{pmatrix}$ is
\par\smallskip\noindent\begin{tabularx}{\linewidth}{@{}*{4}{L}@{}}
\optlabel{A}\,$0.58$ & \optlabel{B}\,$1$ & \correctlabel{C}\,$0.76$ & \optlabel{D}\,$0.52$
\end{tabularx}
\end{qitem}
\begin{solbox}{C}
\textbf{\color{brand}Correct answer.} \textcolor{okgreen}{\textbf{(C)}} $0.76$

\textbf{\color{brand}Concepts required.} Purity $\operatorname{Tr}\rho^2$ distinguishes pure (1) from mixed ($<1$) states; it equals the sum of squared moduli of all entries.

\textbf{\color{brand}Relevant formulae.} $\operatorname{Tr}\rho^2 = \sum_{ij}|\rho_{ij}|^2 = \frac{1 + |\vec r|^2}{2}$ for a qubit.

\textbf{\color{brand}Step-by-step explanation.} $\operatorname{Tr}\rho^2 = \sum_{ij}|\rho_{ij}|^2 = 0.7^2 + 0.3^2 + 0.3^2 + 0.3^2 = 0.49 + 0.09 + 0.09 + 0.09 = 0.76$.
\emph{Check} with the Bloch vector: $x = 2\operatorname{Re}\rho_{01} = 0.6$, $z = \rho_{00} - \rho_{11} = 0.4$, $|\vec r|^2 = 0.52$ and
$\operatorname{Tr}\rho^2 = \frac{1 + |\vec r|^2}{2} = 0.76$. Since $0.76 < 1$ the state is mixed.

\textbf{\color{brand}Why the other options are incorrect.}
\begin{itemize}[leftmargin=2.4em, itemsep=1pt, topsep=2pt]
\item[\textbf{(A)}] $0.58$ keeps only the diagonal ($0.49 + 0.09$) and forgets both off-diagonal terms.
\item[\textbf{(B)}] $1$ would mean a pure state, but $\det\rho = 0.12 \neq 0$.
\item[\textbf{(D)}] $0.52$ is $|\vec r|^2$, not the purity $\frac{1 + |\vec r|^2}{2}$.
\end{itemize}
\textbf{\color{brand}Key takeaway.} Purity = sum of $|\rho_{ij}|^2$; the off-diagonals count twice.
\end{solbox}
\qsep

\needspace{13\baselineskip}
\begin{qitem}{3}
\qmeta{Single-correct MCQ}{Bloch vector from $\rho$}
The Bloch vector $(r_x, r_y, r_z)$ of $\rho = \begin{pmatrix}0.8 & 0.2 - 0.1i\\ 0.2 + 0.1i & 0.2\end{pmatrix}$ is
\par\smallskip\noindent\begin{tabularx}{\linewidth}{@{}*{2}{L}@{}}
\optlabel{A}\,$(0.2,\ 0.1,\ 0.6)$ & \correctlabel{B}\,$(0.4,\ 0.2,\ 0.6)$ \\[3pt]
\optlabel{C}\,$(0.4,\ -0.2,\ 0.6)$ & \optlabel{D}\,$(0.4,\ 0.2,\ 0.8)$
\end{tabularx}
\end{qitem}
\begin{solbox}{B}
\textbf{\color{brand}Correct answer.} \textcolor{okgreen}{\textbf{(B)}} $(0.4,\ 0.2,\ 0.6)$

\textbf{\color{brand}Concepts required.} Every qubit density matrix is $\frac12(I + \vec r\cdot\vec\sigma)$; the Bloch vector components are Pauli expectation values.

\textbf{\color{brand}Relevant formulae.} $r_x = 2\operatorname{Re}\rho_{10}$, $r_y = 2\operatorname{Im}\rho_{10}$, $r_z = \rho_{00} - \rho_{11}$; $\rho_{01} = \frac{r_x - ir_y}{2}$.

\textbf{\color{brand}Step-by-step explanation.} Write $\rho = \frac12(I + r_xX + r_yY + r_zZ) = \frac12\begin{pmatrix}1 + r_z & r_x - ir_y\\ r_x + ir_y & 1 - r_z\end{pmatrix}$.
Then $r_z = \rho_{00} - \rho_{11} = 0.6$, and $\rho_{10} = \frac12(r_x + ir_y) = 0.2 + 0.1i$ gives $r_x = 0.4$, $r_y = 0.2$.
Equivalently $r_k = \operatorname{Tr}(\rho\sigma_k)$. $|\vec r|^2 = 0.56 < 1$: a valid mixed state.

\textbf{\color{brand}Why the other options are incorrect.}
\begin{itemize}[leftmargin=2.4em, itemsep=1pt, topsep=2pt]
\item[\textbf{(A)}] This forgets the factor 2 between off-diagonal entries and Bloch components.
\item[\textbf{(C)}] The sign of $r_y$ is read from $\rho_{01}$ instead of $\rho_{10}$: $\rho_{01} = \frac{r_x - ir_y}{2}$.
\item[\textbf{(D)}] $r_z = 0.8$ uses $\rho_{00}$ instead of $\rho_{00} - \rho_{11}$.
\end{itemize}
\textbf{\color{brand}Key takeaway.} $\vec r = (\operatorname{Tr}\rho X, \operatorname{Tr}\rho Y, \operatorname{Tr}\rho Z)$.
\end{solbox}
\qsep

\needspace{13\baselineskip}
\begin{qitem}{4}
\qmeta{Single-correct MCQ}{Validity of density matrices}
Which of the following matrices is \emph{not} a valid density matrix?
\par\smallskip\noindent\begin{tabularx}{\linewidth}{@{}*{4}{L}@{}}
\optlabel{A}\,$\begin{pmatrix}0.5 & 0.5\\ 0.5 & 0.5\end{pmatrix}$ & \optlabel{B}\,$\begin{pmatrix}0.5 & -0.3i\\ 0.3i & 0.5\end{pmatrix}$ & \optlabel{C}\,$\begin{pmatrix}0.9 & 0\\ 0 & 0.1\end{pmatrix}$ & \correctlabel{D}\,$\begin{pmatrix}0.6 & 0.5\\ 0.5 & 0.4\end{pmatrix}$
\end{tabularx}
\end{qitem}
\begin{solbox}{D}
\textbf{\color{brand}Correct answer.} \textcolor{okgreen}{\textbf{(D)}} $\begin{pmatrix}0.6 & 0.5\\ 0.5 & 0.4\end{pmatrix}$

\textbf{\color{brand}Concepts required.} Valid density matrices: Hermitian, unit trace, positive semi-definite (non-negative eigenvalues).

\textbf{\color{brand}Relevant formulae.} For $2\times2$ Hermitian matrices with trace 1: PSD $\Leftrightarrow \det\rho \ge 0 \Leftrightarrow |\vec r| \le 1$.

\textbf{\color{brand}Step-by-step explanation.} A density matrix must be Hermitian, have trace 1 and be positive semi-definite. All four are Hermitian with trace 1; for a
$2\times2$ matrix PSD $\Leftrightarrow$ $\det\rho \ge 0$. $\det\begin{psmallmatrix}0.6&0.5\\0.5&0.4\end{psmallmatrix} = 0.24 - 0.25 = -0.01 < 0$:
one eigenvalue is negative (a ``probability'' below zero). The others have determinants $0$ (pure $\ket+$), $0.09$ and $0.25 - 0.09 = 0.16$.
Equivalently the Bloch vector of the first has length $\sqrt{1^2 + 0.2^2} > 1$.

\textbf{\color{brand}Why the other options are incorrect.}
\begin{itemize}[leftmargin=2.4em, itemsep=1pt, topsep=2pt]
\item[\textbf{(A)}] $\ket+\!\bra+$: $\det = 0$, eigenvalues 1 and 0, a valid pure state.
\item[\textbf{(B)}] $\det = 0.25 - 0.09 = 0.16 > 0$; eigenvalues $0.8$ and $0.2$: valid.
\item[\textbf{(C)}] A diagonal matrix with entries $0.9, 0.1$ is a valid mixed state.
\end{itemize}
\textbf{\color{brand}Key takeaway.} Check the determinant (or $|\vec r| \le 1$), not just trace and Hermiticity.
\end{solbox}
\qsep

\needspace{13\baselineskip}
\begin{qitem}{5}
\qmeta{Single-correct MCQ}{Trace preservation of Kraus sets}
The operators $K_0 = \begin{pmatrix}1 & 0\\ 0 & 0.8\end{pmatrix}$ and $K_1 = \begin{pmatrix}0 & b\\ 0 & 0\end{pmatrix}$ ($b \ge 0$) form a
trace-preserving set of Kraus operators exactly when $b$ equals
\par\smallskip\noindent\begin{tabularx}{\linewidth}{@{}*{4}{L}@{}}
\correctlabel{A}\,$0.6$ & \optlabel{B}\,$0.36$ & \optlabel{C}\,$0.8$ & \optlabel{D}\,$0.4$
\end{tabularx}
\end{qitem}
\begin{solbox}{A}
\textbf{\color{brand}Correct answer.} \textcolor{okgreen}{\textbf{(A)}} $0.6$

\textbf{\color{brand}Concepts required.} Kraus operators of a channel must satisfy the completeness (trace-preservation) relation.

\textbf{\color{brand}Relevant formulae.} $\sum_kK_k^\dagger K_k = I$; amplitude damping $K_0 = \operatorname{diag}(1, \sqrt{1-\gamma})$, $K_1 = \sqrt\gamma\ket0\!\bra1$.

\textbf{\color{brand}Step-by-step explanation.} Trace preservation requires $\sum_kK_k^\dagger K_k = I$. $K_0^\dagger K_0 = \operatorname{diag}(1, 0.64)$ and
$K_1^\dagger K_1 = \begin{psmallmatrix}0&0\\b&0\end{psmallmatrix}\begin{psmallmatrix}0&b\\0&0\end{psmallmatrix} = \operatorname{diag}(0, b^2)$.
So $0.64 + b^2 = 1$, $b = 0.6$. This is amplitude damping with $\gamma = 0.36$: $K_0 = \operatorname{diag}(1, \sqrt{1-\gamma})$, $K_1 = \sqrt\gamma\ket0\!\bra1$.

\textbf{\color{brand}Why the other options are incorrect.}
\begin{itemize}[leftmargin=2.4em, itemsep=1pt, topsep=2pt]
\item[\textbf{(B)}] $b^2 = 0.36$ is required, so $b = 0.6$; choosing $b = 0.36$ confuses $b$ with $b^2$ and gives $0.64 + 0.13 = 0.77 \neq 1$.
\item[\textbf{(C)}] $b = 0.8$ gives $K_0^\dagger K_0 + K_1^\dagger K_1 = \operatorname{diag}(1, 0.64 + 0.64) = \operatorname{diag}(1, 1.28)$: probability would be created.
\item[\textbf{(D)}] $b = 0.4$ gives $\operatorname{diag}(1, 0.64 + 0.16) = \operatorname{diag}(1, 0.80)$: 20\,\% of the $\ket1$ population would disappear.
\end{itemize}
\textbf{\color{brand}Key takeaway.} Completeness: $K_0^\dagger K_0 + K_1^\dagger K_1 = I$ fixes $b^2 = 1 - 0.64$.
\end{solbox}
\qsep

\needspace{13\baselineskip}
\begin{qitem}{6}
\qmeta{Single-correct MCQ}{Random phases $\to$ dephasing}
A qubit is prepared in $\frac{1}{\sqrt2}(\ket0 + e^{i\varphi}\ket1)$ where the phase $\varphi$ is random and uniformly distributed in
$[-\frac\pi2, \frac\pi2]$. The average value of $\langle X\rangle$ is
\par\smallskip\noindent\begin{tabularx}{\linewidth}{@{}*{4}{L}@{}}
\optlabel{A}\,$\pi/4 \approx 0.79$ & \optlabel{B}\,$1/2$ & \correctlabel{C}\,$2/\pi \approx 0.64$ & \optlabel{D}\,$0$
\end{tabularx}
\end{qitem}
\begin{solbox}{C}
\textbf{\color{brand}Correct answer.} \textcolor{okgreen}{\textbf{(C)}} $2/\pi \approx 0.64$

\textbf{\color{brand}Concepts required.} Averaging over a random phase reduces coherence: the origin of dephasing.

\textbf{\color{brand}Relevant formulae.} $\langle X\rangle = \cos\varphi$; $\mathbb E\cos\varphi = \frac{1}{2\Delta}\int_{-\Delta}^{\Delta}\cos\varphi\,d\varphi = \frac{\sin\Delta}{\Delta}$.

\textbf{\color{brand}Step-by-step explanation.} For fixed $\varphi$, $\langle X\rangle = 2\operatorname{Re}(\alpha^*\beta) = \cos\varphi$. Averaging over the uniform distribution:
$\frac1\pi\int_{-\pi/2}^{\pi/2}\cos\varphi\,d\varphi = \frac{1}{\pi}[\sin\varphi]_{-\pi/2}^{\pi/2} = \frac2\pi \approx 0.637$.
The averaged state $\bar\rho$ has its off-diagonal element reduced from $\frac12$ to $\frac1\pi$: random phases are the origin of
\emph{dephasing}. Populations ($\langle Z\rangle = 0$) are unchanged.

\textbf{\color{brand}Why the other options are incorrect.}
\begin{itemize}[leftmargin=2.4em, itemsep=1pt, topsep=2pt]
\item[\textbf{(A)}] $\frac\pi4$ is the inverse of the correct ratio ($\frac{\Delta}{\sin\Delta}$ instead of $\frac{\sin\Delta}{\Delta}$ with $\Delta = \frac\pi2$) and exceeds the correct value.
\item[\textbf{(B)}] $\frac12$ is not the mean of $\cos\varphi$ on $[-\frac\pi2, \frac\pi2]$; it would be the mean of $\cos^2\varphi$ (the average of $P(+)$-like quantities), not of $\langle X\rangle$.
\item[\textbf{(D)}] $0$ would require the phase to be uniform over the whole circle $[-\pi, \pi]$, where positive and negative $\cos\varphi$ cancel; here $\cos\varphi \ge 0$ throughout.
\end{itemize}
\textbf{\color{brand}Key takeaway.} Random phases shrink $\langle X\rangle$, not the populations.
\end{solbox}
\qsep

\needspace{13\baselineskip}
\begin{qitem}{7}
\qmeta{Single-correct MCQ}{Error-free probability $(1-p)^n$}
A circuit has 20 gates; each fails independently with probability $p = 0.01$. The probability that \emph{at least one} gate fails is
\par\smallskip\noindent\begin{tabularx}{\linewidth}{@{}*{4}{L}@{}}
\correctlabel{A}\,$0.182$ & \optlabel{B}\,$0.818$ & \optlabel{C}\,$0.200$ & \optlabel{D}\,$0.165$
\end{tabularx}
\end{qitem}
\begin{solbox}{A}
\textbf{\color{brand}Correct answer.} \textcolor{okgreen}{\textbf{(A)}} $0.182$

\textbf{\color{brand}Concepts required.} Independent failures; complement rule.

\textbf{\color{brand}Relevant formulae.} $P(\ge1) = 1 - (1-p)^n$; $P(\text{exactly }k) = \binom nkp^k(1-p)^{n-k}$.

\textbf{\color{brand}Step-by-step explanation.} $P(\text{no failure}) = (1 - p)^{20} = 0.99^{20} = e^{20\ln0.99} = e^{-0.2010} = 0.818$. Hence
$P(\ge1) = 1 - 0.818 = 0.182$. The linear estimate $20p = 0.200$ over-counts runs with two or more failures;
$P(\text{exactly one}) = 20(0.01)(0.99)^{19} = 0.165$.

\textbf{\color{brand}Why the other options are incorrect.}
\begin{itemize}[leftmargin=2.4em, itemsep=1pt, topsep=2pt]
\item[\textbf{(B)}] $0.818$ is the probability of \emph{no} failure.
\item[\textbf{(C)}] $np = 0.2$ is a first-order approximation that over-counts.
\item[\textbf{(D)}] $0.165$ is the probability of exactly one failure.
\end{itemize}
\textbf{\color{brand}Key takeaway.} Error-free probability $(1-p)^G$ --- the notebook's lower bound.
\end{solbox}
\qsep

\needspace{13\baselineskip}
\begin{qitem}{8}
\qmeta{Single-correct MCQ}{Shot-noise standard error}
An outcome has true probability $0.9$. It is estimated from $10\,000$ shots. The standard error of the estimate is
\par\smallskip\noindent\begin{tabularx}{\linewidth}{@{}*{4}{L}@{}}
\optlabel{A}\,$0.03$ & \correctlabel{B}\,$0.003$ & \optlabel{C}\,$0.0009$ & \optlabel{D}\,$0.009$
\end{tabularx}
\end{qitem}
\begin{solbox}{B}
\textbf{\color{brand}Correct answer.} \textcolor{okgreen}{\textbf{(B)}} $0.003$

\textbf{\color{brand}Concepts required.} Binomial sampling error of an estimated probability.

\textbf{\color{brand}Relevant formulae.} $\mathrm{SE} = \sqrt{P(1-P)/S}$.

\textbf{\color{brand}Step-by-step explanation.} The count is binomial; the estimated frequency has $\mathrm{SE} = \sqrt{P(1 - P)/S} = \sqrt{0.9\cdot0.1/10\,000} = \sqrt{9\times10^{-6}} = 0.003$.
Shot noise is a \emph{statistical} error of the estimate, not a property of the quantum state: it shrinks like $1/\sqrt S$ even on a noiseless simulator.

\textbf{\color{brand}Why the other options are incorrect.}
\begin{itemize}[leftmargin=2.4em, itemsep=1pt, topsep=2pt]
\item[\textbf{(A)}] $0.03 = \sqrt{0.09/100}$ is the standard error for only 100 shots; with $10^4$ shots it is ten times smaller.
\item[\textbf{(C)}] $0.0009 = \frac{0.09}{100}$ is neither the variance $9\times10^{-6}$ nor its square root $0.003$.
\item[\textbf{(D)}] $0.009$ omits the square root: $\frac{0.9\times0.1}{10} = 0.009$ mixes the variance with a wrong divisor.
\end{itemize}
\textbf{\color{brand}Key takeaway.} Shot noise $\propto 1/\sqrt S$: 10\,000 shots give $\pm0.003$.
\end{solbox}
\qsep

\needspace{13\baselineskip}
\begin{qitem}{9}
\qmeta{Single-correct MCQ}{Coherent angle error}
An $X$ gate is implemented as $R_X(\pi + \varepsilon)$ instead of $R_X(\pi)$, with $\varepsilon = 0.1$ rad. Applied to $\ket0$, the
probability of measuring $0$ (an error) is closest to
\par\smallskip\noindent\begin{tabularx}{\linewidth}{@{}*{4}{L}@{}}
\correctlabel{A}\,$0.0025$ & \optlabel{B}\,$0.01$ & \optlabel{C}\,$0.1$ & \optlabel{D}\,$0.05$
\end{tabularx}
\end{qitem}
\begin{solbox}{A}
\textbf{\color{brand}Correct answer.} \textcolor{okgreen}{\textbf{(A)}} $0.0025$

\textbf{\color{brand}Concepts required.} Rotation gates and coherent angle errors.

\textbf{\color{brand}Relevant formulae.} $R_X(\theta) = \cos\frac\theta2I - i\sin\frac\theta2X$; $P(0) = \cos^2\frac{\pi + \varepsilon}{2} = \sin^2\frac\varepsilon2 \approx \frac{\varepsilon^2}{4}$.

\textbf{\color{brand}Step-by-step explanation.} $R_X(\theta)\ket0 = \cos\frac\theta2\ket0 - i\sin\frac\theta2\ket1$, so $P(0) = \cos^2\frac{\pi + \varepsilon}{2} = \sin^2\frac\varepsilon2$.
For $\varepsilon = 0.1$: $\sin^2(0.05) = 0.0025$. A small \emph{coherent} angle error gives an error probability quadratic in $\varepsilon$
(but such errors add up in amplitude when repeated --- see Section D).

\textbf{\color{brand}Why the other options are incorrect.}
\begin{itemize}[leftmargin=2.4em, itemsep=1pt, topsep=2pt]
\item[\textbf{(B)}] $0.01 = \varepsilon^2$ forgets the half angle.
\item[\textbf{(C)}] $0.1 = \varepsilon$ treats the angle as a probability.
\item[\textbf{(D)}] $0.05 = \varepsilon/2$ is an amplitude, not a probability.
\end{itemize}
\textbf{\color{brand}Key takeaway.} A coherent error $\varepsilon$ costs $\approx\varepsilon^2/4$ in probability per gate.
\end{solbox}
\qsep

\needspace{13\baselineskip}
\begin{qitem}{10}
\qmeta{Single-correct MCQ}{Fidelity with a mixed state}
The fidelity $F = \braket{+|\rho|+}$ between the target $\ket+$ and the mixed state $\rho = 0.8\ket+\!\bra+ + 0.2\,\frac I2$ is
\par\smallskip\noindent\begin{tabularx}{\linewidth}{@{}*{4}{L}@{}}
\optlabel{A}\,$0.6$ & \correctlabel{B}\,$0.9$ & \optlabel{C}\,$0.8$ & \optlabel{D}\,$0.95$
\end{tabularx}
\end{qitem}
\begin{solbox}{B}
\textbf{\color{brand}Correct answer.} \textcolor{okgreen}{\textbf{(B)}} $0.9$

\textbf{\color{brand}Concepts required.} Fidelity of a mixed state with a pure target is a linear function of $\rho$.

\textbf{\color{brand}Relevant formulae.} $F = \braket{\psi|\rho|\psi}$; $\braket{\psi|\frac I2|\psi} = \frac12$.

\textbf{\color{brand}Step-by-step explanation.} By linearity, $\braket{+|\rho|+} = 0.8\braket{+|+}\braket{+|+} + 0.2\cdot\frac12\braket{+|+} = 0.8 + 0.1 = 0.9$.
The maximally mixed part still overlaps the target with probability $\frac12$. (For a pure target, $F = \braket{\psi|\rho|\psi}$; some texts use the square root.)

\textbf{\color{brand}Why the other options are incorrect.}
\begin{itemize}[leftmargin=2.4em, itemsep=1pt, topsep=2pt]
\item[\textbf{(A)}] $0.6$ subtracts the noisy part instead of adding its overlap.
\item[\textbf{(C)}] $0.8$ ignores that the mixed part still overlaps the target.
\item[\textbf{(D)}] $0.95 \approx \sqrt{0.9}$ is the root fidelity convention.
\end{itemize}
\textbf{\color{brand}Key takeaway.} $F = (1-q) + q/d$ for a target mixed with white noise.
\end{solbox}
\qsep

\needspace{13\baselineskip}
\begin{qitem}{11}
\qmeta{Single-correct MCQ}{Partial trace; entanglement and mixedness}
For the two-qubit state $\ket\psi = \sqrt{0.8}\ket{00} + \sqrt{0.2}\ket{11}$, the purity $\operatorname{Tr}\rho_A^2$ of the reduced state of qubit A is
\par\smallskip\noindent\begin{tabularx}{\linewidth}{@{}*{4}{L}@{}}
\optlabel{A}\,$0.8$ & \optlabel{B}\,$0.5$ & \correctlabel{C}\,$0.68$ & \optlabel{D}\,$1$
\end{tabularx}
\end{qitem}
\begin{solbox}{C}
\textbf{\color{brand}Correct answer.} \textcolor{okgreen}{\textbf{(C)}} $0.68$

\textbf{\color{brand}Concepts required.} Partial trace; entanglement makes subsystems mixed.

\textbf{\color{brand}Relevant formulae.} $\rho_A = \operatorname{Tr}_B\ket\psi\!\bra\psi$; for $\sum_i\sqrt{\lambda_i}\ket{ii}$: $\rho_A = \operatorname{diag}(\lambda_i)$.

\textbf{\color{brand}Step-by-step explanation.} $\rho_A = \operatorname{Tr}_B\ket\psi\!\bra\psi = 0.8\ket0\!\bra0 + 0.2\ket1\!\bra1$ (the cross terms $\ket{00}\!\bra{11}$ have orthogonal B parts and vanish).
$\operatorname{Tr}\rho_A^2 = 0.64 + 0.04 = 0.68$. The global state is pure, but its parts are mixed because they are entangled ---
the same mechanism by which a qubit entangled with its environment decoheres.

\textbf{\color{brand}Why the other options are incorrect.}
\begin{itemize}[leftmargin=2.4em, itemsep=1pt, topsep=2pt]
\item[\textbf{(A)}] $0.8$ is the population $\rho_{A,00}$, not the purity $\operatorname{Tr}\rho_A^2 = 0.8^2 + 0.2^2$.
\item[\textbf{(B)}] $\frac12$ is the purity for the maximally entangled state $\frac{1}{\sqrt2}(\ket{00} + \ket{11})$, where $\rho_A = \frac I2$; here the weights are $0.8$ and $0.2$.
\item[\textbf{(D)}] $1$ would hold for a product state such as $\ket{00}$; entanglement with qubit B makes $\rho_A$ mixed even though the global state is pure.
\end{itemize}
\textbf{\color{brand}Key takeaway.} Decoherence = entanglement with an environment we trace out.
\end{solbox}
\qsep

\needspace{13\baselineskip}
\begin{qitem}{12}
\qmeta{Single-correct MCQ}{Eigenvalues of a qubit $\rho$}
The larger eigenvalue of $\rho = \begin{pmatrix}0.7 & 0.3\\ 0.3 & 0.3\end{pmatrix}$ is closest to
\par\smallskip\noindent\begin{tabularx}{\linewidth}{@{}*{4}{L}@{}}
\optlabel{A}\,$0.700$ & \correctlabel{B}\,$0.861$ & \optlabel{C}\,$1.000$ & \optlabel{D}\,$0.760$
\end{tabularx}
\end{qitem}
\begin{solbox}{B}
\textbf{\color{brand}Correct answer.} \textcolor{okgreen}{\textbf{(B)}} $0.861$

\textbf{\color{brand}Concepts required.} Eigenvalues of a qubit density matrix from its Bloch length.

\textbf{\color{brand}Relevant formulae.} $\lambda_\pm = \frac{1 \pm |\vec r|}{2}$; $\lambda_+\lambda_- = \det\rho$.

\textbf{\color{brand}Step-by-step explanation.} For a qubit, $\lambda_\pm = \frac{1 \pm |\vec r|}{2}$. Here $\vec r = (0.6, 0, 0.4)$, $|\vec r| = \sqrt{0.52} = 0.7211$, so
$\lambda_+ = 0.8606$, $\lambda_- = 0.1394$. \emph{Check}: $\lambda_+ + \lambda_- = 1 = \operatorname{Tr}\rho$ and $\lambda_+\lambda_- = 0.12 = \det\rho = 0.21 - 0.09$;
$\lambda_+^2 + \lambda_-^2 = 0.76 = \operatorname{Tr}\rho^2$.

\textbf{\color{brand}Why the other options are incorrect.}
\begin{itemize}[leftmargin=2.4em, itemsep=1pt, topsep=2pt]
\item[\textbf{(A)}] $0.700$ is the diagonal entry $\rho_{00}$; diagonal entries equal eigenvalues only when the off-diagonal elements vanish, and here $\rho_{01} = 0.3$.
\item[\textbf{(C)}] An eigenvalue $1$ would make the other eigenvalue $0$ and the state pure, but $\det\rho = 0.12 \neq 0$.
\item[\textbf{(D)}] $0.760$ is the purity $\lambda_+^2 + \lambda_-^2$, not an eigenvalue.
\end{itemize}
\textbf{\color{brand}Key takeaway.} Diagonal entries are not eigenvalues unless the coherence is zero.
\end{solbox}
\qsep

\needspace{13\baselineskip}
\begin{qitem}{13}
\qmeta{Single-correct MCQ}{Independent qubits: tensor product}
Two independent qubits are in the states $\rho_A = \operatorname{diag}(0.9, 0.1)$ and $\rho_B = \operatorname{diag}(0.8, 0.2)$. In the
state $\rho_A\otimes\rho_B$, the probability that both qubits give the \emph{same} outcome is
\par\smallskip\noindent\begin{tabularx}{\linewidth}{@{}*{4}{L}@{}}
\correctlabel{A}\,$0.74$ & \optlabel{B}\,$0.50$ & \optlabel{C}\,$0.26$ & \optlabel{D}\,$0.72$
\end{tabularx}
\end{qitem}
\begin{solbox}{A}
\textbf{\color{brand}Correct answer.} \textcolor{okgreen}{\textbf{(A)}} $0.74$

\textbf{\color{brand}Concepts required.} Product states of independent qubits; tensor products of diagonal matrices.

\textbf{\color{brand}Relevant formulae.} $P(ab) = P_A(a)P_B(b)$.

\textbf{\color{brand}Step-by-step explanation.} $\rho_A\otimes\rho_B = \operatorname{diag}(0.9\cdot0.8,\ 0.9\cdot0.2,\ 0.1\cdot0.8,\ 0.1\cdot0.2) = \operatorname{diag}(0.72, 0.18, 0.08, 0.02)$
in the order $\ket{00}, \ket{01}, \ket{10}, \ket{11}$. Same outcome: $0.72 + 0.02 = 0.74$.

\textbf{\color{brand}Why the other options are incorrect.}
\begin{itemize}[leftmargin=2.4em, itemsep=1pt, topsep=2pt]
\item[\textbf{(B)}] $0.5$ would hold for unbiased qubits; both qubits are strongly biased towards $0$.
\item[\textbf{(C)}] $0.26 = 0.18 + 0.08$ is the probability of \emph{different} outcomes ($01$ or $10$), the complement of the answer.
\item[\textbf{(D)}] $0.72$ counts only the outcome $00$; the outcome $11$ (probability $0.1\times0.2 = 0.02$) also has both qubits equal.
\end{itemize}
\textbf{\color{brand}Key takeaway.} Independent noise multiplies probabilities.
\end{solbox}
\qsep

\needspace{13\baselineskip}
\begin{qitem}{14}
\qmeta{Single-correct MCQ}{Conjugation $X\rho X$ on the Bloch sphere}
A qubit with Bloch vector $(r_x, r_y, r_z)$ is acted on by $\rho \mapsto X\rho X$. Its new Bloch vector is
\par\smallskip\noindent\begin{tabularx}{\linewidth}{@{}*{4}{L}@{}}
\optlabel{A}\,$(-r_x,\ -r_y,\ -r_z)$ & \optlabel{B}\,$(-r_x,\ r_y,\ r_z)$ & \correctlabel{C}\,$(r_x,\ -r_y,\ -r_z)$ & \optlabel{D}\,$(r_x,\ r_y,\ -r_z)$
\end{tabularx}
\end{qitem}
\begin{solbox}{C}
\textbf{\color{brand}Correct answer.} \textcolor{okgreen}{\textbf{(C)}} $(r_x,\ -r_y,\ -r_z)$

\textbf{\color{brand}Concepts required.} Conjugation by a Pauli is a $\pi$ rotation of the Bloch sphere.

\textbf{\color{brand}Relevant formulae.} $XXX = X$, $XYX = -Y$, $XZX = -Z$.

\textbf{\color{brand}Step-by-step explanation.} $X\rho X = \frac12(I + r_xXXX + r_yXYX + r_zXZX)$. Since $X$ commutes with itself and anticommutes with $Y$ and $Z$:
$XXX = X$, $XYX = -Y$, $XZX = -Z$. So $(r_x, r_y, r_z) \mapsto (r_x, -r_y, -r_z)$: a rotation by $\pi$ about the $x$-axis.
The bit-flip channel $(1 - p)\rho + pX\rho X$ therefore keeps $r_x$ and multiplies $r_y, r_z$ by $1 - 2p$.

\textbf{\color{brand}Why the other options are incorrect.}
\begin{itemize}[leftmargin=2.4em, itemsep=1pt, topsep=2pt]
\item[\textbf{(A)}] Inverting all three components is not a rotation (not a unitary map).
\item[\textbf{(B)}] This is conjugation by $Z$ or by a rotation about another axis, not by $X$.
\item[\textbf{(D)}] Only $r_z$ flipped would be a reflection, not a unitary rotation.
\end{itemize}
\textbf{\color{brand}Key takeaway.} $X\rho X$: keep $r_x$, flip $r_y$, $r_z$.
\end{solbox}
\qsep

\needspace{13\baselineskip}
\begin{qitem}{15}
\qmeta{Matching lists}{Bloch vectors of states}
Match each state (List-I) with its Bloch vector (List-II).
\par\smallskip\noindent{\renewcommand{\arraystretch}{1.35}%
\begin{tabularx}{\linewidth}{@{}c>{\raggedright\arraybackslash}p{0.33\linewidth}@{\hspace{1.2em}}cL@{}}
\toprule
\multicolumn{2}{@{}l}{\textbf{List-I}} & \multicolumn{2}{l@{}}{\textbf{List-II}} \\ \midrule
\textbf{(P)} & $\ket1$ & \textbf{(1)} & $(0,\ 0,\ 0.5)$ \\
\textbf{(Q)} & $\frac I2$ & \textbf{(2)} & $(0,\ 0,\ 0)$ \\
\textbf{(R)} & $\frac12\begin{psmallmatrix}1 & -i\\ i & 1\end{psmallmatrix}$ & \textbf{(3)} & $(0,\ 0,\ -1)$ \\
\textbf{(S)} & $\operatorname{diag}(0.75,\ 0.25)$ & \textbf{(4)} & $(0,\ 1,\ 0)$ \\
\bottomrule
\end{tabularx}}
\par\smallskip\noindent\begin{tabularx}{\linewidth}{@{}*{2}{L}@{}}
\optlabel{A}\,P$\to$3, Q$\to$1, R$\to$4, S$\to$2 & \optlabel{B}\,P$\to$3, Q$\to$2, R$\to$1, S$\to$4 \\[3pt]
\optlabel{C}\,P$\to$1, Q$\to$2, R$\to$4, S$\to$3 & \correctlabel{D}\,P$\to$3, Q$\to$2, R$\to$4, S$\to$1
\end{tabularx}
\end{qitem}
\begin{solbox}{D}
\textbf{\color{brand}Correct answer.} \textcolor{okgreen}{\textbf{(D)}} P$\to$3, Q$\to$2, R$\to$4, S$\to$1

\textbf{\color{brand}Concepts required.} Bloch vectors of pure, mixed and maximally mixed states.

\textbf{\color{brand}Relevant formulae.} $r_z = \rho_{00} - \rho_{11}$, $r_x + ir_y = 2\rho_{10}$.

\textbf{\color{brand}Step-by-step explanation.} Use $r_z = \rho_{00} - \rho_{11}$, $r_x + ir_y = 2\rho_{10}$.
\begin{itemize}\itemsep0pt
\item $\ket1$: south pole $(0, 0, -1)$ (3).
\item $I/2$: the centre of the ball (2) --- maximally mixed.
\item $\rho_{10} = \frac i2$: $r_y = 1$, i.e.\ $(0, 1, 0)$ (4), the pure state $\frac{1}{\sqrt2}(\ket0 + i\ket1)$.
\item $\operatorname{diag}(0.75, 0.25)$: $r_z = 0.5$ (1), a mixed state on the $z$-axis.
\end{itemize}
Hence P$\to$3, Q$\to$2, R$\to$4, S$\to$1.

\textbf{\color{brand}Why the other options are incorrect.}
\begin{itemize}[leftmargin=2.4em, itemsep=1pt, topsep=2pt]
\item[\textbf{(A)}] Q$\to$1, S$\to$2 confuses $I/2$ (centre) with a partially mixed state.
\item[\textbf{(B)}] R$\to$1, S$\to$4 gives a diagonal matrix a $y$-component, which needs an imaginary off-diagonal.
\item[\textbf{(C)}] P$\to$1, S$\to$3 swaps $\ket1$ and the mixed state: $\ket1$ is pure, length 1.
\end{itemize}
\textbf{\color{brand}Key takeaway.} Length = purity information; direction = which state.
\end{solbox}
\qsep

\needspace{13\baselineskip}
\begin{qitem}{16}
\qmeta{Matching lists}{Hermitian, unitary, projector, Kraus}
Match each matrix (List-I) with its property (List-II).
\par\smallskip\noindent{\renewcommand{\arraystretch}{1.35}%
\begin{tabularx}{\linewidth}{@{}c>{\raggedright\arraybackslash}p{0.33\linewidth}@{\hspace{1.2em}}cL@{}}
\toprule
\multicolumn{2}{@{}l}{\textbf{List-I}} & \multicolumn{2}{l@{}}{\textbf{List-II}} \\ \midrule
\textbf{(P)} & $\begin{psmallmatrix}0 & 1\\ 1 & 0\end{psmallmatrix}$ & \textbf{(1)} & unitary but not Hermitian \\
\textbf{(Q)} & $\frac12\begin{psmallmatrix}1 & 1\\ 1 & 1\end{psmallmatrix}$ & \textbf{(2)} & Hermitian and unitary \\
\textbf{(R)} & $\begin{psmallmatrix}1 & 0\\ 0 & i\end{psmallmatrix}$ & \textbf{(3)} & neither Hermitian nor unitary (a Kraus operator of amplitude damping, up to a factor) \\
\textbf{(S)} & $\begin{psmallmatrix}0 & 1\\ 0 & 0\end{psmallmatrix}$ & \textbf{(4)} & a rank-1 projector and a valid density matrix \\
\bottomrule
\end{tabularx}}
\par\smallskip\noindent\begin{tabularx}{\linewidth}{@{}*{2}{L}@{}}
\optlabel{A}\,P$\to$4, Q$\to$2, R$\to$1, S$\to$3 & \optlabel{B}\,P$\to$2, Q$\to$4, R$\to$3, S$\to$1 \\[3pt]
\optlabel{C}\,P$\to$1, Q$\to$4, R$\to$2, S$\to$3 & \correctlabel{D}\,P$\to$2, Q$\to$4, R$\to$1, S$\to$3
\end{tabularx}
\end{qitem}
\begin{solbox}{D}
\textbf{\color{brand}Correct answer.} \textcolor{okgreen}{\textbf{(D)}} P$\to$2, Q$\to$4, R$\to$1, S$\to$3

\textbf{\color{brand}Concepts required.} Hermitian, unitary, projector and density-matrix properties.

\textbf{\color{brand}Relevant formulae.} Hermitian $A = A^\dagger$; unitary $U^\dagger U = I$; projector $P^2 = P = P^\dagger$.

\textbf{\color{brand}Step-by-step explanation.} \begin{itemize}\itemsep0pt
\item $X$: $X^\dagger = X$ and $X^\dagger X = I$ (2); it has trace 0, so it is not a density matrix.
\item $\ket+\!\bra+$: Hermitian, $P^2 = P$, trace 1, PSD (4); not unitary ($\det = 0$).
\item $S = \operatorname{diag}(1, i)$: $S^\dagger S = I$ but $S^\dagger = \operatorname{diag}(1, -i) \neq S$ (1).
\item $\ket0\!\bra1$: nilpotent, $\neq$ its adjoint, not invertible (3); $\sqrt\gamma\ket0\!\bra1$ is the ``jump'' of amplitude damping.
\end{itemize}
Hence P$\to$2, Q$\to$4, R$\to$1, S$\to$3.

\textbf{\color{brand}Why the other options are incorrect.}
\begin{itemize}[leftmargin=2.4em, itemsep=1pt, topsep=2pt]
\item[\textbf{(A)}] P$\to$4 makes $X$ a density matrix, but $\operatorname{Tr}X = 0$ and $X^2 = I \neq X$ (not a projector).
\item[\textbf{(B)}] R$\to$3 calls $S$ non-unitary, but $S^\dagger S = \operatorname{diag}(1, -i\cdot i) = I$.
\item[\textbf{(C)}] P$\to$1 calls $X$ non-Hermitian, but $X^\dagger = X$; and R$\to$2 calls $S$ Hermitian, but $S^\dagger = \operatorname{diag}(1, -i) \neq S$.
\end{itemize}
\textbf{\color{brand}Key takeaway.} Kraus operators need not be Hermitian or unitary; only $\sum K^\dagger K = I$.
\end{solbox}
\qsep

\needspace{13\baselineskip}
\begin{qitem}{17}
\qmeta{Assertion--Reason}{Mixed states inside the Bloch ball}
\textbf{Assertion (A):} Every mixed single-qubit state lies strictly inside the Bloch ball ($|\vec r| < 1$).\\[1mm]
\textbf{Reason (R):} For $\rho = \frac12(I + \vec r\cdot\vec\sigma)$ one has $\operatorname{Tr}\rho^2 = \frac{1 + |\vec r|^2}{2}$, which equals 1 only when $|\vec r| = 1$.
\begin{itemize}[leftmargin=2.6em, labelsep=0.4em, itemsep=2pt, topsep=3pt]
\item[\correctlabel{A}] Both Assertion (A) and Reason (R) are true, and R is the correct explanation of A
\item[\optlabel{B}] Both A and R are true, but R is \emph{not} the correct explanation of A
\item[\optlabel{C}] A is true, but R is false
\item[\optlabel{D}] A is false, but R is true
\end{itemize}
\end{qitem}
\begin{solbox}{A}
\textbf{\color{brand}Correct answer.} \textcolor{okgreen}{\textbf{(A)}} Both Assertion (A) and Reason (R) are true, and R is the correct explanation of A

\textbf{\color{brand}Concepts required.} Purity and Bloch length.

\textbf{\color{brand}Relevant formulae.} $(\vec r\cdot\vec\sigma)^2 = |\vec r|^2I$; $\operatorname{Tr}\rho^2 = \frac{1 + |\vec r|^2}{2}$.

\textbf{\color{brand}Step-by-step explanation.} \textbf{R is true}: $\rho^2 = \frac14(I + 2\vec r\cdot\vec\sigma + (\vec r\cdot\vec\sigma)^2) = \frac14((1 + |\vec r|^2)I + 2\vec r\cdot\vec\sigma)$ because
$(\vec r\cdot\vec\sigma)^2 = |\vec r|^2I$; the trace is $\frac{1 + |\vec r|^2}{2}$. \textbf{A is true}: ``mixed'' means $\operatorname{Tr}\rho^2 < 1$,
which by R is exactly $|\vec r| < 1$. R is the reason for A. Answer: (A).

\textbf{\color{brand}Why the other options are incorrect.}
\begin{itemize}[leftmargin=2.4em, itemsep=1pt, topsep=2pt]
\item[\textbf{(B)}] requires that R does not explain A; but A follows directly from R, since ``mixed'' means $\operatorname{Tr}\rho^2 < 1$.
\item[\textbf{(C)}] requires R to be false, but $\operatorname{Tr}\rho^2 = \frac{1 + |\vec r|^2}{2}$ follows from $(\vec r\cdot\vec\sigma)^2 = |\vec r|^2I$ and $\operatorname{Tr}\sigma_k = 0$.
\item[\textbf{(D)}] requires A to be false, but no mixed state can reach the surface: $|\vec r| = 1$ forces $\operatorname{Tr}\rho^2 = 1$.
\end{itemize}
\textbf{\color{brand}Key takeaway.} Mixed $\Leftrightarrow$ inside the Bloch ball; noise shrinks Bloch vectors.
\end{solbox}
\qsep

\needspace{13\baselineskip}
\begin{qitem}{18}
\qmeta{Statement I / II}{Density-matrix conditions; Born rule}
\textbf{Statement I:} Every $2\times2$ Hermitian matrix with trace 1 is a valid density matrix.\\[1mm]
\textbf{Statement II:} The diagonal entries of a density matrix in the computational basis are the probabilities of the outcomes of a
computational-basis measurement.
\begin{itemize}[leftmargin=2.6em, labelsep=0.4em, itemsep=2pt, topsep=3pt]
\item[\optlabel{A}] Both Statement I and Statement II are true
\item[\optlabel{B}] Both Statement I and Statement II are false
\item[\optlabel{C}] Statement I is true, but Statement II is false
\item[\correctlabel{D}] Statement I is false, but Statement II is true
\end{itemize}
\end{qitem}
\begin{solbox}{D}
\textbf{\color{brand}Correct answer.} \textcolor{okgreen}{\textbf{(D)}} Statement I is false, but Statement II is true

\textbf{\color{brand}Concepts required.} Density-matrix conditions; Born rule for mixed states.

\textbf{\color{brand}Relevant formulae.} $P(k) = \operatorname{Tr}(\Pi_k\rho) = \rho_{kk}$.

\textbf{\color{brand}Step-by-step explanation.} \textbf{I is false}: positivity is also required; $\operatorname{diag}(1.2, -0.2)$ is Hermitian with trace 1 but has a negative eigenvalue.
\textbf{II is true}: $P(k) = \operatorname{Tr}(\ket k\!\bra k\rho) = \rho_{kk}$ (Born rule for mixed states). This is how the notebook turns an exact noisy
$\rho$ into outcome probabilities on the density-matrix path. Answer: (D).

\textbf{\color{brand}Why the other options are incorrect.}
\begin{itemize}[leftmargin=2.4em, itemsep=1pt, topsep=2pt]
\item[\textbf{(A)}] requires Statement I to be true, but I omits positivity: $\operatorname{diag}(1.2, -0.2)$ is Hermitian with trace 1 and is not a state.
\item[\textbf{(B)}] requires Statement II to be false, but $\rho_{kk} = \operatorname{Tr}(\ket k\!\bra k\rho)$ is exactly the Born-rule probability of outcome $k$.
\item[\textbf{(C)}] asserts I true and II false --- exactly the reverse of the correct judgement.
\end{itemize}
\textbf{\color{brand}Key takeaway.} Probabilities are on the diagonal; positivity is not automatic.
\end{solbox}
\qsep

\needspace{20\baselineskip}
\begin{passagebox}{Paragraph for Questions 19 and 20}
A qubit should have been prepared in $\ket+ = \frac{1}{\sqrt2}(\ket0 + \ket1)$. State tomography of the real (noisy)
preparation gives the density matrix
\[ \rho = \begin{pmatrix}0.55 & 0.35\\ 0.35 & 0.45\end{pmatrix}. \]
\end{passagebox}
\needspace{13\baselineskip}
\begin{qitem}{19}
\qmeta{Paragraph based}{Tomography: Bloch vector, purity}
The Bloch vector and the purity of the prepared state are
\par\smallskip\noindent\begin{tabularx}{\linewidth}{@{}*{2}{L}@{}}
\optlabel{A}\,$(0.7,\ 0,\ 0.1)$ and $0.50$ & \optlabel{B}\,$(0.7,\ 0,\ 0.55)$ and $0.85$ \\[3pt]
\optlabel{C}\,$(0.35,\ 0,\ 0.1)$ and $0.75$ & \correctlabel{D}\,$(0.7,\ 0,\ 0.1)$ and $0.75$
\end{tabularx}
\end{qitem}
\begin{solbox}{D}
\textbf{\color{brand}Correct answer.} \textcolor{okgreen}{\textbf{(D)}} $(0.7,\ 0,\ 0.1)$ and $0.75$

\textbf{\color{brand}Concepts required.} Bloch vector and purity of a measured density matrix.

\textbf{\color{brand}Relevant formulae.} $r_x = 2\operatorname{Re}\rho_{01}$, $r_z = \rho_{00} - \rho_{11}$; $\operatorname{Tr}\rho^2 = \frac{1 + |\vec r|^2}{2}$.

\textbf{\color{brand}Step-by-step explanation.} $r_x = 2\operatorname{Re}\rho_{01} = 0.7$, $r_y = 0$ (real matrix), $r_z = 0.55 - 0.45 = 0.1$.
$\operatorname{Tr}\rho^2 = 0.55^2 + 0.45^2 + 2(0.35)^2 = 0.3025 + 0.2025 + 0.245 = 0.75 = \frac{1 + 0.49 + 0.01}{2}$. The vector is shorter than 1
(mixed: decoherence) and tilted towards $\ket0$ (a population imbalance, e.g.\ relaxation).

\textbf{\color{brand}Why the other options are incorrect.}
\begin{itemize}[leftmargin=2.4em, itemsep=1pt, topsep=2pt]
\item[\textbf{(A)}] Purity $0.50$ belongs to the maximally mixed state $I/2$; here the Bloch vector has length $\sqrt{0.5} \approx 0.71$, so the purity is $0.75$.
\item[\textbf{(B)}] $r_z = 0.55$ uses $\rho_{00}$ instead of $\rho_{00} - \rho_{11}$, and $0.85$ is the fidelity with $\ket+$, not the purity.
\item[\textbf{(C)}] $r_x = 0.35$ forgets the factor 2 in $r_x = 2\operatorname{Re}\rho_{01}$ (and would give a different purity).
\end{itemize}
\textbf{\color{brand}Key takeaway.} Read $\vec r$ off $\rho$, then purity from $|\vec r|$.
\end{solbox}
\qsep

\needspace{13\baselineskip}
\begin{qitem}{20}
\qmeta{Paragraph based}{Tomography: fidelity, probabilities}
The fidelity $\braket{+|\rho|+}$ of the preparation, and the probability that a computational-basis measurement gives $1$, are
\par\smallskip\noindent\begin{tabularx}{\linewidth}{@{}*{4}{L}@{}}
\optlabel{A}\,$0.70$ and $0.45$ & \correctlabel{B}\,$0.85$ and $0.45$ & \optlabel{C}\,$0.92$ and $0.45$ & \optlabel{D}\,$0.85$ and $0.55$
\end{tabularx}
\end{qitem}
\begin{solbox}{B}
\textbf{\color{brand}Correct answer.} \textcolor{okgreen}{\textbf{(B)}} $0.85$ and $0.45$

\textbf{\color{brand}Concepts required.} Fidelity with $\ket+$ and computational-basis probabilities.

\textbf{\color{brand}Relevant formulae.} $\braket{+|\rho|+} = \frac{1 + r_x}{2}$; $P(1) = \rho_{11}$.

\textbf{\color{brand}Step-by-step explanation.} $\braket{+|\rho|+} = \frac12(\rho_{00} + \rho_{01} + \rho_{10} + \rho_{11}) = \frac12(0.55 + 0.35 + 0.35 + 0.45) = 0.85 = \frac{1 + r_x}{2}$.
$P(1) = \rho_{11} = 0.45$. ($0.92 = \sqrt{0.85}$ is the root fidelity used in some texts; $0.70$ is $r_x$ itself.)

\textbf{\color{brand}Why the other options are incorrect.}
\begin{itemize}[leftmargin=2.4em, itemsep=1pt, topsep=2pt]
\item[\textbf{(A)}] $0.70$ is the Bloch component $r_x$; the fidelity is $\frac{1 + r_x}{2} = 0.85$.
\item[\textbf{(C)}] $0.92 = \sqrt{0.85}$ is the ``root fidelity'' convention; the question defines $F = \braket{+|\rho|+}$.
\item[\textbf{(D)}] $0.55 = \rho_{00}$ is the probability of outcome 0; the probability of 1 is $\rho_{11} = 0.45$.
\end{itemize}
\textbf{\color{brand}Key takeaway.} $F(\ket+, \rho) = \frac{1 + r_x}{2}$.
\end{solbox}
\qsep

\newpage

\begin{sectionbanner}SECTION B \enspace$\cdot$\enspace Quantum Noise Fundamentals \hfill Q.21--40\end{sectionbanner}

\needspace{13\baselineskip}
\begin{qitem}{21}
\qmeta{Single-correct MCQ}{Fixed points of bit flip}
Which single-qubit state is left \emph{unchanged} by the bit-flip channel $\rho \mapsto (1 - p)\rho + pX\rho X$ for every $p$?
\par\smallskip\noindent\begin{tabularx}{\linewidth}{@{}*{2}{L}@{}}
\optlabel{A}\,$\frac{1}{\sqrt2}(\ket0 + i\ket1)$ & \optlabel{B}\,$\ket1$ \\[3pt]
\optlabel{C}\,$\ket0$ & \correctlabel{D}\,$\ket+$
\end{tabularx}
\end{qitem}
\begin{solbox}{D}
\textbf{\color{brand}Correct answer.} \textcolor{okgreen}{\textbf{(D)}} $\ket+$

\textbf{\color{brand}Concepts required.} Fixed points of a Pauli channel: eigenstates of the error operator.

\textbf{\color{brand}Relevant formulae.} Bit flip: $(r_x, r_y, r_z) \mapsto (r_x, (1-2p)r_y, (1-2p)r_z)$.

\textbf{\color{brand}Step-by-step explanation.} A state is invariant iff $X\rho X = \rho$, i.e.\ its Bloch vector satisfies $(r_x, -r_y, -r_z) = (r_x, r_y, r_z)$: $r_y = r_z = 0$.
$\ket+$ has $\vec r = (1, 0, 0)$ (indeed $X\ket+ = \ket+$). $\ket0$, $\ket1$ have $r_z = \pm1$ and become mixed; $\frac{1}{\sqrt2}(\ket0 + i\ket1)$ has
$r_y = 1$, which shrinks to $1 - 2p$. The eigenstates of the error operator are immune to that error.

\textbf{\color{brand}Why the other options are incorrect.}
\begin{itemize}[leftmargin=2.4em, itemsep=1pt, topsep=2pt]
\item[\textbf{(A)}] This state has $r_y = 1$, which shrinks to $1 - 2p$.
\item[\textbf{(B)}] $\ket1$ has $r_z = -1$, which shrinks as well.
\item[\textbf{(C)}] $\ket0$ has $r_z = 1$, which shrinks to $1 - 2p$.
\end{itemize}
\textbf{\color{brand}Key takeaway.} States along the error's axis are immune to it.
\end{solbox}
\qsep

\needspace{13\baselineskip}
\begin{qitem}{22}
\qmeta{Single-correct MCQ}{Phase flip on $\ket+$}
A phase-flip channel with $p = 0.1$ acts on $\ket+$. The state is then measured in the $X$ basis. The probability of the
outcome ``$+$'' is
\par\smallskip\noindent\begin{tabularx}{\linewidth}{@{}*{4}{L}@{}}
\optlabel{A}\,$0.95$ & \optlabel{B}\,$1.0$ & \optlabel{C}\,$0.8$ & \correctlabel{D}\,$0.9$
\end{tabularx}
\end{qitem}
\begin{solbox}{D}
\textbf{\color{brand}Correct answer.} \textcolor{okgreen}{\textbf{(D)}} $0.9$

\textbf{\color{brand}Concepts required.} Phase flips act like $\ket+ \leftrightarrow \ket-$.

\textbf{\color{brand}Relevant formulae.} $\rho' = (1-p)\rho + pZ\rho Z$; $\langle X\rangle \mapsto (1-2p)\langle X\rangle$; $P(+) = \frac{1 + \langle X\rangle}{2}$.

\textbf{\color{brand}Step-by-step explanation.} $\rho' = 0.9\ket+\!\bra+ + 0.1Z\ket+\!\bra+Z = 0.9\ket+\!\bra+ + 0.1\ket-\!\bra-$ (since $Z\ket+ = \ket-$). So $P(+) = 0.9$ and
$\langle X\rangle = 0.9 - 0.1 = 1 - 2p = 0.8$. In the $Z$ basis nothing changes ($P(0) = P(1) = \frac12$): phase errors are invisible there.

\textbf{\color{brand}Why the other options are incorrect.}
\begin{itemize}[leftmargin=2.4em, itemsep=1pt, topsep=2pt]
\item[\textbf{(A)}] $0.95$ uses $p/2$ as the error probability; a $Z$ applied with probability $p$ turns $\ket+$ into $\ket-$ with probability $p$.
\item[\textbf{(B)}] $1.0$ would mean the phase flip had no effect; that is true only for $Z$-basis measurements, not in the $X$ basis.
\item[\textbf{(C)}] $0.8 = 1 - 2p$ is the expectation value $\langle X\rangle = P(+) - P(-)$, not the probability $P(+)$.
\end{itemize}
\textbf{\color{brand}Key takeaway.} Phase errors appear only in a rotated basis.
\end{solbox}
\qsep

\needspace{13\baselineskip}
\begin{qitem}{23}
\qmeta{Single-correct MCQ}{Depolarizing flip probability}
A depolarizing channel $\rho \mapsto (1 - p)\rho + \frac p3(X\rho X + Y\rho Y + Z\rho Z)$ with $p = 0.03$ acts on $\ket0$. The probability of
then measuring $1$ is
\par\smallskip\noindent\begin{tabularx}{\linewidth}{@{}*{4}{L}@{}}
\optlabel{A}\,$0.03$ & \optlabel{B}\,$0.04$ & \correctlabel{C}\,$0.02$ & \optlabel{D}\,$0.01$
\end{tabularx}
\end{qitem}
\begin{solbox}{C}
\textbf{\color{brand}Correct answer.} \textcolor{okgreen}{\textbf{(C)}} $0.02$

\textbf{\color{brand}Concepts required.} Depolarizing channel in CUDA-Q's convention.

\textbf{\color{brand}Relevant formulae.} $\mathcal E(\rho) = (1-p)\rho + \frac p3(X\rho X + Y\rho Y + Z\rho Z)$, $\vec r \mapsto (1 - \frac{4p}{3})\vec r$.

\textbf{\color{brand}Step-by-step explanation.} $X$ and $Y$ flip $\ket0$ to $\ket1$ (up to a phase), $Z$ does not. So $P(1) = \frac p3 + \frac p3 = \frac{2p}{3} = 0.02$.
Equivalently the Bloch vector shrinks by $1 - \frac{4p}{3} = 0.96$: $r_z = 0.96$, $P(1) = \frac{1 - r_z}{2} = 0.02$.
(This is CUDA-Q's \texttt{DepolarizationChannel(p)} convention; the same channel is $\rho \mapsto (1 - \lambda)\rho + \lambda\frac I2$ with $\lambda = \frac{4p}{3}$.)

\textbf{\color{brand}Why the other options are incorrect.}
\begin{itemize}[leftmargin=2.4em, itemsep=1pt, topsep=2pt]
\item[\textbf{(A)}] $0.03 = p$ treats all three Paulis as flips, but $Z\ket0 = \ket0$ does not change the outcome.
\item[\textbf{(B)}] $0.04 = \frac{4p}{3}$ is the reduction $1 - r_z$ of the Bloch component, twice the flip probability $\frac{1 - r_z}{2}$.
\item[\textbf{(D)}] $0.01 = p/3$ counts only $X$; $Y = iXZ$ also maps $\ket0$ to $\ket1$ (up to a phase).
\end{itemize}
\textbf{\color{brand}Key takeaway.} Two of the three Paulis flip a $Z$ outcome: $P = \frac{2p}{3}$.
\end{solbox}
\qsep

\needspace{13\baselineskip}
\begin{qitem}{24}
\qmeta{Single-correct MCQ}{Amplitude damping on $\ket+$}
Amplitude damping with $\gamma = 0.36$ acts on $\ket+$. The resulting density matrix is
\par\smallskip\noindent\begin{tabularx}{\linewidth}{@{}*{4}{L}@{}}
\correctlabel{A}\,$\begin{pmatrix}0.68 & 0.40\\ 0.40 & 0.32\end{pmatrix}$ & \optlabel{B}\,$\begin{pmatrix}0.68 & 0.32\\ 0.32 & 0.32\end{pmatrix}$ & \optlabel{C}\,$\begin{pmatrix}0.50 & 0.40\\ 0.40 & 0.50\end{pmatrix}$ & \optlabel{D}\,$\begin{pmatrix}0.68 & 0.50\\ 0.50 & 0.32\end{pmatrix}$
\end{tabularx}
\end{qitem}
\begin{solbox}{A}
\textbf{\color{brand}Correct answer.} \textcolor{okgreen}{\textbf{(A)}} $\begin{pmatrix}0.68 & 0.40\\ 0.40 & 0.32\end{pmatrix}$

\textbf{\color{brand}Concepts required.} Amplitude damping: energy relaxation to $\ket0$.

\textbf{\color{brand}Relevant formulae.} $\rho_{11} \mapsto (1-\gamma)\rho_{11}$, $\rho_{00} \mapsto \rho_{00} + \gamma\rho_{11}$, $\rho_{01} \mapsto \sqrt{1-\gamma}\rho_{01}$.

\textbf{\color{brand}Step-by-step explanation.} With $K_0 = \operatorname{diag}(1, \sqrt{1-\gamma})$, $K_1 = \sqrt\gamma\ket0\!\bra1$:
$\rho_{11}' = (1 - \gamma)\rho_{11}$, $\rho_{00}' = \rho_{00} + \gamma\rho_{11}$, $\rho_{01}' = \sqrt{1 - \gamma}\,\rho_{01}$.
For $\ket+$ ($\rho_{ij} = \frac12$): $\rho_{11}' = 0.64\cdot0.5 = 0.32$, $\rho_{00}' = 0.5 + 0.18 = 0.68$, $\rho_{01}' = 0.8\cdot0.5 = 0.40$.
Populations move towards $\ket0$ and the coherence shrinks by $\sqrt{1 - \gamma}$, not by $1 - \gamma$.

\textbf{\color{brand}Why the other options are incorrect.}
\begin{itemize}[leftmargin=2.4em, itemsep=1pt, topsep=2pt]
\item[\textbf{(B)}] Scaling the coherence by $1 - \gamma = 0.64$ instead of $\sqrt{1-\gamma} = 0.8$ confuses the population decay with the amplitude decay.
\item[\textbf{(C)}] Leaving the populations at $0.5$ ignores the jump $\ket1 \to \ket0$; that would be phase damping, not amplitude damping.
\item[\textbf{(D)}] Keeping the coherence at $0.5$ ignores that the no-jump operator $\operatorname{diag}(1, \sqrt{1-\gamma})$ also shrinks $\rho_{01}$.
\end{itemize}
\textbf{\color{brand}Key takeaway.} $T_1$ decay also costs coherence, at half the rate.
\end{solbox}
\qsep

\needspace{13\baselineskip}
\begin{qitem}{25}
\qmeta{Single-correct MCQ}{Phase damping $\equiv$ phase flip}
Phase damping with parameter $\lambda$ multiplies the off-diagonal elements by $\sqrt{1 - \lambda}$; a phase-flip channel multiplies them by $1 - 2p$.
The phase-flip probability $p$ that has the same effect as phase damping with $\lambda = 0.36$ is
\par\smallskip\noindent\begin{tabularx}{\linewidth}{@{}*{4}{L}@{}}
\optlabel{A}\,$0.18$ & \optlabel{B}\,$0.36$ & \correctlabel{C}\,$0.10$ & \optlabel{D}\,$0.20$
\end{tabularx}
\end{qitem}
\begin{solbox}{C}
\textbf{\color{brand}Correct answer.} \textcolor{okgreen}{\textbf{(C)}} $0.10$

\textbf{\color{brand}Concepts required.} Different Kraus sets can represent the same channel.

\textbf{\color{brand}Relevant formulae.} Phase flip: coherence $\times(1-2p)$; phase damping: coherence $\times\sqrt{1-\lambda}$.

\textbf{\color{brand}Step-by-step explanation.} Both channels leave the populations unchanged, so they are equal iff they scale the coherence by the same factor:
$1 - 2p = \sqrt{1 - 0.36} = 0.8 \Rightarrow p = 0.1$. Phase damping and the phase-flip channel are the same channel in different Kraus
representations (Kraus operators are not unique).

\textbf{\color{brand}Why the other options are incorrect.}
\begin{itemize}[leftmargin=2.4em, itemsep=1pt, topsep=2pt]
\item[\textbf{(A)}] $0.18 = \lambda/2$ equates $2p$ with $\lambda$, which is only the first-order approximation of $1 - 2p = \sqrt{1-\lambda}$.
\item[\textbf{(B)}] $0.36$ sets $p = \lambda$; then $1 - 2p = 0.28 \neq 0.8$.
\item[\textbf{(D)}] $0.20 = 1 - \sqrt{1 - \lambda}$ forgets the factor 2 in the phase-flip coherence factor $1 - 2p$.
\end{itemize}
\textbf{\color{brand}Key takeaway.} Phase damping $\equiv$ phase flip with $1 - 2p = \sqrt{1-\lambda}$.
\end{solbox}
\qsep

\needspace{13\baselineskip}
\begin{qitem}{26}
\qmeta{Single-correct MCQ}{$T_1$ half-life}
A qubit with $T_1 = 50\,\mu$s starts in $\ket1$. After what time is $P(1) = \frac12$?
\par\smallskip\noindent\begin{tabularx}{\linewidth}{@{}*{4}{L}@{}}
\optlabel{A}\,$50\,\mu$s & \optlabel{B}\,$69.3\,\mu$s & \optlabel{C}\,$25\,\mu$s & \correctlabel{D}\,$34.7\,\mu$s
\end{tabularx}
\end{qitem}
\begin{solbox}{D}
\textbf{\color{brand}Correct answer.} \textcolor{okgreen}{\textbf{(D)}} $34.7\,\mu$s

\textbf{\color{brand}Concepts required.} Exponential $T_1$ decay and half-life.

\textbf{\color{brand}Relevant formulae.} $P_1(t) = e^{-t/T_1}$; $t_{1/2} = T_1\ln2$.

\textbf{\color{brand}Step-by-step explanation.} $P(1) = e^{-t/T_1} = \frac12 \Rightarrow t = T_1\ln2 = 50\times0.693 = 34.7\,\mu$s. At $t = T_1$ the population is $e^{-1} = 0.368$, not $\frac12$;
$69.3\,\mu$s would give $\frac14$.

\textbf{\color{brand}Why the other options are incorrect.}
\begin{itemize}[leftmargin=2.4em, itemsep=1pt, topsep=2pt]
\item[\textbf{(A)}] $50\,\mu$s $= T_1$ gives $P(1) = e^{-1} = 0.368$, not $\frac12$.
\item[\textbf{(B)}] $69.3\,\mu$s $= 2T_1\ln2$ is two half-lives: $P(1) = \frac14$.
\item[\textbf{(C)}] $25\,\mu$s $= T_1/2$ assumes linear decay; exponential decay gives $e^{-0.5} = 0.61$ at that time.
\end{itemize}
\textbf{\color{brand}Key takeaway.} Half-life $= 0.693\,T_1$.
\end{solbox}
\qsep

\needspace{13\baselineskip}
\begin{qitem}{27}
\qmeta{Single-correct MCQ}{$T_2$ from $T_1$ and $T_\varphi$}
A qubit has $T_1 = 100\,\mu$s and pure-dephasing time $T_\varphi = 200\,\mu$s. Using $\frac{1}{T_2} = \frac{1}{2T_1} + \frac{1}{T_\varphi}$, $T_2$ equals
\par\smallskip\noindent\begin{tabularx}{\linewidth}{@{}*{4}{L}@{}}
\optlabel{A}\,$150\,\mu$s & \optlabel{B}\,$66.7\,\mu$s & \optlabel{C}\,$200\,\mu$s & \correctlabel{D}\,$100\,\mu$s
\end{tabularx}
\end{qitem}
\begin{solbox}{D}
\textbf{\color{brand}Correct answer.} \textcolor{okgreen}{\textbf{(D)}} $100\,\mu$s

\textbf{\color{brand}Concepts required.} Relationship between $T_1$, $T_2$ and pure dephasing.

\textbf{\color{brand}Relevant formulae.} $\frac{1}{T_2} = \frac{1}{2T_1} + \frac{1}{T_\varphi}$.

\textbf{\color{brand}Step-by-step explanation.} $\frac{1}{T_2} = \frac{1}{200} + \frac{1}{200} = \frac{1}{100}\ \mu\text{s}^{-1}$, so $T_2 = 100\,\mu$s. Relaxation contributes only $\frac{1}{2T_1}$ because amplitude
damping multiplies the coherence by $\sqrt{1 - \gamma} \approx e^{-t/2T_1}$. Using $\frac1{T_1}$ instead would give $66.7\,\mu$s.

\textbf{\color{brand}Why the other options are incorrect.}
\begin{itemize}[leftmargin=2.4em, itemsep=1pt, topsep=2pt]
\item[\textbf{(A)}] $150\,\mu$s averages the two times; decoherence \emph{rates} add, times do not.
\item[\textbf{(B)}] $66.7\,\mu$s uses $\frac{1}{T_1}$ instead of $\frac{1}{2T_1}$ for the relaxation contribution to dephasing.
\item[\textbf{(C)}] $200\,\mu$s $= T_\varphi$ ignores the dephasing caused by relaxation itself.
\end{itemize}
\textbf{\color{brand}Key takeaway.} Add rates, not times; relaxation contributes half its rate to dephasing.
\end{solbox}
\qsep

\needspace{13\baselineskip}
\begin{qitem}{28}
\qmeta{Single-correct MCQ}{Constraint $T_2 \le 2T_1$}
Which pair $(T_1, T_2)$ is \emph{physically impossible} in the standard model?
\par\smallskip\noindent\begin{tabularx}{\linewidth}{@{}*{2}{L}@{}}
\optlabel{A}\,$T_1 = 100\,\mu$s, $T_2 = 30\,\mu$s & \correctlabel{B}\,$T_1 = 50\,\mu$s, $T_2 = 120\,\mu$s \\[3pt]
\optlabel{C}\,$T_1 = 50\,\mu$s, $T_2 = 100\,\mu$s & \optlabel{D}\,$T_1 = 80\,\mu$s, $T_2 = 80\,\mu$s
\end{tabularx}
\end{qitem}
\begin{solbox}{B}
\textbf{\color{brand}Correct answer.} \textcolor{okgreen}{\textbf{(B)}} $T_1 = 50\,\mu$s, $T_2 = 120\,\mu$s

\textbf{\color{brand}Concepts required.} Physical constraint $T_2 \le 2T_1$.

\textbf{\color{brand}Relevant formulae.} $\frac{1}{T_\varphi} = \frac1{T_2} - \frac{1}{2T_1} \ge 0$.

\textbf{\color{brand}Step-by-step explanation.} Since $\frac1{T_\varphi} = \frac1{T_2} - \frac1{2T_1} \ge 0$, every qubit satisfies $T_2 \le 2T_1$. $120 > 2\times50 = 100$ violates it. $T_2 = 2T_1$
(no pure dephasing) is the allowed limit; $T_2 < T_1$ is common in practice.

\textbf{\color{brand}Why the other options are incorrect.}
\begin{itemize}[leftmargin=2.4em, itemsep=1pt, topsep=2pt]
\item[\textbf{(A)}] $T_2 = 30\,\mu$s $< T_1 = 100\,\mu$s is common in real devices: strong pure dephasing.
\item[\textbf{(C)}] $T_2 = 2T_1 = 100\,\mu$s is the allowed limit of a qubit without pure dephasing ($T_\varphi = \infty$).
\item[\textbf{(D)}] $T_2 = T_1$ satisfies $T_2 \le 2T_1$ with $\frac1{T_\varphi} = \frac{1}{2T_1} > 0$.
\end{itemize}
\textbf{\color{brand}Key takeaway.} $T_2$ can never exceed $2T_1$.
\end{solbox}
\qsep

\needspace{13\baselineskip}
\begin{qitem}{29}
\qmeta{Single-correct MCQ}{Coherence-limited depth}
Gates take $50$ ns and the qubit has $T_2 = 100\,\mu$s. Requiring the coherence factor $e^{-t/T_2}$ to stay at least $0.9$, the largest number of
sequential gates is about
\par\smallskip\noindent\begin{tabularx}{\linewidth}{@{}*{4}{L}@{}}
\correctlabel{A}\,$210$ & \optlabel{B}\,$21$ & \optlabel{C}\,$2000$ & \optlabel{D}\,$105$
\end{tabularx}
\end{qitem}
\begin{solbox}{A}
\textbf{\color{brand}Correct answer.} \textcolor{okgreen}{\textbf{(A)}} $210$

\textbf{\color{brand}Concepts required.} Coherence-limited circuit depth.

\textbf{\color{brand}Relevant formulae.} $e^{-nt_g/T_2} \ge F \Rightarrow n \le \frac{T_2\ln(1/F)}{t_g}$.

\textbf{\color{brand}Step-by-step explanation.} $e^{-nt_g/T_2} \ge 0.9 \Leftrightarrow n \le \frac{T_2\ln(1/0.9)}{t_g} = \frac{100\,\mu\text{s}\times0.1054}{0.05\,\mu\text{s}} = 210.7$, so about $210$ gates.
$T_2/t_g = 2000$ is the number of gates in one $T_2$, when the coherence has already fallen to $e^{-1} = 0.37$.

\textbf{\color{brand}Why the other options are incorrect.}
\begin{itemize}[leftmargin=2.4em, itemsep=1pt, topsep=2pt]
\item[\textbf{(B)}] $21$ is ten times too small: it corresponds to a target coherence of about $0.99$ ($\ln(1/0.99) = 0.01$).
\item[\textbf{(C)}] $2000 = T_2/t_g$ is the number of gates in one $T_2$; by then the coherence has decayed to $e^{-1} = 0.37$, far below $0.9$.
\item[\textbf{(D)}] $105$ uses $\frac12\ln(1/0.9)$, i.e.\ treats the coherence as decaying twice as fast as $e^{-t/T_2}$.
\end{itemize}
\textbf{\color{brand}Key takeaway.} Depth budget $\approx (T_2/t_g)\ln(1/F)$.
\end{solbox}
\qsep

\needspace{13\baselineskip}
\begin{qitem}{30}
\qmeta{Single-correct MCQ}{Readout error probabilities}
A detector reads $1$ when the qubit is in $\ket0$ with probability $0.02$, and reads $0$ when it is in $\ket1$ with probability $0.05$. For a perfectly
prepared $\ket+$, the probability of reading $1$ is
\par\smallskip\noindent\begin{tabularx}{\linewidth}{@{}*{4}{L}@{}}
\optlabel{A}\,$0.515$ & \optlabel{B}\,$0.475$ & \correctlabel{C}\,$0.485$ & \optlabel{D}\,$0.500$
\end{tabularx}
\end{qitem}
\begin{solbox}{C}
\textbf{\color{brand}Correct answer.} \textcolor{okgreen}{\textbf{(C)}} $0.485$

\textbf{\color{brand}Concepts required.} Classical readout error as a confusion matrix.

\textbf{\color{brand}Relevant formulae.} $P(\text{read }1) = \sum_sP(s)P(1|s)$.

\textbf{\color{brand}Step-by-step explanation.} Law of total probability: $P(\text{read }1) = P(0)P(1|0) + P(1)P(1|1) = 0.5(0.02) + 0.5(0.95) = 0.01 + 0.475 = 0.485$.
Readout errors are classical (they act on the recorded bit) and are usually \emph{asymmetric}: $\ket1$ can relax during the measurement.

\textbf{\color{brand}Why the other options are incorrect.}
\begin{itemize}[leftmargin=2.4em, itemsep=1pt, topsep=2pt]
\item[\textbf{(A)}] $0.515 = 0.5(0.05) + 0.5(0.98)$ swaps the roles of the two error rates.
\item[\textbf{(B)}] $0.475 = 0.5\times0.95$ forgets the contribution $0.5\times0.02$ of $\ket0$ shots misread as 1.
\item[\textbf{(D)}] $0.5$ assumes a perfect detector; the two error rates are unequal, so they do not cancel.
\end{itemize}
\textbf{\color{brand}Key takeaway.} Readout errors mix outcome probabilities through a stochastic matrix.
\end{solbox}
\qsep

\needspace{13\baselineskip}
\begin{qitem}{31}
\qmeta{Single-correct MCQ}{Readout errors on GHZ}
A perfect 3-qubit GHZ state is measured with the readout errors of the previous question ($P(1|0) = 0.02$, $P(0|1) = 0.05$, independent per qubit).
The probability of observing $000$ or $111$ is about
\par\smallskip\noindent\begin{tabularx}{\linewidth}{@{}*{4}{L}@{}}
\correctlabel{A}\,$0.899$ & \optlabel{B}\,$0.950$ & \optlabel{C}\,$0.857$ & \optlabel{D}\,$0.941$
\end{tabularx}
\end{qitem}
\begin{solbox}{A}
\textbf{\color{brand}Correct answer.} \textcolor{okgreen}{\textbf{(A)}} $0.899$

\textbf{\color{brand}Concepts required.} Readout errors on several qubits multiply.

\textbf{\color{brand}Relevant formulae.} $P(\text{all correct}) = \prod_i(1 - e_i)$.

\textbf{\color{brand}Step-by-step explanation.} $P(000\text{ read correctly}) = 0.98^3 = 0.9412$, $P(111\text{ read correctly}) = 0.95^3 = 0.8574$. With probability $\frac12$ each:
$\frac12(0.9412 + 0.8574) = 0.8993$; the cross terms ($000 \to 111$: $0.02^3$, $111 \to 000$: $0.05^3$) add only $7\times10^{-5}$. Readout
errors grow with the number of measured qubits.

\textbf{\color{brand}Why the other options are incorrect.}
\begin{itemize}[leftmargin=2.4em, itemsep=1pt, topsep=2pt]
\item[\textbf{(B)}] $0.95$ is the single-qubit probability of reading $\ket1$ correctly; three qubits must all be read correctly.
\item[\textbf{(C)}] $0.857 = 0.95^3$ is the probability for the $111$ branch only.
\item[\textbf{(D)}] $0.941 = 0.98^3$ is the probability for the $000$ branch only, but GHZ gives $111$ half of the time.
\end{itemize}
\textbf{\color{brand}Key takeaway.} Readout fidelity decays exponentially with qubit number.
\end{solbox}
\qsep

\needspace{13\baselineskip}
\begin{qitem}{32}
\qmeta{Single-correct MCQ}{Coherent vs incoherent errors}
Which of the following is a \emph{coherent} (unitary) error rather than decoherence or a measurement error?
\begin{itemize}[leftmargin=2.6em, labelsep=0.4em, itemsep=2pt, topsep=3pt]
\item[\optlabel{A}] the discriminator labelling a $\ket1$ outcome as $0$
\item[\correctlabel{B}] a systematic over-rotation of every $X$ gate caused by a miscalibrated pulse amplitude
\item[\optlabel{C}] a random phase that differs from shot to shot and averages to dephasing
\item[\optlabel{D}] energy relaxation of $\ket1$ to $\ket0$ (a $T_1$ process)
\end{itemize}
\end{qitem}
\begin{solbox}{B}
\textbf{\color{brand}Correct answer.} \textcolor{okgreen}{\textbf{(B)}} a systematic over-rotation of every $X$ gate caused by a miscalibrated pulse amplitude

\textbf{\color{brand}Concepts required.} Coherent vs incoherent errors vs readout errors.

\textbf{\color{brand}Relevant formulae.} Coherent: $\rho \mapsto U_\varepsilon\rho U_\varepsilon^\dagger$; stochastic: $\sum_kp_kP_k\rho P_k$; readout: confusion matrix on outcomes.

\textbf{\color{brand}Step-by-step explanation.} A coherent error is a deterministic, unitary deviation $U_{\text{actual}} = U\,e^{-i\varepsilon H}$, identical in every run; it keeps the state pure
and its effects can add up in \emph{amplitude} (quadratically in probability). $T_1$ relaxation is non-unitary decoherence; a phase that is random from run
to run produces a mixed state (stochastic dephasing); a misassigned outcome is a classical readout error.

\textbf{\color{brand}Why the other options are incorrect.}
\begin{itemize}[leftmargin=2.4em, itemsep=1pt, topsep=2pt]
\item[\textbf{(A)}] Misassignment happens after the quantum measurement.
\item[\textbf{(C)}] A phase random from shot to shot is stochastic dephasing.
\item[\textbf{(D)}] $T_1$ relaxation is non-unitary (amplitude damping).
\end{itemize}
\textbf{\color{brand}Key takeaway.} Systematic and the same every time $\Rightarrow$ coherent.
\end{solbox}
\qsep

\needspace{13\baselineskip}
\begin{qitem}{33}
\qmeta{Single-correct MCQ}{Two-qubit depolarizing marginal}
After a CNOT, CUDA-Q's \texttt{Depolarization2(p)} applies each of the 15 non-identity two-qubit Paulis with probability $p/15$. For $p = 0.015$, the
probability that the outcome of a $Z$-basis measurement of the control qubit is flipped by this error is
\par\smallskip\noindent\begin{tabularx}{\linewidth}{@{}*{4}{L}@{}}
\optlabel{A}\,$0.010$ & \correctlabel{B}\,$0.008$ & \optlabel{C}\,$0.0075$ & \optlabel{D}\,$0.015$
\end{tabularx}
\end{qitem}
\begin{solbox}{B}
\textbf{\color{brand}Correct answer.} \textcolor{okgreen}{\textbf{(B)}} $0.008$

\textbf{\color{brand}Concepts required.} Two-qubit depolarizing channel and its marginal effect.

\textbf{\color{brand}Relevant formulae.} $\mathcal E_2(\rho) = (1-p)\rho + \frac{p}{15}\sum_{P\neq II}P\rho P$.

\textbf{\color{brand}Step-by-step explanation.} The $Z$ outcome of the control flips iff the Pauli on the control is $X$ or $Y$. Of the 15 non-identity pairs $P_c\otimes P_t$, those with
$P_c \in \{X, Y\}$ are $2\times4 = 8$. Probability $= 8\cdot\frac{p}{15} = \frac{8(0.015)}{15} = 0.008$.

\textbf{\color{brand}Why the other options are incorrect.}
\begin{itemize}[leftmargin=2.4em, itemsep=1pt, topsep=2pt]
\item[\textbf{(A)}] $0.010 = \frac{2p}{3}$ is the flip probability of the \emph{single-qubit} depolarizing channel, not the marginal of the two-qubit one ($\frac{8}{15}p$).
\item[\textbf{(C)}] $0.0075 = \frac p2$ assumes half of the Paulis flip the control; the correct fraction is $\frac{8}{15}$.
\item[\textbf{(D)}] $0.015 = p$ counts all 15 Paulis as flipping the control, but $IX, IY, IZ, ZI, \dots$ ($7$ of them) leave its $Z$ outcome unchanged.
\end{itemize}
\textbf{\color{brand}Key takeaway.} Count the Paulis that anticommute with the measured observable.
\end{solbox}
\qsep

\needspace{13\baselineskip}
\begin{qitem}{34}
\qmeta{Single-correct MCQ}{Phase errors and the $X$ basis}
A 15-qubit GHZ state suffers only phase-flip noise. Which measurement reveals the errors?
\begin{itemize}[leftmargin=2.6em, labelsep=0.4em, itemsep=2pt, topsep=3pt]
\item[\correctlabel{A}] the parity $\langle X^{\otimes15}\rangle$ measured after a Hadamard on every qubit
\item[\optlabel{B}] the probability $P(0\cdots0)$ alone
\item[\optlabel{C}] the number of distinct outcomes in the $Z$ basis
\item[\optlabel{D}] the $Z$-basis population $P(0\cdots0) + P(1\cdots1)$
\end{itemize}
\end{qitem}
\begin{solbox}{A}
\textbf{\color{brand}Correct answer.} \textcolor{okgreen}{\textbf{(A)}} the parity $\langle X^{\otimes15}\rangle$ measured after a Hadamard on every qubit

\textbf{\color{brand}Concepts required.} Which errors each measurement basis can see.

\textbf{\color{brand}Relevant formulae.} $[Z, Z_i] = 0$; $Z_i(\ket{0\cdots0} + \ket{1\cdots1}) = \ket{0\cdots0} - \ket{1\cdots1}$; $\langle X^{\otimes N}\rangle \to -\langle X^{\otimes N}\rangle$.

\textbf{\color{brand}Step-by-step explanation.} $Z$ errors commute with $Z$-basis measurements, so all $Z$-basis statistics are unchanged: the population stays 1 and only the outcomes $0\cdots0$,
$1\cdots1$ occur. A single $Z$ error turns $\ket{0\cdots0} + \ket{1\cdots1}$ into $\ket{0\cdots0} - \ket{1\cdots1}$, which flips the sign of $\langle X^{\otimes N}\rangle$.
This is why \texttt{Noise\_Simulation.ipynb} also measures in the $X$ basis (\texttt{MEASURE\_X\_BASIS = True}).

\textbf{\color{brand}Why the other options are incorrect.}
\begin{itemize}[leftmargin=2.4em, itemsep=1pt, topsep=2pt]
\item[\textbf{(B)}] $P(0\cdots0)$ stays $\frac12$; $\ket{0\cdots0} \pm \ket{1\cdots1}$ have identical $Z$ statistics.
\item[\textbf{(C)}] The number of distinct $Z$-basis outcomes stays 2 ($0\cdots0$ and $1\cdots1$), whatever the phase errors.
\item[\textbf{(D)}] The $Z$-basis population stays exactly 1: $Z$ errors only change phases.
\end{itemize}
\textbf{\color{brand}Key takeaway.} Measure in a second basis to see phase errors.
\end{solbox}
\qsep

\needspace{13\baselineskip}
\begin{qitem}{35}
\qmeta{Matching lists}{Kraus operators of channels}
Match each channel (List-I) with a set of Kraus operators (List-II).
\par\smallskip\noindent{\renewcommand{\arraystretch}{1.35}%
\begin{tabularx}{\linewidth}{@{}c>{\raggedright\arraybackslash}p{0.33\linewidth}@{\hspace{1.2em}}cL@{}}
\toprule
\multicolumn{2}{@{}l}{\textbf{List-I}} & \multicolumn{2}{l@{}}{\textbf{List-II}} \\ \midrule
\textbf{(P)} & bit flip & \textbf{(1)} & $\operatorname{diag}(1, \sqrt{1-\lambda})$, $\operatorname{diag}(0, \sqrt\lambda)$ \\
\textbf{(Q)} & phase flip & \textbf{(2)} & $\sqrt{1-p}\,I$, $\sqrt p\,Z$ \\
\textbf{(R)} & amplitude damping & \textbf{(3)} & $\begin{psmallmatrix}1&0\\0&\sqrt{1-\gamma}\end{psmallmatrix}$, $\begin{psmallmatrix}0&\sqrt\gamma\\0&0\end{psmallmatrix}$ \\
\textbf{(S)} & phase damping & \textbf{(4)} & $\sqrt{1-p}\,I$, $\sqrt p\,X$ \\
\bottomrule
\end{tabularx}}
\par\smallskip\noindent\begin{tabularx}{\linewidth}{@{}*{2}{L}@{}}
\optlabel{A}\,P$\to$4, Q$\to$1, R$\to$3, S$\to$2 & \correctlabel{B}\,P$\to$4, Q$\to$2, R$\to$3, S$\to$1 \\[3pt]
\optlabel{C}\,P$\to$2, Q$\to$4, R$\to$3, S$\to$1 & \optlabel{D}\,P$\to$4, Q$\to$2, R$\to$1, S$\to$3
\end{tabularx}
\end{qitem}
\begin{solbox}{B}
\textbf{\color{brand}Correct answer.} \textcolor{okgreen}{\textbf{(B)}} P$\to$4, Q$\to$2, R$\to$3, S$\to$1

\textbf{\color{brand}Concepts required.} Kraus representations of the standard channels.

\textbf{\color{brand}Relevant formulae.} Completeness $\sum K^\dagger K = I$ for each set.

\textbf{\color{brand}Step-by-step explanation.} \begin{itemize}\itemsep0pt
\item bit flip: apply $X$ with probability $p$ (4);
\item phase flip: apply $Z$ with probability $p$ (2);
\item amplitude damping: the ``jump'' $\sqrt\gamma\ket0\!\bra1$ and the no-jump $K_0$ (3);
\item phase damping: diagonal Kraus operators, populations untouched, coherence $\times\sqrt{1-\lambda}$ (1).
\end{itemize}
Each set satisfies $\sum K^\dagger K = I$. Hence P$\to$4, Q$\to$2, R$\to$3, S$\to$1.

\textbf{\color{brand}Why the other options are incorrect.}
\begin{itemize}[leftmargin=2.4em, itemsep=1pt, topsep=2pt]
\item[\textbf{(A)}] Q$\to$1, S$\to$2 gives the phase flip the diagonal damping operators and vice versa; the channels are equivalent but the question asks for the stated Kraus forms.
\item[\textbf{(C)}] P$\to$2, Q$\to$4 swaps $X$ and $Z$: the bit flip applies $X$, the phase flip applies $Z$.
\item[\textbf{(D)}] R$\to$1, S$\to$3 swaps amplitude and phase damping; only amplitude damping has the off-diagonal jump $\ket0\!\bra1$.
\end{itemize}
\textbf{\color{brand}Key takeaway.} Bit flip $X$, phase flip $Z$, AD jump $\ket0\!\bra1$, PD diagonal.
\end{solbox}
\qsep

\needspace{13\baselineskip}
\begin{qitem}{36}
\qmeta{Matching lists}{Channels on the Bloch ball}
Match each channel (List-I) with its action on the Bloch vector $(r_x, r_y, r_z)$ (List-II).
\par\smallskip\noindent{\renewcommand{\arraystretch}{1.35}%
\begin{tabularx}{\linewidth}{@{}c>{\raggedright\arraybackslash}p{0.33\linewidth}@{\hspace{1.2em}}cL@{}}
\toprule
\multicolumn{2}{@{}l}{\textbf{List-I}} & \multicolumn{2}{l@{}}{\textbf{List-II}} \\ \midrule
\textbf{(P)} & depolarizing (CUDA-Q convention, $p$) & \textbf{(1)} & $r_x, r_y$ scaled, $r_z$ unchanged \\
\textbf{(Q)} & amplitude damping ($\gamma$) & \textbf{(2)} & all three components multiplied by the same factor $1 - \frac{4p}{3}$ \\
\textbf{(R)} & phase damping ($\lambda$) & \textbf{(3)} & $r_x$ unchanged, $r_y$ and $r_z$ multiplied by $1 - 2p$ \\
\textbf{(S)} & bit flip ($p$) & \textbf{(4)} & $r_x, r_y$ multiplied by $\sqrt{1-\gamma}$, $r_z \mapsto \gamma + (1-\gamma)r_z$ \\
\bottomrule
\end{tabularx}}
\par\smallskip\noindent\begin{tabularx}{\linewidth}{@{}*{2}{L}@{}}
\optlabel{A}\,P$\to$2, Q$\to$1, R$\to$4, S$\to$3 & \optlabel{B}\,P$\to$2, Q$\to$4, R$\to$3, S$\to$1 \\[3pt]
\optlabel{C}\,P$\to$3, Q$\to$4, R$\to$1, S$\to$2 & \correctlabel{D}\,P$\to$2, Q$\to$4, R$\to$1, S$\to$3
\end{tabularx}
\end{qitem}
\begin{solbox}{D}
\textbf{\color{brand}Correct answer.} \textcolor{okgreen}{\textbf{(D)}} P$\to$2, Q$\to$4, R$\to$1, S$\to$3

\textbf{\color{brand}Concepts required.} Geometric action of channels on the Bloch ball.

\textbf{\color{brand}Relevant formulae.} Depolarizing $\vec r \mapsto (1 - \frac{4p}{3})\vec r$; AD $r_z \mapsto \gamma + (1-\gamma)r_z$.

\textbf{\color{brand}Step-by-step explanation.} \begin{itemize}\itemsep0pt
\item depolarizing: an isotropic shrink of the whole ball (2);
\item amplitude damping: shrinks and pushes the ball towards the north pole $\ket0$ --- the only non-unital one here (4);
\item phase damping: shrinks towards the $z$-axis (1);
\item bit flip: shrinks towards the $x$-axis (3).
\end{itemize}
Hence P$\to$2, Q$\to$4, R$\to$1, S$\to$3.

\textbf{\color{brand}Why the other options are incorrect.}
\begin{itemize}[leftmargin=2.4em, itemsep=1pt, topsep=2pt]
\item[\textbf{(A)}] Q$\to$1, R$\to$4 swaps AD and PD: only AD shifts the centre.
\item[\textbf{(B)}] R$\to$3, S$\to$1: PD keeps $r_z$, bit flip keeps $r_x$.
\item[\textbf{(C)}] P$\to$3, S$\to$2 gives depolarizing an axis; it is isotropic.
\end{itemize}
\textbf{\color{brand}Key takeaway.} Only amplitude damping moves the centre of the ball.
\end{solbox}
\qsep

\needspace{13\baselineskip}
\begin{qitem}{37}
\qmeta{Assertion--Reason}{Non-unital amplitude damping}
\textbf{Assertion (A):} Amplitude damping is not unital: it does not map the maximally mixed state $I/2$ to itself.\\[1mm]
\textbf{Reason (R):} Amplitude damping multiplies the off-diagonal elements of $\rho$ by $\sqrt{1 - \gamma}$.
\begin{itemize}[leftmargin=2.6em, labelsep=0.4em, itemsep=2pt, topsep=3pt]
\item[\optlabel{A}] Both Assertion (A) and Reason (R) are true, and R is the correct explanation of A
\item[\correctlabel{B}] Both A and R are true, but R is \emph{not} the correct explanation of A
\item[\optlabel{C}] A is true, but R is false
\item[\optlabel{D}] A is false, but R is true
\end{itemize}
\end{qitem}
\begin{solbox}{B}
\textbf{\color{brand}Correct answer.} \textcolor{okgreen}{\textbf{(B)}} Both A and R are true, but R is \emph{not} the correct explanation of A

\textbf{\color{brand}Concepts required.} Unital vs non-unital channels.

\textbf{\color{brand}Relevant formulae.} Unital: $\mathcal E(I) = I$; AD: $\mathcal E(I/2) = \operatorname{diag}(\frac{1+\gamma}2, \frac{1-\gamma}2)$.

\textbf{\color{brand}Step-by-step explanation.} \textbf{A is true}: $I/2 \mapsto \operatorname{diag}(\frac{1+\gamma}{2}, \frac{1-\gamma}{2}) \neq I/2$ --- energy flows to $\ket0$.
\textbf{R is true} as well, but it does not explain A: phase damping shrinks the off-diagonals in the same way and \emph{is} unital. Non-unitality
comes from the population transfer $\rho_{00} \mapsto \rho_{00} + \gamma\rho_{11}$ (the jump operator $\ket0\!\bra1$). Answer: (B).

\textbf{\color{brand}Why the other options are incorrect.}
\begin{itemize}[leftmargin=2.4em, itemsep=1pt, topsep=2pt]
\item[\textbf{(A)}] needs R to explain A; phase damping shows that coherence decay does not imply non-unitality.
\item[\textbf{(C)}] needs R false; the $\sqrt{1-\gamma}$ factor is correct.
\item[\textbf{(D)}] needs A false; AD is the standard non-unital channel.
\end{itemize}
\textbf{\color{brand}Key takeaway.} Non-unital = drives towards a preferred state (here $\ket0$).
\end{solbox}
\qsep

\needspace{13\baselineskip}
\begin{qitem}{38}
\qmeta{Statement I / II}{Complete dephasing vs unitary $Z$}
\textbf{Statement I:} The phase-flip channel with $p = \frac12$ maps every single-qubit state to a state with zero off-diagonal elements.\\[1mm]
\textbf{Statement II:} The phase-flip channel with $p = 1$ destroys all quantum information.
\begin{itemize}[leftmargin=2.6em, labelsep=0.4em, itemsep=2pt, topsep=3pt]
\item[\optlabel{A}] Both Statement I and Statement II are true
\item[\optlabel{B}] Both Statement I and Statement II are false
\item[\correctlabel{C}] Statement I is true, but Statement II is false
\item[\optlabel{D}] Statement I is false, but Statement II is true
\end{itemize}
\end{qitem}
\begin{solbox}{C}
\textbf{\color{brand}Correct answer.} \textcolor{okgreen}{\textbf{(C)}} Statement I is true, but Statement II is false

\textbf{\color{brand}Concepts required.} Unitary vs random errors; complete dephasing.

\textbf{\color{brand}Relevant formulae.} Coherence $\times(1 - 2p)$; $p = 1$: $\rho \mapsto Z\rho Z$.

\textbf{\color{brand}Step-by-step explanation.} \textbf{I is true}: the coherence is multiplied by $1 - 2p = 0$ (complete dephasing). \textbf{II is false}: with $p = 1$ the channel is
$\rho \mapsto Z\rho Z$, a \emph{unitary} that can be undone by applying $Z$ again. Information is lost only when the error is random ($0 < p < 1$),
most for $p = \frac12$. Answer: (C).

\textbf{\color{brand}Why the other options are incorrect.}
\begin{itemize}[leftmargin=2.4em, itemsep=1pt, topsep=2pt]
\item[\textbf{(A)}] requires II to be true, but $\rho \mapsto Z\rho Z$ is reversible by applying $Z$ again.
\item[\textbf{(B)}] requires I to be false, but at $p = \frac12$ the coherence factor $1 - 2p$ is exactly 0.
\item[\textbf{(D)}] asserts I false and II true --- exactly the reverse of the correct judgement.
\end{itemize}
\textbf{\color{brand}Key takeaway.} Information loss needs randomness, maximal at $p = \frac12$.
\end{solbox}
\qsep

\needspace{20\baselineskip}
\begin{passagebox}{Paragraph for Questions 39 and 40}
A superconducting qubit has $T_1 = 80\,\mu$s and $T_2 = 60\,\mu$s. It is prepared in $\ket+$ and then idles for $t = 20\,\mu$s.
Use the standard model: populations relax towards $\ket0$ as $\rho_{11}(t) = \rho_{11}(0)e^{-t/T_1}$, and the coherence decays as
$|\rho_{01}(t)| = |\rho_{01}(0)|e^{-t/T_2}$, with $\frac{1}{T_2} = \frac{1}{2T_1} + \frac{1}{T_\varphi}$ ($T_\varphi$ = pure-dephasing time).
\end{passagebox}
\needspace{13\baselineskip}
\begin{qitem}{39}
\qmeta{Paragraph based}{$T_1$/$T_2$ decay of $\ket+$}
After the idle period, $\rho_{11}$ and $|\rho_{01}|$ are approximately
\par\smallskip\noindent\begin{tabularx}{\linewidth}{@{}*{4}{L}@{}}
\correctlabel{A}\,$0.389$ and $0.358$ & \optlabel{B}\,$0.500$ and $0.358$ & \optlabel{C}\,$0.389$ and $0.316$ & \optlabel{D}\,$0.389$ and $0.389$
\end{tabularx}
\end{qitem}
\begin{solbox}{A}
\textbf{\color{brand}Correct answer.} \textcolor{okgreen}{\textbf{(A)}} $0.389$ and $0.358$

\textbf{\color{brand}Concepts required.} Applying $T_1$ and $T_2$ decay to a state.

\textbf{\color{brand}Relevant formulae.} $\rho_{11}(t) = \rho_{11}(0)e^{-t/T_1}$, $|\rho_{01}(t)| = |\rho_{01}(0)|e^{-t/T_2}$.

\textbf{\color{brand}Step-by-step explanation.} $\rho_{11}(0) = |\rho_{01}(0)| = \frac12$ for $\ket+$. $\rho_{11} = \frac12e^{-20/80} = \frac12e^{-0.25} = 0.389$.
$|\rho_{01}| = \frac12e^{-20/60} = \frac12e^{-1/3} = 0.358$. The $T_2$ decay already \emph{includes} the relaxation contribution; multiplying by
$e^{-t/2T_1}$ again would double-count it ($0.316$).

\textbf{\color{brand}Why the other options are incorrect.}
\begin{itemize}[leftmargin=2.4em, itemsep=1pt, topsep=2pt]
\item[\textbf{(B)}] $\rho_{11} = 0.5$ assumes no relaxation; with $T_1 = 80\,\mu$s, $e^{-20/80} = 0.78$ of the excited population survives.
\item[\textbf{(C)}] $0.316$ multiplies the $T_2$ decay by an extra $e^{-t/2T_1}$; $T_2$ already contains the relaxation contribution.
\item[\textbf{(D)}] $0.389$ for the coherence uses $e^{-t/T_1}$; coherence decays with $T_2$, not $T_1$.
\end{itemize}
\textbf{\color{brand}Key takeaway.} $T_1$ for populations, $T_2$ for coherences.
\end{solbox}
\qsep

\needspace{13\baselineskip}
\begin{qitem}{40}
\qmeta{Paragraph based}{Fidelity and $T_\varphi$}
The fidelity $\braket{+|\rho|+}$ after the idle period, and the pure-dephasing time $T_\varphi$, are
\par\smallskip\noindent\begin{tabularx}{\linewidth}{@{}*{2}{L}@{}}
\optlabel{A}\,$0.889$ and $96\,\mu$s & \optlabel{B}\,$0.716$ and $96\,\mu$s \\[3pt]
\correctlabel{C}\,$0.858$ and $96\,\mu$s & \optlabel{D}\,$0.858$ and $43.6\,\mu$s
\end{tabularx}
\end{qitem}
\begin{solbox}{C}
\textbf{\color{brand}Correct answer.} \textcolor{okgreen}{\textbf{(C)}} $0.858$ and $96\,\mu$s

\textbf{\color{brand}Concepts required.} Fidelity with $\ket+$; extracting $T_\varphi$.

\textbf{\color{brand}Relevant formulae.} $F = \frac12 + \operatorname{Re}\rho_{01}$; $\frac1{T_\varphi} = \frac1{T_2} - \frac1{2T_1}$.

\textbf{\color{brand}Step-by-step explanation.} $\braket{+|\rho|+} = \frac12(\rho_{00} + \rho_{11}) + \operatorname{Re}\rho_{01} = \frac12 + 0.358 = 0.858$ (populations sum to 1 whatever their split).
$\frac{1}{T_\varphi} = \frac{1}{T_2} - \frac{1}{2T_1} = \frac{1}{60} - \frac{1}{160} = 0.01667 - 0.00625 = 0.01042\ \mu\text{s}^{-1}$, so $T_\varphi = 96\,\mu$s.
Adding instead of subtracting gives $43.6\,\mu$s.

\textbf{\color{brand}Why the other options are incorrect.}
\begin{itemize}[leftmargin=2.4em, itemsep=1pt, topsep=2pt]
\item[\textbf{(A)}] $0.889 = \frac12 + \rho_{11}$ uses a population instead of the coherence $\operatorname{Re}\rho_{01}$.
\item[\textbf{(B)}] $0.716 = 2|\rho_{01}|$ forgets the constant $\frac12$ from $\frac12(\rho_{00} + \rho_{11})$ and halving the coherence term.
\item[\textbf{(D)}] $43.6\,\mu$s adds the rates $\frac1{T_2} + \frac{1}{2T_1}$ instead of subtracting them.
\end{itemize}
\textbf{\color{brand}Key takeaway.} Fidelity with $\ket+$ depends only on the coherence.
\end{solbox}
\qsep

\newpage

\begin{sectionbanner}SECTION C \enspace$\cdot$\enspace Noise Models and Simulation \hfill Q.41--60\end{sectionbanner}

\needspace{13\baselineskip}
\begin{qitem}{41}
\qmeta{Single-correct MCQ}{CUDA-Q depolarizing convention}
On the target \texttt{density-matrix-cpu} a kernel prepares $\ket0$ and applies \texttt{rx(t, q[0])}, with \texttt{t = 0.0} passed as a kernel
argument. The noise model attaches \texttt{cudaq.DepolarizationChannel(0.3)} to every \texttt{rx}. \texttt{cudaq.observe} of $Z$ returns
\par\smallskip\noindent\begin{tabularx}{\linewidth}{@{}*{4}{L}@{}}
\optlabel{A}\,$0.7$ & \correctlabel{B}\,$0.6$ & \optlabel{C}\,$0.4$ & \optlabel{D}\,$1.0$
\end{tabularx}
\end{qitem}
\begin{solbox}{B}
\textbf{\color{brand}Correct answer.} \textcolor{okgreen}{\textbf{(B)}} $0.6$

\textbf{\color{brand}Concepts required.} CUDA-Q's depolarizing convention; channels act after the gate.

\textbf{\color{brand}Relevant formulae.} $\vec r \mapsto (1 - \frac{4p}{3})\vec r$.

\textbf{\color{brand}Step-by-step explanation.} \texttt{rx(0)} is the identity, so the state before the channel is $\ket0$ ($r_z = 1$). CUDA-Q's depolarizing channel has Kraus operators
$\sqrt{1-p}\,I$, $\sqrt{p/3}\,X$, $\sqrt{p/3}\,Y$, $\sqrt{p/3}\,Z$; it multiplies the Bloch vector by $1 - \frac{4p}{3} = 1 - 0.4 = 0.6$.
So $\langle Z\rangle = 0.6$ (measured with CUDA-Q 0.16 for the lecture). $\rho_{11} = \frac{2p}{3} = 0.2$.

\textbf{\color{brand}Why the other options are incorrect.}
\begin{itemize}[leftmargin=2.4em, itemsep=1pt, topsep=2pt]
\item[\textbf{(A)}] $0.7 = 1 - p$ is the weight of the identity Kraus operator, not $\langle Z\rangle$; the $X$ and $Y$ branches flip $Z$.
\item[\textbf{(C)}] $0.4 = 1 - 2p$ is the factor of a bit-flip (or phase-flip) channel with probability $p$, not of the depolarizing channel.
\item[\textbf{(D)}] $1.0$ would mean the noise never acted --- which is what happens when the identity gate is removed by the compiler (Question Q.42).
\end{itemize}
\textbf{\color{brand}Key takeaway.} $\langle Z\rangle = 1 - \frac{4p}{3}$ on $\ket0$.
\end{solbox}
\qsep

\needspace{13\baselineskip}
\begin{qitem}{42}
\qmeta{Single-correct MCQ}{Constant identity gates and noise}
The same experiment is repeated, but the kernel now contains \texttt{rx(0.0, q[0])} with the angle written as a \emph{literal constant}. With CUDA-Q 0.16
the result is $\langle Z\rangle = 1.0$. The best explanation is
\begin{itemize}[leftmargin=2.6em, labelsep=0.4em, itemsep=2pt, topsep=3pt]
\item[\correctlabel{A}] the compiler removes the constant identity rotation, so there is no \texttt{rx} gate left to carry the noise channel
\item[\optlabel{B}] noise channels are applied only before a gate, never after it
\item[\optlabel{C}] \texttt{density-matrix-cpu} ignores noise models
\item[\optlabel{D}] the depolarizing channel acts only on non-zero rotation angles by definition
\end{itemize}
\end{qitem}
\begin{solbox}{A}
\textbf{\color{brand}Correct answer.} \textcolor{okgreen}{\textbf{(A)}} the compiler removes the constant identity rotation, so there is no \texttt{rx} gate left to carry the noise channel

\textbf{\color{brand}Concepts required.} Noise channels attach to gates of the compiled circuit.

\textbf{\color{brand}Relevant formulae.} Channel fires $\Leftrightarrow$ the gate survives compilation.

\textbf{\color{brand}Step-by-step explanation.} Noise models attach channels to \emph{gate occurrences} in the compiled circuit. A rotation by a constant $0$ is an identity and is optimised away
during compilation, so the channel never fires; with a run-time argument the gate survives (Question Q.41). The channel itself does not depend on the
angle, the density-matrix target does apply noise, and channels are applied \emph{after} the gate. The lecture's experiments therefore pass idle angles
as arguments.

\textbf{\color{brand}Why the other options are incorrect.}
\begin{itemize}[leftmargin=2.4em, itemsep=1pt, topsep=2pt]
\item[\textbf{(B)}] CUDA-Q applies the channel after each matching gate.
\item[\textbf{(C)}] The same target applies noise when the angle is an argument.
\item[\textbf{(D)}] The channel's Kraus operators do not depend on the angle.
\end{itemize}
\textbf{\color{brand}Key takeaway.} Pass identity angles as arguments, or the idle gate (and its noise) disappears.
\end{solbox}
\qsep

\needspace{13\baselineskip}
\begin{qitem}{43}
\qmeta{Single-correct MCQ}{CPTP check}
Which set of operators defines a valid (completely positive, trace-preserving) quantum channel?
\begin{itemize}[leftmargin=2.6em, labelsep=0.4em, itemsep=2pt, topsep=3pt]
\item[\optlabel{A}] $\Big\{\begin{psmallmatrix}1&0\\0&0.6\end{psmallmatrix},\ \begin{psmallmatrix}0&0.6\\0&0\end{psmallmatrix}\Big\}$
\item[\optlabel{B}] $\{I,\ \sqrt{0.1}\,X\}$
\item[\optlabel{C}] $\{\sqrt{0.5}\,I,\ \sqrt{0.5}\,X,\ \sqrt{0.5}\,Z\}$
\item[\correctlabel{D}] $\{\sqrt{0.7}\,I,\ \sqrt{0.3}\,Y\}$
\end{itemize}
\end{qitem}
\begin{solbox}{D}
\textbf{\color{brand}Correct answer.} \textcolor{okgreen}{\textbf{(D)}} $\{\sqrt{0.7}\,I,\ \sqrt{0.3}\,Y\}$

\textbf{\color{brand}Concepts required.} CPTP condition.

\textbf{\color{brand}Relevant formulae.} $\sum_kK_k^\dagger K_k = I$.

\textbf{\color{brand}Step-by-step explanation.} Any map $\rho \mapsto \sum K\rho K^\dagger$ is completely positive; it is trace preserving iff $\sum K^\dagger K = I$.
$0.7I + 0.3Y^2 = I$ \checkmark. $I + 0.1X^2 = 1.1I$; $\operatorname{diag}(1, 0.36) + \operatorname{diag}(0, 0.36) = \operatorname{diag}(1, 0.72)$; $0.5 + 0.5 + 0.5 = 1.5$:
the other three create or destroy probability. The first is the bit-phase-flip channel with $p = 0.3$.

\textbf{\color{brand}Why the other options are incorrect.}
\begin{itemize}[leftmargin=2.4em, itemsep=1pt, topsep=2pt]
\item[\textbf{(A)}] Here $\sum K^\dagger K = \operatorname{diag}(1, 0.36 + 0.36) = \operatorname{diag}(1, 0.72)$: probability is lost from $\ket1$ (it would need $0.8$ in the jump operator).
\item[\textbf{(B)}] $\{I, \sqrt{0.1}X\}$: $\sum K^\dagger K = I + 0.1I = 1.1I$, so the trace grows by 10\,\%.
\item[\textbf{(C)}] $0.5 + 0.5 + 0.5 = 1.5$: $\sum K^\dagger K = 1.5I$ creates probability.
\end{itemize}
\textbf{\color{brand}Key takeaway.} Check $\sum K^\dagger K = I$.
\end{solbox}
\qsep

\needspace{13\baselineskip}
\begin{qitem}{44}
\qmeta{Single-correct MCQ}{Applying the bit-flip channel}
The bit-flip channel with $p = 0.2$ acts on $\rho = \begin{pmatrix}0.9 & 0.3\\ 0.3 & 0.1\end{pmatrix}$. The output is
\par\smallskip\noindent\begin{tabularx}{\linewidth}{@{}*{4}{L}@{}}
\optlabel{A}\,$\begin{pmatrix}0.74 & 0.18\\ 0.18 & 0.26\end{pmatrix}$ & \correctlabel{B}\,$\begin{pmatrix}0.74 & 0.30\\ 0.30 & 0.26\end{pmatrix}$ & \optlabel{C}\,$\begin{pmatrix}0.82 & 0.30\\ 0.30 & 0.18\end{pmatrix}$ & \optlabel{D}\,$\begin{pmatrix}0.90 & 0.18\\ 0.18 & 0.10\end{pmatrix}$
\end{tabularx}
\end{qitem}
\begin{solbox}{B}
\textbf{\color{brand}Correct answer.} \textcolor{okgreen}{\textbf{(B)}} $\begin{pmatrix}0.74 & 0.30\\ 0.30 & 0.26\end{pmatrix}$

\textbf{\color{brand}Concepts required.} Applying a Kraus channel to a matrix.

\textbf{\color{brand}Relevant formulae.} $\mathcal E(\rho) = (1-p)\rho + pX\rho X$; $X\rho X$ swaps rows and columns.

\textbf{\color{brand}Step-by-step explanation.} $X\rho X$ swaps both rows and columns: $\begin{psmallmatrix}0.1&0.3\\0.3&0.9\end{psmallmatrix}$. Then
$0.8\rho + 0.2X\rho X = \begin{psmallmatrix}0.72 + 0.02 & 0.24 + 0.06\\ 0.24 + 0.06 & 0.08 + 0.18\end{psmallmatrix} = \begin{psmallmatrix}0.74 & 0.30\\ 0.30 & 0.26\end{psmallmatrix}$.
The real off-diagonal ($r_x = 0.6$) is untouched; $r_z$ shrinks from $0.8$ to $0.8(1 - 2p) = 0.48$.

\textbf{\color{brand}Why the other options are incorrect.}
\begin{itemize}[leftmargin=2.4em, itemsep=1pt, topsep=2pt]
\item[\textbf{(A)}] Multiplying the off-diagonal by $0.6$ is what a phase flip (or $r_x$ shrink) would do; a bit flip keeps $r_x$ and the real off-diagonal.
\item[\textbf{(C)}] This corresponds to $p = 0.1$: $0.9\rho + 0.1X\rho X = \begin{psmallmatrix}0.82&0.3\\0.3&0.18\end{psmallmatrix}$.
\item[\textbf{(D)}] This is the phase-flip result: populations unchanged and coherence $\times(1 - 2p)$.
\end{itemize}
\textbf{\color{brand}Key takeaway.} Bit flips mix the populations, keep $r_x$.
\end{solbox}
\qsep

\needspace{13\baselineskip}
\begin{qitem}{45}
\qmeta{Single-correct MCQ}{Composing depolarizing channels}
Two depolarizing channels (CUDA-Q convention) with $p_1 = 0.03$ and $p_2 = 0.06$ act one after the other. They equal a single depolarizing channel with $p$ equal to
\par\smallskip\noindent\begin{tabularx}{\linewidth}{@{}*{4}{L}@{}}
\optlabel{A}\,$0.1168$ & \optlabel{B}\,$0.0882$ & \correctlabel{C}\,$0.0876$ & \optlabel{D}\,$0.0900$
\end{tabularx}
\end{qitem}
\begin{solbox}{C}
\textbf{\color{brand}Correct answer.} \textcolor{okgreen}{\textbf{(C)}} $0.0876$

\textbf{\color{brand}Concepts required.} Composition of depolarizing channels.

\textbf{\color{brand}Relevant formulae.} Shrink factors multiply: $(1 - \frac{4p_1}{3})(1 - \frac{4p_2}{3})$.

\textbf{\color{brand}Step-by-step explanation.} Each channel multiplies the Bloch vector by $1 - \frac{4p}{3}$: $0.96$ and $0.92$. Composed: $0.96\times0.92 = 0.8832 = 1 - \frac{4p}{3}$, so
$p = \frac34(1 - 0.8832) = 0.0876$. Probabilities do not simply add ($0.09$); composing ``no-error'' probabilities, $1 - 0.97\cdot0.94 = 0.0882$, ignores that two
Pauli errors can cancel.

\textbf{\color{brand}Why the other options are incorrect.}
\begin{itemize}[leftmargin=2.4em, itemsep=1pt, topsep=2pt]
\item[\textbf{(A)}] $0.1168 = 1 - 0.8832$ is the reduction of the Bloch length, i.e.\ $\frac{4p}{3}$, not $p$ itself.
\item[\textbf{(B)}] $0.0882 = 1 - 0.97\times0.94$ composes the \emph{no-error} probabilities and ignores that, e.g., $X$ followed by $X$ gives no error.
\item[\textbf{(D)}] $0.09 = p_1 + p_2$: error probabilities do not simply add, because two errors can combine into the identity or another Pauli.
\end{itemize}
\textbf{\color{brand}Key takeaway.} Compose channels by multiplying their Bloch shrink factors.
\end{solbox}
\qsep

\needspace{13\baselineskip}
\begin{qitem}{46}
\qmeta{Single-correct MCQ}{Notebook: depolarizing $(1-p)^G$}
In \texttt{Noise\_Simulation.ipynb} the 15-qubit GHZ circuit has $G$ noisy gate locations (1 Hadamard and 14 CNOTs, each followed by a depolarizing channel).
The notebook's no-error probability $(1 - p)^G$ at $p = 0.02$ is
\par\smallskip\noindent\begin{tabularx}{\linewidth}{@{}*{4}{L}@{}}
\optlabel{A}\,$0.7000$ & \optlabel{B}\,$0.9800$ & \correctlabel{C}\,$0.7386$ & \optlabel{D}\,$0.5455$
\end{tabularx}
\end{qitem}
\begin{solbox}{C}
\textbf{\color{brand}Correct answer.} \textcolor{okgreen}{\textbf{(C)}} $0.7386$

\textbf{\color{brand}Concepts required.} Error-free probability for the notebook's depolarizing GHZ.

\textbf{\color{brand}Relevant formulae.} $(1-p)^{g_1}(1-p)^{g_2}$ with $g_1 = 1$, $g_2 = 14$.

\textbf{\color{brand}Step-by-step explanation.} For depolarizing noise each gate location (one- or two-qubit) fails with probability $p$, so $G = 1 + 14 = 15$ and
$(0.98)^{15} = e^{15\ln0.98} = e^{-0.3030} = 0.7386$ (printed in the notebook). $1 - 15p = 0.70$ is the linear approximation; $0.98^{30} = 0.5455$
would count two failure chances per CNOT, which is the rule the notebook uses for bit-flip and phase-flip noise, not for depolarizing noise.

\textbf{\color{brand}Why the other options are incorrect.}
\begin{itemize}[leftmargin=2.4em, itemsep=1pt, topsep=2pt]
\item[\textbf{(A)}] $0.70 = 1 - Gp$ is the first-order approximation $(1-p)^G \approx 1 - Gp$, inaccurate at $Gp = 0.3$.
\item[\textbf{(B)}] $0.98$ counts a single gate location instead of $G = 15$.
\item[\textbf{(D)}] $0.5455 = 0.98^{30}$ gives each CNOT two failure chances; for depolarizing noise the notebook uses one joint channel per CNOT.
\end{itemize}
\textbf{\color{brand}Key takeaway.} Depolarizing: one failure chance per gate location.
\end{solbox}
\qsep

\needspace{13\baselineskip}
\begin{qitem}{47}
\qmeta{Single-correct MCQ}{Notebook: bit-flip $(1-p)^{29}$}
For \texttt{NOISE\_MODEL = "bit\_flip"}, the notebook's CNOT channel is \texttt{Pauli2} with probabilities $p(1-p)$ for $IX$ and $XI$ and $p^2$ for $XX$. The
no-error probability of the 15-qubit GHZ circuit at $p = 0.01$ is
\par\smallskip\noindent\begin{tabularx}{\linewidth}{@{}*{4}{L}@{}}
\correctlabel{A}\,$0.7472$ & \optlabel{B}\,$0.8601$ & \optlabel{C}\,$0.7397$ & \optlabel{D}\,$0.9801$
\end{tabularx}
\end{qitem}
\begin{solbox}{A}
\textbf{\color{brand}Correct answer.} \textcolor{okgreen}{\textbf{(A)}} $0.7472$

\textbf{\color{brand}Concepts required.} Two-qubit bit-flip channel as independent flips.

\textbf{\color{brand}Relevant formulae.} $p(1-p) + p(1-p) + p^2 = 1 - (1-p)^2$.

\textbf{\color{brand}Step-by-step explanation.} The CNOT channel equals independent bit flips on both qubits: no error with probability $(1-p)^2$. With one Hadamard and 14 CNOTs:
$(1-p)\cdot(1-p)^{2\cdot14} = (1-p)^{29} = 0.99^{29} = e^{-0.2915} = 0.7472$. This is the notebook's \texttt{no\_error\_probability} for non-depolarizing
kinds ($1 - (1-p)^2$ per two-qubit gate). $0.8601$ is the depolarizing value.

\textbf{\color{brand}Why the other options are incorrect.}
\begin{itemize}[leftmargin=2.4em, itemsep=1pt, topsep=2pt]
\item[\textbf{(B)}] $0.8601 = 0.99^{15}$ is the depolarizing result (one failure chance per CNOT).
\item[\textbf{(C)}] $0.99^{30}$ counts the Hadamard as a two-qubit location; it has one failure chance, giving $1 + 2\times14 = 29$.
\item[\textbf{(D)}] $0.9801 = 0.99^2$ counts only one CNOT (two independent flips) and ignores the other 13 CNOTs and the Hadamard.
\end{itemize}
\textbf{\color{brand}Key takeaway.} Independent flips on both CNOT qubits: $(1-p)^{29}$ in total.
\end{solbox}
\qsep

\needspace{13\baselineskip}
\begin{qitem}{48}
\qmeta{Single-correct MCQ}{Density-matrix memory}
The density-matrix method stores $4^N$ complex numbers of 16 bytes. The memory for $N = 13$ qubits is about
\par\smallskip\noindent\begin{tabularx}{\linewidth}{@{}*{4}{L}@{}}
\correctlabel{A}\,$1.07$ GB & \optlabel{B}\,$131$ kB & \optlabel{C}\,$268$ MB & \optlabel{D}\,$4.29$ GB
\end{tabularx}
\end{qitem}
\begin{solbox}{A}
\textbf{\color{brand}Correct answer.} \textcolor{okgreen}{\textbf{(A)}} $1.07$ GB

\textbf{\color{brand}Concepts required.} Density-matrix memory.

\textbf{\color{brand}Relevant formulae.} $16\cdot4^N$ bytes.

\textbf{\color{brand}Step-by-step explanation.} $16\times4^{13} = 16\times67\,108\,864 = 1.074\times10^9$ bytes $\approx 1.07$ GB. $131$ kB is a state vector ($16\cdot2^{13}$); $268$ MB is $N = 12$ and
$4.29$ GB is $N = 14$ (values quoted in the notebook). Every extra qubit multiplies the memory by 4.

\textbf{\color{brand}Why the other options are incorrect.}
\begin{itemize}[leftmargin=2.4em, itemsep=1pt, topsep=2pt]
\item[\textbf{(B)}] $131$ kB $= 16\cdot2^{13}$ bytes is the memory of a 13-qubit \emph{state vector}, not a density matrix.
\item[\textbf{(C)}] $268$ MB $= 16\cdot4^{12}$ is the 12-qubit value.
\item[\textbf{(D)}] $4.29$ GB $= 16\cdot4^{14}$ is the 14-qubit value.
\end{itemize}
\textbf{\color{brand}Key takeaway.} Each qubit multiplies density-matrix memory by 4.
\end{solbox}
\qsep

\needspace{13\baselineskip}
\begin{qitem}{49}
\qmeta{Single-correct MCQ}{What \texttt{stim} can simulate}
Which simulation can the stabilizer target \texttt{stim} perform \emph{correctly}?
\begin{itemize}[leftmargin=2.6em, labelsep=0.4em, itemsep=2pt, topsep=3pt]
\item[\optlabel{A}] a GHZ circuit with amplitude-damping noise
\item[\optlabel{B}] a GHZ circuit followed by a $T$ gate, with bit-flip noise
\item[\correctlabel{C}] a GHZ circuit ($H$, CNOT) with depolarizing noise after every gate
\item[\optlabel{D}] the notebook's mirror circuit with random $R_Y$, $R_Z$ rotations and depolarizing noise
\end{itemize}
\end{qitem}
\begin{solbox}{C}
\textbf{\color{brand}Correct answer.} \textcolor{okgreen}{\textbf{(C)}} a GHZ circuit ($H$, CNOT) with depolarizing noise after every gate

\textbf{\color{brand}Concepts required.} Gottesman--Knill theorem.

\textbf{\color{brand}Relevant formulae.} Clifford gates + Pauli channels + Pauli measurements: polynomial simulation.

\textbf{\color{brand}Step-by-step explanation.} The stabilizer formalism (Gottesman--Knill) efficiently simulates Clifford circuits ($H$, $S$, CNOT, Paulis, measurement) with \emph{Pauli} noise, which
maps stabilizer states to probabilistic mixtures of stabilizer states. Arbitrary rotations and $T$ are non-Clifford; amplitude damping is not a Pauli channel.
The notebook raises an error in those cases instead of silently ignoring the noise.

\textbf{\color{brand}Why the other options are incorrect.}
\begin{itemize}[leftmargin=2.4em, itemsep=1pt, topsep=2pt]
\item[\textbf{(A)}] Amplitude damping is not a mixture of Paulis; \texttt{stim} cannot represent it (the notebook refuses this combination).
\item[\textbf{(B)}] The $T$ gate is not a Clifford gate; it maps stabilizer states outside the stabilizer set.
\item[\textbf{(D)}] The mirror circuit uses $R_Y$, $R_Z$ with random angles: non-Clifford gates that the stabilizer formalism cannot represent.
\end{itemize}
\textbf{\color{brand}Key takeaway.} \texttt{stim}: Clifford + Pauli noise only.
\end{solbox}
\qsep

\needspace{13\baselineskip}
\begin{qitem}{50}
\qmeta{Single-correct MCQ}{\texttt{as\_counts\_from\_rho}}
On the density-matrix path the notebook calls \texttt{as\_counts\_from\_rho(rho, n, shots)}. What does this function do?
\begin{itemize}[leftmargin=2.6em, labelsep=0.4em, itemsep=2pt, topsep=3pt]
\item[\optlabel{A}] it takes the diagonal of $\rho$ and returns it unchanged as counts, without sampling
\item[\correctlabel{B}] it takes the diagonal of the exact noisy $\rho$ as outcome probabilities, reorders the indices so that qubit 0 is written first, and draws \texttt{shots} counts from a multinomial distribution
\item[\optlabel{C}] it measures the density-matrix simulator \texttt{shots} times with \texttt{cudaq.sample}
\item[\optlabel{D}] it rounds every entry of $\rho$ to the nearest multiple of $1/\texttt{shots}$
\end{itemize}
\end{qitem}
\begin{solbox}{B}
\textbf{\color{brand}Correct answer.} \textcolor{okgreen}{\textbf{(B)}} it takes the diagonal of the exact noisy $\rho$ as outcome probabilities, reorders the indices so that qubit 0 is written first, and draws \texttt{shots} counts from a multinomial distribution

\textbf{\color{brand}Concepts required.} Sampling from the exact diagonal of $\rho$.

\textbf{\color{brand}Relevant formulae.} Counts $\sim\mathrm{Multinomial}(S, \operatorname{diag}\rho)$.

\textbf{\color{brand}Step-by-step explanation.} The function clips and normalises \texttt{np.real(np.diag(rho))}, applies \texttt{bit\_reverse\_order} (CUDA-Q stores qubit 0 as the least significant
bit), draws \texttt{counts\_rng.multinomial(shots, probs)} with a seeded NumPy generator, and also returns the exact distribution. Because measurement
probabilities are exactly the diagonal of $\rho$, this is statistically identical to noisy sampling --- and much faster than noisy \texttt{cudaq.sample} on the
density-matrix target in CUDA-Q 0.16 (about 115 s for 10 qubits, according to the README).

\textbf{\color{brand}Why the other options are incorrect.}
\begin{itemize}[leftmargin=2.4em, itemsep=1pt, topsep=2pt]
\item[\textbf{(A)}] It does sample (\texttt{multinomial}); the exact distribution is returned separately as a second value.
\item[\textbf{(C)}] The function never calls \texttt{cudaq.sample}; noisy sampling on the density-matrix target was too slow in CUDA-Q 0.16.
\item[\textbf{(D)}] There is no rounding: counts are random multinomial draws from the exact probabilities.
\end{itemize}
\textbf{\color{brand}Key takeaway.} Exact $\rho$ + multinomial draws = noisy shots, fast.
\end{solbox}
\qsep

\needspace{13\baselineskip}
\begin{qitem}{51}
\qmeta{Single-correct MCQ}{Hellinger fidelity and TVD}
The ideal distribution of a 2-qubit circuit is $P(00) = P(11) = 0.5$. The noisy distribution is $Q(00) = Q(11) = 0.45$, $Q(01) = Q(10) = 0.05$. The Hellinger
fidelity $F_H = \big(\sum_x\sqrt{P_xQ_x}\big)^2$ and the total variation distance $\frac12\sum_x|P_x - Q_x|$ are
\par\smallskip\noindent\begin{tabularx}{\linewidth}{@{}*{4}{L}@{}}
\optlabel{A}\,$0.900$ and $0.200$ & \optlabel{B}\,$0.949$ and $0.100$ & \optlabel{C}\,$0.810$ and $0.100$ & \correctlabel{D}\,$0.900$ and $0.100$
\end{tabularx}
\end{qitem}
\begin{solbox}{D}
\textbf{\color{brand}Correct answer.} \textcolor{okgreen}{\textbf{(D)}} $0.900$ and $0.100$

\textbf{\color{brand}Concepts required.} Distances between probability distributions.

\textbf{\color{brand}Relevant formulae.} $F_H = (\sum\sqrt{P_xQ_x})^2$, TVD $= \frac12\sum|P_x - Q_x|$.

\textbf{\color{brand}Step-by-step explanation.} $\sum_x\sqrt{P_xQ_x} = 2\sqrt{0.5\times0.45} = 2(0.4743) = 0.9487$; squared: $F_H = 0.900$. Only the two ideal outcomes contribute.
TVD $= \frac12(0.05 + 0.05 + 0.05 + 0.05) = 0.100$. For a balanced two-peak ideal distribution $F_H$ equals the population $Q(00) + Q(11)$ --- which is why the
notebook's Hellinger fidelity and GHZ population columns coincide.

\textbf{\color{brand}Why the other options are incorrect.}
\begin{itemize}[leftmargin=2.4em, itemsep=1pt, topsep=2pt]
\item[\textbf{(A)}] $0.200$ forgets the factor $\frac12$ in the total variation distance.
\item[\textbf{(B)}] $0.949$ is the Bhattacharyya coefficient $\sum\sqrt{P_xQ_x}$ before squaring.
\item[\textbf{(C)}] $0.810$ squares the result once more ($0.9^2$).
\end{itemize}
\textbf{\color{brand}Key takeaway.} For a two-peak ideal distribution, $F_H$ = population.
\end{solbox}
\qsep

\needspace{13\baselineskip}
\begin{qitem}{52}
\qmeta{Single-correct MCQ}{Parity from $X$-basis counts}
A 2-qubit GHZ state is measured after Hadamards on both qubits ($X$ basis): counts \texttt{\{00: 450, 11: 400, 01: 80, 10: 70\}}. The parity
$\langle X\otimes X\rangle$ is
\par\smallskip\noindent\begin{tabularx}{\linewidth}{@{}*{4}{L}@{}}
\optlabel{A}\,$0.40$ & \optlabel{B}\,$0.15$ & \correctlabel{C}\,$0.70$ & \optlabel{D}\,$0.85$
\end{tabularx}
\end{qitem}
\begin{solbox}{C}
\textbf{\color{brand}Correct answer.} \textcolor{okgreen}{\textbf{(C)}} $0.70$

\textbf{\color{brand}Concepts required.} Parity from counts.

\textbf{\color{brand}Relevant formulae.} $\langle X^{\otimes N}\rangle = P(\text{even}) - P(\text{odd})$ in the $X$ basis.

\textbf{\color{brand}Step-by-step explanation.} Even parity ($00$, $11$) counts $+1$, odd parity ($01$, $10$) counts $-1$: $\langle XX\rangle = \frac{(450 + 400) - (80 + 70)}{1000} = \frac{850 - 150}{1000} = 0.70$.
This is the notebook's \texttt{parity\_expectation}. The ideal GHZ value is 1; phase and bit errors both lower it.

\textbf{\color{brand}Why the other options are incorrect.}
\begin{itemize}[leftmargin=2.4em, itemsep=1pt, topsep=2pt]
\item[\textbf{(A)}] $0.40$ subtracts the counts of $11$ from those of $00$ and ignores the odd outcomes.
\item[\textbf{(B)}] $0.15 = 150/1000$ is $P(\text{odd parity})$, one of the two terms, not the difference $P(\text{even}) - P(\text{odd})$.
\item[\textbf{(D)}] $0.85$ is $P(\text{even parity})$, not $P(\text{even}) - P(\text{odd})$.
\end{itemize}
\textbf{\color{brand}Key takeaway.} Parity = even minus odd.
\end{solbox}
\qsep

\needspace{13\baselineskip}
\begin{qitem}{53}
\qmeta{Single-correct MCQ}{Trajectories vs density matrix}
For $N = 20$ qubits, by what factor does the memory of exact density-matrix simulation exceed that of one trajectory (state vector) in the same precision?
\par\smallskip\noindent\begin{tabularx}{\linewidth}{@{}*{4}{L}@{}}
\optlabel{A}\,$20$ & \optlabel{B}\,about $10^{12}$ & \optlabel{C}\,$2$ & \correctlabel{D}\,about $10^6$
\end{tabularx}
\end{qitem}
\begin{solbox}{D}
\textbf{\color{brand}Correct answer.} \textcolor{okgreen}{\textbf{(D)}} about $10^6$

\textbf{\color{brand}Concepts required.} Memory of density matrices vs trajectories.

\textbf{\color{brand}Relevant formulae.} $4^N/2^N = 2^N$.

\textbf{\color{brand}Step-by-step explanation.} $\frac{16\cdot4^{20}}{16\cdot2^{20}} = 2^{20} \approx 1.05\times10^6$. Trajectory (Monte-Carlo wave-function) simulation keeps one $2^N$ state and inserts random
Kraus jumps; averaging many trajectories reproduces $\rho$ statistically. It trades memory for repetitions --- CUDA-Q runs it on GPUs (\texttt{nvidia} target with a
noise model), which was not executed for these materials.

\textbf{\color{brand}Why the other options are incorrect.}
\begin{itemize}[leftmargin=2.4em, itemsep=1pt, topsep=2pt]
\item[\textbf{(A)}] $20 = N$ assumes the memory ratio grows linearly with the number of qubits.
\item[\textbf{(B)}] $10^{12} \approx 4^{20}$ is the number of density-matrix entries, not the ratio to the $2^{20}$ state vector.
\item[\textbf{(C)}] $2$ is the ratio of the bases ($4/2$), not of the sizes $4^N/2^N$.
\end{itemize}
\textbf{\color{brand}Key takeaway.} Trajectories trade $2^N$ memory for many repetitions.
\end{solbox}
\qsep

\needspace{13\baselineskip}
\begin{qitem}{54}
\qmeta{Single-correct MCQ}{\texttt{add\_all\_qubit\_channel}}
In \texttt{build\_noise\_model} the line \texttt{model.add\_all\_qubit\_channel("x", two, num\_controls=1)} means
\begin{itemize}[leftmargin=2.6em, labelsep=0.4em, itemsep=2pt, topsep=3pt]
\item[\optlabel{A}] the channel \texttt{two} is applied once, at the end of the circuit, to the first two qubits
\item[\optlabel{B}] the channel \texttt{two} replaces the CNOT gate
\item[\optlabel{C}] after every single-qubit \texttt{x} gate, the channel \texttt{two} acts on that qubit and its neighbour
\item[\correctlabel{D}] after every CNOT (an \texttt{x} gate with one control), the two-qubit channel \texttt{two} acts on its control and target qubits
\end{itemize}
\end{qitem}
\begin{solbox}{D}
\textbf{\color{brand}Correct answer.} \textcolor{okgreen}{\textbf{(D)}} after every CNOT (an \texttt{x} gate with one control), the two-qubit channel \texttt{two} acts on its control and target qubits

\textbf{\color{brand}Concepts required.} \texttt{NoiseModel.add\_all\_qubit\_channel} semantics.

\textbf{\color{brand}Relevant formulae.} (gate name, channel, number of controls) $\to$ channel after every matching gate on its qubits.

\textbf{\color{brand}Step-by-step explanation.} \texttt{add\_all\_qubit\_channel(name, channel, num\_controls)} attaches a channel to every occurrence of the named gate with that number of controls, on all qubits it
acts on. Plain \texttt{x} gates (no control) get the single-qubit channel \texttt{one} from the loop over \texttt{("h", "x", "ry", "rz")}. The channel is applied
\emph{in addition to}, and after, the ideal gate.

\textbf{\color{brand}Why the other options are incorrect.}
\begin{itemize}[leftmargin=2.4em, itemsep=1pt, topsep=2pt]
\item[\textbf{(A)}] The channel is attached to every gate occurrence of that type, not applied once at the end.
\item[\textbf{(B)}] The ideal CNOT is kept; the channel is applied after it.
\item[\textbf{(C)}] Plain \texttt{x} gates (no controls) get the single-qubit channel \texttt{one}; \texttt{num\_controls=1} restricts this line to CNOTs.
\end{itemize}
\textbf{\color{brand}Key takeaway.} \texttt{num\_controls=1} selects CNOTs.
\end{solbox}
\qsep

\needspace{13\baselineskip}
\begin{qitem}{55}
\qmeta{Matching lists}{Noise-capable CUDA-Q targets}
Match each CUDA-Q simulation target (List-I) with its property for noisy simulation (List-II).
\par\smallskip\noindent{\renewcommand{\arraystretch}{1.35}%
\begin{tabularx}{\linewidth}{@{}c>{\raggedright\arraybackslash}p{0.33\linewidth}@{\hspace{1.2em}}cL@{}}
\toprule
\multicolumn{2}{@{}l}{\textbf{List-I}} & \multicolumn{2}{l@{}}{\textbf{List-II}} \\ \midrule
\textbf{(P)} & \texttt{stim} & \textbf{(1)} & Monte-Carlo trajectories on a GPU, memory $\propto 2^N$, statistical \\
\textbf{(Q)} & \texttt{density-matrix-cpu} & \textbf{(2)} & noiseless state vector: not offered by the noise notebook (it ignores noise models) \\
\textbf{(R)} & \texttt{nvidia} with a noise model & \textbf{(3)} & Clifford circuits with Pauli noise, polynomial cost, hundreds of qubits \\
\textbf{(S)} & \texttt{qpp-cpu} & \textbf{(4)} & exact $\rho$, any Kraus channel, memory $16\cdot4^N$ bytes \\
\bottomrule
\end{tabularx}}
\par\smallskip\noindent\begin{tabularx}{\linewidth}{@{}*{2}{L}@{}}
\correctlabel{A}\,P$\to$3, Q$\to$4, R$\to$1, S$\to$2 & \optlabel{B}\,P$\to$3, Q$\to$4, R$\to$2, S$\to$1 \\[3pt]
\optlabel{C}\,P$\to$4, Q$\to$3, R$\to$1, S$\to$2 & \optlabel{D}\,P$\to$3, Q$\to$1, R$\to$4, S$\to$2
\end{tabularx}
\end{qitem}
\begin{solbox}{A}
\textbf{\color{brand}Correct answer.} \textcolor{okgreen}{\textbf{(A)}} P$\to$3, Q$\to$4, R$\to$1, S$\to$2

\textbf{\color{brand}Concepts required.} Noise-capable CUDA-Q targets.

\textbf{\color{brand}Relevant formulae.} stabilizer: poly; trajectory: $16\cdot2^N$; density matrix: $16\cdot4^N$.

\textbf{\color{brand}Step-by-step explanation.} \begin{itemize}\itemsep0pt
\item \texttt{stim}: the default on a CPU for GHZ with Pauli noise (recorded run) (3);
\item \texttt{density-matrix-cpu}: exact but $4^N$ (4);
\item \texttt{nvidia} + noise model: trajectory method, selected automatically when a GPU is visible (1);
\item \texttt{qpp-cpu}: a noiseless state-vector simulator; the notebook refuses it for noise (2).
\end{itemize}
Hence P$\to$3, Q$\to$4, R$\to$1, S$\to$2.

\textbf{\color{brand}Why the other options are incorrect.}
\begin{itemize}[leftmargin=2.4em, itemsep=1pt, topsep=2pt]
\item[\textbf{(B)}] R$\to$2, S$\to$1: \texttt{nvidia} with a noise model does simulate noise.
\item[\textbf{(C)}] P$\to$4, Q$\to$3 swaps stim and the density-matrix simulator.
\item[\textbf{(D)}] Q$\to$1, R$\to$4 swaps exact and trajectory simulation.
\end{itemize}
\textbf{\color{brand}Key takeaway.} Choose the method by circuit type, noise type and size.
\end{solbox}
\qsep

\needspace{13\baselineskip}
\begin{qitem}{56}
\qmeta{Matching lists}{Figures of merit}
Match each figure of merit of the notebook (List-I) with what it is sensitive to or requires (List-II).
\par\smallskip\noindent{\renewcommand{\arraystretch}{1.35}%
\begin{tabularx}{\linewidth}{@{}c>{\raggedright\arraybackslash}p{0.33\linewidth}@{\hspace{1.2em}}cL@{}}
\toprule
\multicolumn{2}{@{}l}{\textbf{List-I}} & \multicolumn{2}{l@{}}{\textbf{List-II}} \\ \midrule
\textbf{(P)} & $Z$-basis GHZ population & \textbf{(1)} & needs the full density matrix; also the off-diagonal (coherence) elements \\
\textbf{(Q)} & $X$-basis parity $\langle X^{\otimes N}\rangle$ & \textbf{(2)} & bit flips (errors that change computational-basis outcomes) \\
\textbf{(R)} & state fidelity $\braket{\psi|\rho|\psi}$ & \textbf{(3)} & phase errors, which are invisible in the $Z$ basis \\
\textbf{(S)} & purity $\operatorname{Tr}\rho^2$ & \textbf{(4)} & mixedness of the state, whatever the target \\
\bottomrule
\end{tabularx}}
\par\smallskip\noindent\begin{tabularx}{\linewidth}{@{}*{2}{L}@{}}
\optlabel{A}\,P$\to$2, Q$\to$3, R$\to$4, S$\to$1 & \optlabel{B}\,P$\to$2, Q$\to$1, R$\to$3, S$\to$4 \\[3pt]
\correctlabel{C}\,P$\to$2, Q$\to$3, R$\to$1, S$\to$4 & \optlabel{D}\,P$\to$3, Q$\to$2, R$\to$1, S$\to$4
\end{tabularx}
\end{qitem}
\begin{solbox}{C}
\textbf{\color{brand}Correct answer.} \textcolor{okgreen}{\textbf{(C)}} P$\to$2, Q$\to$3, R$\to$1, S$\to$4

\textbf{\color{brand}Concepts required.} What each figure of merit can detect.

\textbf{\color{brand}Relevant formulae.} $F_{\mathrm{GHZ}} = \frac12(\rho_{00} + \rho_{dd}) + \operatorname{Re}\rho_{0d}$.

\textbf{\color{brand}Step-by-step explanation.} \begin{itemize}\itemsep0pt
\item population: only diagonal elements; drops with bit flips (2);
\item $X$ parity: $Z$ errors flip its sign (3);
\item state fidelity: for GHZ $F = \frac12(\rho_{0,0} + \rho_{d,d} + 2\operatorname{Re}\rho_{0,d})$ --- needs the coherence $\rho_{0\cdots0,1\cdots1}$, available only on the density-matrix path (1);
\item purity: 1 for pure states, $2^{-N}$ for the maximally mixed state; independent of the target (4).
\end{itemize}
Hence P$\to$2, Q$\to$3, R$\to$1, S$\to$4.

\textbf{\color{brand}Why the other options are incorrect.}
\begin{itemize}[leftmargin=2.4em, itemsep=1pt, topsep=2pt]
\item[\textbf{(A)}] R$\to$4, S$\to$1 swaps fidelity and purity: purity needs no target state, fidelity does.
\item[\textbf{(B)}] Q$\to$1, R$\to$3: the parity is estimated from measured counts, without the density matrix; the fidelity needs $\rho$.
\item[\textbf{(D)}] P$\to$3, Q$\to$2 swaps the bases: phase errors are seen by the $X$-basis parity, bit flips by the $Z$-basis population.
\end{itemize}
\textbf{\color{brand}Key takeaway.} Population sees bit flips, parity sees phase flips, fidelity needs $\rho$.
\end{solbox}
\qsep

\needspace{13\baselineskip}
\begin{qitem}{57}
\qmeta{Assertion--Reason}{$(1-p)^G$: bound, not exact value}
\textbf{Assertion (A):} For depolarizing noise, the GHZ state fidelity $\braket{\mathrm{GHZ}|\rho|\mathrm{GHZ}}$ is \emph{exactly} equal to the
no-error probability $(1 - p)^G$.\\[1mm]
\textbf{Reason (R):} For Pauli noise, a run without any error produces the ideal state, and no error occurs with probability $(1 - p)^G$.
\begin{itemize}[leftmargin=2.6em, labelsep=0.4em, itemsep=2pt, topsep=3pt]
\item[\optlabel{A}] Both Assertion (A) and Reason (R) are true, and R is the correct explanation of A
\item[\optlabel{B}] Both A and R are true, but R is \emph{not} the correct explanation of A
\item[\optlabel{C}] A is true, but R is false
\item[\correctlabel{D}] A is false, but R is true
\end{itemize}
\end{qitem}
\begin{solbox}{D}
\textbf{\color{brand}Correct answer.} \textcolor{okgreen}{\textbf{(D)}} A is false, but R is true

\textbf{\color{brand}Concepts required.} Lower bound vs exact value of the fidelity under Pauli noise.

\textbf{\color{brand}Relevant formulae.} Pauli channel $= (1-p)\,\mathrm{id} + p\,(\text{Pauli errors})$; $F \ge (1-p)^G$; $Z_iZ_j\ket{\mathrm{GHZ}} = \ket{\mathrm{GHZ}}$.

\textbf{\color{brand}Step-by-step explanation.} \textbf{R is true}: a Pauli channel is a mixture ``ideal with probability $1 - p$, Pauli error otherwise''; the error-free branch alone contributes
$(1-p)^G$ to the fidelity, which makes $(1-p)^G$ a \emph{lower bound}. \textbf{A is false}: some error branches also leave the GHZ state unchanged ---
e.g.\ $Z\otimes Z$ after a CNOT ($ZZ\ket{\mathrm{GHZ}} = \ket{\mathrm{GHZ}}$) --- so the fidelity is strictly larger. The lecture's 6-qubit density-matrix run gives
$F = 0.9492$ at $p = 0.01$, while $(0.99)^6 = 0.9415$. (For non-Pauli noise such as amplitude damping the notebook does not claim the bound at all.) Answer: (D).

\textbf{\color{brand}Why the other options are incorrect.}
\begin{itemize}[leftmargin=2.4em, itemsep=1pt, topsep=2pt]
\item[\textbf{(A)}] requires A to be true, but errors that leave the GHZ state invariant add to the fidelity, so $F > (1-p)^G$.
\item[\textbf{(B)}] also requires A to be true (with R not explaining it); A is false because the bound is not tight.
\item[\textbf{(C)}] requires R to be false, but R is exactly the argument that makes $(1-p)^G$ a lower bound.
\end{itemize}
\textbf{\color{brand}Key takeaway.} $(1-p)^G$ is a lower bound for Pauli noise, not the fidelity itself.
\end{solbox}
\qsep

\needspace{13\baselineskip}
\begin{qitem}{58}
\qmeta{Statement I / II}{Density-matrix path in CUDA-Q 0.16}
\textbf{Statement I:} In CUDA-Q 0.16 the density-matrix path of the notebook calls \texttt{set\_noise} before \texttt{get\_state},
because \texttt{get\_state} has no \texttt{noise\_model} argument.\\[1mm]
\textbf{Statement II:} The density-matrix method needs $16\cdot4^N$ bytes, so $N = 14$ needs about $4.3$ GB.
\begin{itemize}[leftmargin=2.6em, labelsep=0.4em, itemsep=2pt, topsep=3pt]
\item[\correctlabel{A}] Both Statement I and Statement II are true
\item[\optlabel{B}] Both Statement I and Statement II are false
\item[\optlabel{C}] Statement I is true, but Statement II is false
\item[\optlabel{D}] Statement I is false, but Statement II is true
\end{itemize}
\end{qitem}
\begin{solbox}{A}
\textbf{\color{brand}Correct answer.} \textcolor{okgreen}{\textbf{(A)}} Both Statement I and Statement II are true

\textbf{\color{brand}Concepts required.} The notebook's density-matrix path in CUDA-Q 0.16.

\textbf{\color{brand}Relevant formulae.} $16\cdot4^{14} = 4.29\times10^9$ B.

\textbf{\color{brand}Step-by-step explanation.} \textbf{I is true}: the notebook says so in a comment and uses \texttt{set\_noise} / \texttt{unset\_noise} in a \texttt{try/finally} block (with a deprecation
warning filtered). \textbf{II is true}: $16\times4^{14} = 4.29\times10^9$ bytes. Answer: (A).

\textbf{\color{brand}Why the other options are incorrect.}
\begin{itemize}[leftmargin=2.4em, itemsep=1pt, topsep=2pt]
\item[\textbf{(B)}] requires both statements to be false; both are true.
\item[\textbf{(C)}] requires II to be false, but $16\cdot4^{14} = 4.29\times10^9$ bytes.
\item[\textbf{(D)}] requires I to be false, but the notebook does call \texttt{set\_noise} because \texttt{get\_state} has no noise argument in 0.16.
\end{itemize}
\textbf{\color{brand}Key takeaway.} \texttt{set\_noise} + \texttt{get\_state}; memory $4^N$.
\end{solbox}
\qsep

\needspace{20\baselineskip}
\begin{passagebox}{Paragraph for Questions 59 and 60}
The previously recorded run of \texttt{Noise\_Simulation.ipynb} simulates the 15-qubit GHZ circuit ($H$ on qubit 0 and a chain of
14 CNOTs) with depolarizing noise after every gate on the stabilizer target \texttt{stim} (10\,000 shots, seed 42). Part of its table:
\begin{center}\small
\begin{tabular}{rcccr}\toprule
$p$ & population $P(0\cdots0) + P(1\cdots1)$ & parity $\langle X^{\otimes15}\rangle$ & $(1-p)^G$ & distinct outcomes\\ \midrule
0.01 & 0.8936 & 0.6910 & 0.8601 & 96\\
0.1 & 0.3232 & 0.0160 & 0.2059 & 1165\\ \bottomrule
\end{tabular}
\end{center}
\end{passagebox}
\needspace{13\baselineskip}
\begin{qitem}{59}
\qmeta{Paragraph based}{Recorded run: leaked probability}
The probability mass found \emph{outside} the two ideal GHZ outcomes at $p = 0.01$ and at $p = 0.1$ is
\par\smallskip\noindent\begin{tabularx}{\linewidth}{@{}*{2}{L}@{}}
\optlabel{A}\,$0.1399$ and $0.7941$ & \correctlabel{B}\,$0.1064$ and $0.6768$ \\[3pt]
\optlabel{C}\,$0.3090$ and $0.9840$ & \optlabel{D}\,$0.8936$ and $0.3232$
\end{tabularx}
\end{qitem}
\begin{solbox}{B}
\textbf{\color{brand}Correct answer.} \textcolor{okgreen}{\textbf{(B)}} $0.1064$ and $0.6768$

\textbf{\color{brand}Concepts required.} Leaked probability vs error probability.

\textbf{\color{brand}Relevant formulae.} leak $= 1 - $ population; $P(\text{some error}) = 1 - (1-p)^G$.

\textbf{\color{brand}Step-by-step explanation.} It is $1 - \text{population}$: $1 - 0.8936 = 0.1064$ and $1 - 0.3232 = 0.6768$ (this also equals the total variation distance to the ideal distribution
for a balanced population). $1 - (1-p)^G$ ($0.1399$, $0.7941$) is the probability that \emph{some} error occurred --- larger, because many errors (e.g.\ $Z$ errors) do
not move probability out of the two peaks. $1 - $ parity measures phase damage.

\textbf{\color{brand}Why the other options are incorrect.}
\begin{itemize}[leftmargin=2.4em, itemsep=1pt, topsep=2pt]
\item[\textbf{(A)}] $1 - (1-p)^G$ is the probability that \emph{some} error occurred; many errors ($Z$-type) do not move probability out of the two peaks.
\item[\textbf{(C)}] $1 - $ parity measures the phase damage seen in the $X$ basis, not the leaked $Z$-basis probability.
\item[\textbf{(D)}] $0.8936$ and $0.3232$ are the populations themselves, i.e.\ the mass \emph{inside} the peaks.
\end{itemize}
\textbf{\color{brand}Key takeaway.} Not every error leaves the ideal outcomes.
\end{solbox}
\qsep

\needspace{13\baselineskip}
\begin{qitem}{60}
\qmeta{Paragraph based}{Recorded run: population vs parity}
Why is the population at $p = 0.01$ ($0.8936$) \emph{larger} than $(1 - p)^G = 0.8601$, while the parity ($0.6910$) is smaller?
\begin{itemize}[leftmargin=2.6em, labelsep=0.4em, itemsep=2pt, topsep=3pt]
\item[\optlabel{A}] $(1-p)^G$ is an upper bound on the population, so the population must lie below it --- the table contains a mistake
\item[\correctlabel{B}] $(1-p)^G$ counts every error, but $Z$-type errors leave the $Z$-basis population intact; the parity is also sensitive to those phase errors (and to Hadamard noise in the $X$-basis circuit)
\item[\optlabel{C}] the stabilizer simulator under-reports errors because it cannot simulate depolarizing noise
\item[\optlabel{D}] 10\,000 shots are too few: with more shots the population would fall to exactly $0.8601$
\end{itemize}
\end{qitem}
\begin{solbox}{B}
\textbf{\color{brand}Correct answer.} \textcolor{okgreen}{\textbf{(B)}} $(1-p)^G$ counts every error, but $Z$-type errors leave the $Z$-basis population intact; the parity is also sensitive to those phase errors (and to Hadamard noise in the $X$-basis circuit)

\textbf{\color{brand}Concepts required.} Interpreting population, parity and bound together.

\textbf{\color{brand}Relevant formulae.} Pauli noise: $F \ge (1-p)^G$; $Z$ errors invisible to $Z$-basis populations.

\textbf{\color{brand}Step-by-step explanation.} $(1-p)^G$ is the weight of the error-free branch, a \emph{lower} bound for Pauli noise. Among the depolarizing errors, $Z$-type errors (e.g.\ $Z$ on one qubit, or $ZZ$ after a CNOT) keep the outcome in $\{0\cdots0, 1\cdots1\}$ --- so does an $X$ error on qubit 0 right after the Hadamard, which the CNOT chain copies to every qubit ---
so the population exceeds the bound. The parity measurement
needs extra Hadamards (also noisy), and $Z$-type errors flip its sign, so it is the more demanding test. The statistical error at 10\,000 shots is only
$\sqrt{0.89\cdot0.11/10^4} \approx 0.003$; \texttt{stim} simulates Pauli channels exactly.

\textbf{\color{brand}Why the other options are incorrect.}
\begin{itemize}[leftmargin=2.4em, itemsep=1pt, topsep=2pt]
\item[\textbf{(A)}] $(1-p)^G$ is a lower bound, not an upper one.
\item[\textbf{(C)}] \texttt{stim} simulates Pauli (incl.\ depolarizing) noise exactly.
\item[\textbf{(D)}] Shot noise is about $\pm0.003$, far below the difference.
\end{itemize}
\textbf{\color{brand}Key takeaway.} Use several figures of merit; each sees different errors.
\end{solbox}
\qsep

\newpage

\begin{sectionbanner}SECTION D \enspace$\cdot$\enspace Advanced Noise Analysis and Implementation \hfill Q.61--80\end{sectionbanner}

\needspace{13\baselineskip}
\begin{qitem}{61}
\qmeta{Single-correct MCQ}{Process vs average gate fidelity}
For a single-qubit depolarizing channel (CUDA-Q convention: Kraus operators $\sqrt{1-p}\,I$, $\sqrt{p/3}\,X, Y, Z$) the process (entanglement) fidelity is
$F_{\text{pro}} = 1 - p$. Using $F_{\text{avg}} = \frac{dF_{\text{pro}} + 1}{d + 1}$, the average gate fidelity for $p = 0.003$ is
\par\smallskip\noindent\begin{tabularx}{\linewidth}{@{}*{4}{L}@{}}
\correctlabel{A}\,$0.9980$ & \optlabel{B}\,$0.9990$ & \optlabel{C}\,$0.9970$ & \optlabel{D}\,$0.99775$
\end{tabularx}
\end{qitem}
\begin{solbox}{A}
\textbf{\color{brand}Correct answer.} \textcolor{okgreen}{\textbf{(A)}} $0.9980$

\textbf{\color{brand}Concepts required.} Process vs average gate fidelity.

\textbf{\color{brand}Relevant formulae.} $F_{\text{avg}} = \frac{dF_{\text{pro}} + 1}{d + 1}$.

\textbf{\color{brand}Step-by-step explanation.} $d = 2$: $F_{\text{avg}} = \frac{2(1 - p) + 1}{3} = 1 - \frac{2p}{3} = 1 - 0.002 = 0.9980$; the average gate infidelity (``error rate'') is $r = 0.002$.
$F_{\text{pro}} = 0.997$ itself is lower because it does not average over input states. $1 - \frac{3p}{4}$ mixes up the $\lambda = \frac{4p}{3}$ convention.

\textbf{\color{brand}Why the other options are incorrect.}
\begin{itemize}[leftmargin=2.4em, itemsep=1pt, topsep=2pt]
\item[\textbf{(B)}] $0.9990 = 1 - \frac p3$ counts only one of the three Pauli errors.
\item[\textbf{(C)}] $0.9970 = 1 - p$ is the process (entanglement) fidelity, which does not average over input states.
\item[\textbf{(D)}] $0.99775 = 1 - \frac{3p}{4}$ uses the convention $\rho \mapsto (1-p)\rho + p\frac I2$, not CUDA-Q's.
\end{itemize}
\textbf{\color{brand}Key takeaway.} Single-qubit depolarizing: $r = \frac{2p}{3}$.
\end{solbox}
\qsep

\needspace{13\baselineskip}
\begin{qitem}{62}
\qmeta{Single-correct MCQ}{Randomized benchmarking}
Single-qubit randomized benchmarking fits the survival probability to $A\alpha^m + B$ ($m$ = sequence length) and finds $\alpha = 0.99$. The error per Clifford
$r = \frac{(d - 1)(1 - \alpha)}{d}$ is
\par\smallskip\noindent\begin{tabularx}{\linewidth}{@{}*{4}{L}@{}}
\optlabel{A}\,$0.010$ & \optlabel{B}\,$0.0067$ & \optlabel{C}\,$0.0025$ & \correctlabel{D}\,$0.005$
\end{tabularx}
\end{qitem}
\begin{solbox}{D}
\textbf{\color{brand}Correct answer.} \textcolor{okgreen}{\textbf{(D)}} $0.005$

\textbf{\color{brand}Concepts required.} Randomized benchmarking decay.

\textbf{\color{brand}Relevant formulae.} $r = \frac{(d-1)(1-\alpha)}{d}$.

\textbf{\color{brand}Step-by-step explanation.} $d = 2$: $r = \frac{1\cdot(1 - 0.99)}{2} = 0.005$. RB is robust to state-preparation and measurement errors (absorbed in $A$, $B$) and reports an
\emph{average} error per gate under twirled (effectively depolarizing) noise. $1 - \alpha = 0.01$ forgets the factor $\frac{d-1}{d}$.

\textbf{\color{brand}Why the other options are incorrect.}
\begin{itemize}[leftmargin=2.4em, itemsep=1pt, topsep=2pt]
\item[\textbf{(A)}] $0.010 = 1 - \alpha$ forgets the factor $\frac{d-1}{d} = \frac12$.
\item[\textbf{(B)}] $0.0067 = \frac23(1 - \alpha)$ uses $d = 3$ (a qutrit).
\item[\textbf{(C)}] $0.0025 = \frac14(1-\alpha)$ divides by $d^2 = 4$ instead of using $\frac{d-1}{d}$.
\end{itemize}
\textbf{\color{brand}Key takeaway.} RB measures average error per gate, robust to SPAM.
\end{solbox}
\qsep

\needspace{13\baselineskip}
\begin{qitem}{63}
\qmeta{Single-correct MCQ}{Linear ZNE}
Zero-noise extrapolation measures $E(1) = 0.80$ at the native noise and $E(2) = 0.64$ at doubled noise. The \emph{linear} extrapolation to $\lambda = 0$ gives
\par\smallskip\noindent\begin{tabularx}{\linewidth}{@{}*{4}{L}@{}}
\optlabel{A}\,$0.88$ & \optlabel{B}\,$1.00$ & \optlabel{C}\,$0.80$ & \correctlabel{D}\,$0.96$
\end{tabularx}
\end{qitem}
\begin{solbox}{D}
\textbf{\color{brand}Correct answer.} \textcolor{okgreen}{\textbf{(D)}} $0.96$

\textbf{\color{brand}Concepts required.} Linear zero-noise extrapolation.

\textbf{\color{brand}Relevant formulae.} $E(0) = 2E(1) - E(2)$.

\textbf{\color{brand}Step-by-step explanation.} Line through $(1, 0.80)$ and $(2, 0.64)$: slope $-0.16$, so $E(0) = 0.80 + 0.16 = 0.96$ (formula $2E(1) - E(2)$).
An exponential model $E = ae^{-b\lambda}$ would give $E(1)^2/E(2) = 1.00$: the result depends on the assumed model --- ZNE is a biased estimator.

\textbf{\color{brand}Why the other options are incorrect.}
\begin{itemize}[leftmargin=2.4em, itemsep=1pt, topsep=2pt]
\item[\textbf{(A)}] $0.88$ interpolates to $\lambda = 1.5$ instead of extrapolating to $\lambda = 0$.
\item[\textbf{(B)}] $1.00 = E(1)^2/E(2)$ is the exponential extrapolation, a different model from the requested linear one.
\item[\textbf{(C)}] $0.80$ is the unmitigated value at $\lambda = 1$.
\end{itemize}
\textbf{\color{brand}Key takeaway.} ZNE result depends on the extrapolation model.
\end{solbox}
\qsep

\needspace{13\baselineskip}
\begin{qitem}{64}
\qmeta{Single-correct MCQ}{Richardson ZNE}
ZNE with three noise levels gives $E(1) = 0.80$, $E(2) = 0.65$, $E(3) = 0.53$. The quadratic (Richardson) extrapolation $E(0) = 3E(1) - 3E(2) + E(3)$ equals
\par\smallskip\noindent\begin{tabularx}{\linewidth}{@{}*{4}{L}@{}}
\optlabel{A}\,$0.95$ & \optlabel{B}\,$0.985$ & \correctlabel{C}\,$0.98$ & \optlabel{D}\,$0.53$
\end{tabularx}
\end{qitem}
\begin{solbox}{C}
\textbf{\color{brand}Correct answer.} \textcolor{okgreen}{\textbf{(C)}} $0.98$

\textbf{\color{brand}Concepts required.} Richardson extrapolation with three points.

\textbf{\color{brand}Relevant formulae.} $E(0) = 3E(1) - 3E(2) + E(3)$ for $\lambda = 1, 2, 3$.

\textbf{\color{brand}Step-by-step explanation.} $3(0.80) - 3(0.65) + 0.53 = 2.40 - 1.95 + 0.53 = 0.98$. \emph{Check}: the parabola through the three points has second difference
$0.53 - 2(0.65) + 0.80 = 0.03$, so $E(\lambda) = 0.98 - 0.195\lambda + 0.015\lambda^2$. The linear estimate from the first two points is $0.95$; an exponential fit through $\lambda = 1, 2$ gives $0.8^2/0.65 = 0.985$.
The coefficients $(3, -3, 1)$ amplify the statistical errors: $\sqrt{9 + 9 + 1} \approx 4.4$ times the error of one point.

\textbf{\color{brand}Why the other options are incorrect.}
\begin{itemize}[leftmargin=2.4em, itemsep=1pt, topsep=2pt]
\item[\textbf{(A)}] $0.95 = 2E(1) - E(2)$ is the two-point linear estimate, not the three-point Richardson one.
\item[\textbf{(B)}] $0.985 = E(1)^2/E(2)$ is a two-point exponential fit, not the quadratic Richardson estimate.
\item[\textbf{(D)}] $0.53 = E(3)$ is the most noisy measurement, not an extrapolation.
\end{itemize}
\textbf{\color{brand}Key takeaway.} Higher-order extrapolation: less bias, more variance.
\end{solbox}
\qsep

\needspace{13\baselineskip}
\begin{qitem}{65}
\qmeta{Single-correct MCQ}{Gate folding scale factor}
A circuit has 10 gates of equal error. To amplify noise, 4 of them are ``folded'' ($G \to GG^\dagger G$), the others are left unchanged. Assuming the noise grows with
the number of gates, the effective noise scale factor is
\par\smallskip\noindent\begin{tabularx}{\linewidth}{@{}*{4}{L}@{}}
\optlabel{A}\,$2.2$ & \optlabel{B}\,$1.4$ & \correctlabel{C}\,$1.8$ & \optlabel{D}\,$3.0$
\end{tabularx}
\end{qitem}
\begin{solbox}{C}
\textbf{\color{brand}Correct answer.} \textcolor{okgreen}{\textbf{(C)}} $1.8$

\textbf{\color{brand}Concepts required.} Noise amplification by gate folding.

\textbf{\color{brand}Relevant formulae.} $\lambda = \frac{\text{gates after folding}}{\text{original gates}}$.

\textbf{\color{brand}Step-by-step explanation.} Each folded gate becomes 3 gates: $4\times3 + 6\times1 = 18$ gates instead of 10, so $\lambda = 18/10 = 1.8$. Folding all gates gives $\lambda = 3$.
Folding keeps the ideal unitary unchanged ($GG^\dagger G = G$) while multiplying the noise --- the hardware analogue of the lecture's ``scale $p$'' experiment.

\textbf{\color{brand}Why the other options are incorrect.}
\begin{itemize}[leftmargin=2.4em, itemsep=1pt, topsep=2pt]
\item[\textbf{(A)}] $2.2 = 22/10$ counts three extra gates per fold.
\item[\textbf{(B)}] $1.4 = 14/10$ counts one extra gate per fold; $GG^\dagger G$ adds two.
\item[\textbf{(D)}] $3.0$ assumes all ten gates are folded ($30/10$); only four are folded here.
\end{itemize}
\textbf{\color{brand}Key takeaway.} $GG^\dagger G$ triples the noise of $G$ without changing its action.
\end{solbox}
\qsep

\needspace{13\baselineskip}
\begin{qitem}{66}
\qmeta{Single-correct MCQ}{Readout mitigation}
Readout calibration of one qubit gives $P(\text{read }0|0) = 0.98$ and $P(\text{read }0|1) = 0.05$. An experiment reads $0$ with frequency $0.60$. The mitigated
probability of the true state $\ket0$ is about
\par\smallskip\noindent\begin{tabularx}{\linewidth}{@{}*{4}{L}@{}}
\optlabel{A}\,$0.604$ & \correctlabel{B}\,$0.591$ & \optlabel{C}\,$0.600$ & \optlabel{D}\,$0.612$
\end{tabularx}
\end{qitem}
\begin{solbox}{B}
\textbf{\color{brand}Correct answer.} \textcolor{okgreen}{\textbf{(B)}} $0.591$

\textbf{\color{brand}Concepts required.} Inverting a confusion matrix.

\textbf{\color{brand}Relevant formulae.} $\mathbf q = A\mathbf p \Rightarrow \mathbf p = A^{-1}\mathbf q$.

\textbf{\color{brand}Step-by-step explanation.} Solve $\begin{psmallmatrix}0.98 & 0.05\\ 0.02 & 0.95\end{psmallmatrix}\begin{psmallmatrix}p_0\\ p_1\end{psmallmatrix} = \begin{psmallmatrix}0.60\\ 0.40\end{psmallmatrix}$ with
$p_0 + p_1 = 1$: $0.98p_0 + 0.05(1 - p_0) = 0.60 \Rightarrow 0.93p_0 = 0.55 \Rightarrow p_0 = 0.5914$. $0.6/0.98 = 0.612$ ignores the $1\to0$ errors;
$0.604$ comes from the transposed confusion matrix.

\textbf{\color{brand}Why the other options are incorrect.}
\begin{itemize}[leftmargin=2.4em, itemsep=1pt, topsep=2pt]
\item[\textbf{(A)}] $0.604$ results from the transposed confusion matrix ($0.98p_0 + 0.02p_1 = 0.60$).
\item[\textbf{(C)}] $0.600$ is the raw (unmitigated) frequency; it still contains the misassigned $\ket1$ shots.
\item[\textbf{(D)}] $0.612 = 0.60/0.98$ corrects only the $0 \to 1$ errors and ignores the $1 \to 0$ errors.
\end{itemize}
\textbf{\color{brand}Key takeaway.} Columns of $A$ = true states, rows = read outcomes.
\end{solbox}
\qsep

\needspace{13\baselineskip}
\begin{qitem}{67}
\qmeta{Single-correct MCQ}{Limits of readout mitigation}
Which statement about readout-error mitigation by inverting the calibrated confusion matrix $A$ is correct?
\begin{itemize}[leftmargin=2.6em, labelsep=0.4em, itemsep=2pt, topsep=3pt]
\item[\optlabel{A}] it reduces the number of shots needed, because $A^{-1}$ removes the shot noise
\item[\optlabel{B}] it requires quantum error correction to estimate $A$
\item[\optlabel{C}] it also corrects gate errors that happened before the measurement
\item[\correctlabel{D}] $A^{-1}\mathbf q$ can contain small negative ``quasi-probabilities'', and the full $2^N\times2^N$ matrix becomes impractical for many qubits
\end{itemize}
\end{qitem}
\begin{solbox}{D}
\textbf{\color{brand}Correct answer.} \textcolor{okgreen}{\textbf{(D)}} $A^{-1}\mathbf q$ can contain small negative ``quasi-probabilities'', and the full $2^N\times2^N$ matrix becomes impractical for many qubits

\textbf{\color{brand}Concepts required.} Limits of readout mitigation.

\textbf{\color{brand}Relevant formulae.} $\mathbf p = A^{-1}\mathbf q$; variance grows with $\|A^{-1}\|$.

\textbf{\color{brand}Step-by-step explanation.} Inversion only undoes the classical misassignment of outcomes. Shot noise in $\mathbf q$ propagates (and is amplified) through $A^{-1}$, which can produce slightly
negative entries --- the lecture's 3-qubit example gave a minimum of $-0.0009$; they are usually clipped or handled by constrained least squares. $A$ is measured
by preparing and reading basis states; for many qubits one assumes tensor-product or correlated sparse models instead of the full $4^N$-entry matrix.

\textbf{\color{brand}Why the other options are incorrect.}
\begin{itemize}[leftmargin=2.4em, itemsep=1pt, topsep=2pt]
\item[\textbf{(A)}] Inversion amplifies the shot noise ($A^{-1}$ has norm $> 1$); it cannot reduce the number of shots needed.
\item[\textbf{(B)}] $A$ is estimated by preparing basis states and measuring them --- no error correction is involved.
\item[\textbf{(C)}] Errors in the gates before the measurement change the state; the confusion matrix only models the classical readout step.
\end{itemize}
\textbf{\color{brand}Key takeaway.} Readout mitigation is post-processing; quasi-probabilities can be negative.
\end{solbox}
\qsep

\needspace{13\baselineskip}
\begin{qitem}{68}
\qmeta{Single-correct MCQ}{Ramsey vs Hahn echo}
A qubit in $\ket+$ idles for $t$ with a random but \emph{static} detuning $\delta \sim \mathcal N(0, \sigma^2)$, $\sigma = 1$ rad/$\mu$s, so it acquires the phase $\delta t$. At
$t = 2\,\mu$s the average coherence $\langle X\rangle$ without and with a Hahn echo ($X$ pulse at $t/2$ and again at the end) is
\par\smallskip\noindent\begin{tabularx}{\linewidth}{@{}*{4}{L}@{}}
\optlabel{A}\,$0.135$ and $0.135$ & \optlabel{B}\,$0.368$ and $1.0$ & \optlabel{C}\,$0.607$ and $1.0$ & \correctlabel{D}\,$0.135$ and $1.0$
\end{tabularx}
\end{qitem}
\begin{solbox}{D}
\textbf{\color{brand}Correct answer.} \textcolor{okgreen}{\textbf{(D)}} $0.135$ and $1.0$

\textbf{\color{brand}Concepts required.} Ramsey decay vs spin echo for static noise.

\textbf{\color{brand}Relevant formulae.} $\mathbb E\cos(\delta t) = e^{-\sigma^2t^2/2}$; $XR_Z(a)X = R_Z(-a)$.

\textbf{\color{brand}Step-by-step explanation.} Without echo: $\langle X\rangle = \mathbb E\cos(\delta t) = e^{-\sigma^2t^2/2} = e^{-2} = 0.135$ (Gaussian characteristic function).
With echo: $R_Z(\varphi/2)\,X\,R_Z(\varphi/2) = X$ because $XR_Z(a)X = R_Z(-a)$; the two halves cancel for \emph{every} static $\delta$, so $\langle X\rangle = 1$.
($e^{-1} = 0.368$ uses $e^{-\sigma^2t^2/4}$; $0.607 = e^{-1/2}$ uses $t = 1$.)

\textbf{\color{brand}Why the other options are incorrect.}
\begin{itemize}[leftmargin=2.4em, itemsep=1pt, topsep=2pt]
\item[\textbf{(A)}] $0.135$ with echo ignores that the echo cancels a static phase exactly.
\item[\textbf{(B)}] $0.368 = e^{-1}$ uses $e^{-\sigma^2t^2/4}$; the Gaussian average is $\mathbb E\cos(\delta t) = e^{-\sigma^2t^2/2}$.
\item[\textbf{(C)}] $0.607 = e^{-1/2}$ is the free decay at $t = 1\,\mu$s, not $2\,\mu$s.
\end{itemize}
\textbf{\color{brand}Key takeaway.} Echoes cancel slow noise.
\end{solbox}
\qsep

\needspace{13\baselineskip}
\begin{qitem}{69}
\qmeta{Single-correct MCQ}{DD and Markovian noise}
Now the dephasing is Markovian: an independent phase-flip channel with $p = 0.02$ acts after each of 20 idle steps on $\ket+$. With a Hahn echo inserted at step 10,
$\langle X\rangle$ at the end is
\par\smallskip\noindent\begin{tabularx}{\linewidth}{@{}*{2}{L}@{}}
\optlabel{A}\,$0$ & \optlabel{B}\,$1.000$ \\[3pt]
\correctlabel{C}\,$0.442$, the same as without echo & \optlabel{D}\,$0.668$
\end{tabularx}
\end{qitem}
\begin{solbox}{C}
\textbf{\color{brand}Correct answer.} \textcolor{okgreen}{\textbf{(C)}} $0.442$, the same as without echo

\textbf{\color{brand}Concepts required.} Markovian noise and dynamical decoupling.

\textbf{\color{brand}Relevant formulae.} $\langle X\rangle = (1-2p)^n$; $X$ commutes with the phase-flip channel.

\textbf{\color{brand}Step-by-step explanation.} Each step multiplies $\langle X\rangle$ by $1 - 2p = 0.96$: $0.96^{20} = 0.442$. The $X$ pulses commute with the phase-flip channel up to the sign that the second
$X$ undoes, and independent random flips in the two halves do not cancel. The lecture's density-matrix simulation gives $0.4420$ both with and without echo:
dynamical decoupling only fights noise that is slow (correlated in time) compared with the pulse spacing.

\textbf{\color{brand}Why the other options are incorrect.}
\begin{itemize}[leftmargin=2.4em, itemsep=1pt, topsep=2pt]
\item[\textbf{(A)}] $0$: the coherence decays to $0.96^{20} = 0.44$, it does not vanish.
\item[\textbf{(B)}] $1.000$ assumes the echo cancels the phase flips, but independent random flips in the two halves do not cancel.
\item[\textbf{(D)}] $0.668 = 0.98^{20}$ uses $1 - p$ instead of $1 - 2p$ per step.
\end{itemize}
\textbf{\color{brand}Key takeaway.} DD does not help against white (uncorrelated) noise.
\end{solbox}
\qsep

\needspace{13\baselineskip}
\begin{qitem}{70}
\qmeta{Single-correct MCQ}{3-qubit bit-flip code}
In the 3-qubit bit-flip code each qubit flips independently with probability $p$, and encoding and correction are perfect. The logical error probability at $p = 0.1$, and
the value of $p$ above which encoding is worse than an unencoded qubit, are
\par\smallskip\noindent\begin{tabularx}{\linewidth}{@{}*{2}{L}@{}}
\correctlabel{A}\,$0.028$ and $p = 0.5$ & \optlabel{B}\,$0.010$ and $p = 0.5$ \\[3pt]
\optlabel{C}\,$0.028$ and $p = 1/3$ & \optlabel{D}\,$0.030$ and $p = 0.5$
\end{tabularx}
\end{qitem}
\begin{solbox}{A}
\textbf{\color{brand}Correct answer.} \textcolor{okgreen}{\textbf{(A)}} $0.028$ and $p = 0.5$

\textbf{\color{brand}Concepts required.} Repetition code and break-even.

\textbf{\color{brand}Relevant formulae.} $P_L = 3p^2 - 2p^3$; $P_L < p \Leftrightarrow p < \frac12$.

\textbf{\color{brand}Step-by-step explanation.} Majority vote fails when 2 or 3 qubits flip: $P_L = 3p^2(1-p) + p^3 = 3p^2 - 2p^3 = 0.03 - 0.002 = 0.028$ (the lecture's CUDA-Q density-matrix simulation
reproduces this formula to $10^{-16}$). $P_L < p \Leftrightarrow 3p - 2p^2 < 1 \Leftrightarrow (2p - 1)(p - 1) > 0 \Leftrightarrow p < \frac12$. $3p^2 = 0.030$ drops the
$-2p^3$ term.

\textbf{\color{brand}Why the other options are incorrect.}
\begin{itemize}[leftmargin=2.4em, itemsep=1pt, topsep=2pt]
\item[\textbf{(B)}] $0.010 = p^2$ forgets the three ways of choosing two flipped qubits.
\item[\textbf{(C)}] Break-even at $p = \frac13$ is wrong: $3p^2 - 2p^3 < p \Leftrightarrow p < \frac12$.
\item[\textbf{(D)}] $0.030 = 3p^2$ drops the $-2p^3$ term (the probability of three flips is counted three times).
\end{itemize}
\textbf{\color{brand}Key takeaway.} Encoding helps below a threshold.
\end{solbox}
\qsep

\needspace{13\baselineskip}
\begin{qitem}{71}
\qmeta{Single-correct MCQ}{Surface-code scaling}
A surface-code memory has logical error rate $p_L \approx A(p/p_{\text{th}})^{(d+1)/2}$ for code distance $d$. With $p/p_{\text{th}} = 0.1$, increasing the distance from
$d = 3$ to $d = 5$ reduces $p_L$ by a factor of
\par\smallskip\noindent\begin{tabularx}{\linewidth}{@{}*{2}{L}@{}}
\optlabel{A}\,none: $p_L$ increases by a factor of 10 & \correctlabel{B}\,$10$ \\[3pt]
\optlabel{C}\,$2$ & \optlabel{D}\,$100$
\end{tabularx}
\end{qitem}
\begin{solbox}{B}
\textbf{\color{brand}Correct answer.} \textcolor{okgreen}{\textbf{(B)}} $10$

\textbf{\color{brand}Concepts required.} Below-threshold scaling of surface codes.

\textbf{\color{brand}Relevant formulae.} $p_L \propto (p/p_{\text{th}})^{(d+1)/2}$.

\textbf{\color{brand}Step-by-step explanation.} The exponent goes from $2$ to $3$: $p_L$ is multiplied by $(p/p_{\text{th}})^{1} = 0.1$, i.e.\ reduced by a factor of 10. Below threshold ($p < p_{\text{th}}$) larger
codes suppress errors exponentially in $d$; above threshold they make things worse. QEC needs many physical qubits and fast feedback --- it is not
implemented in this repository.

\textbf{\color{brand}Why the other options are incorrect.}
\begin{itemize}[leftmargin=2.4em, itemsep=1pt, topsep=2pt]
\item[\textbf{(A)}] Below threshold ($p < p_{\text{th}}$) larger codes \emph{reduce} $p_L$; it increases only above threshold.
\item[\textbf{(C)}] $2$ ignores the power law; the reduction factor is $(p/p_{\text{th}})^{-1} = 10$.
\item[\textbf{(D)}] $100$ corresponds to an exponent increase of 2; going from $d = 3$ to $d = 5$ raises $\frac{d+1}{2}$ from 2 to 3, by 1.
\end{itemize}
\textbf{\color{brand}Key takeaway.} Each $+2$ in distance multiplies $p_L$ by $p/p_{\text{th}}$.
\end{solbox}
\qsep

\needspace{13\baselineskip}
\begin{qitem}{72}
\qmeta{Single-correct MCQ}{Coherent vs stochastic accumulation}
Each of $n = 60$ idle gates over-rotates coherently by $\varepsilon = 0.05$ rad about $x$. Compare with a stochastic bit flip of the same per-gate error probability
$p = \sin^2(\varepsilon/2)$. Starting from $\ket0$, the probabilities of finding $1$ after 60 gates are approximately
\par\smallskip\noindent\begin{tabularx}{\linewidth}{@{}*{2}{L}@{}}
\correctlabel{A}\,$0.995$ (coherent) and $0.036$ (stochastic) & \optlabel{B}\,$0.0375$ and $0.036$ \\[3pt]
\optlabel{C}\,$0.995$ and $0.995$ & \optlabel{D}\,$0.036$ and $0.036$
\end{tabularx}
\end{qitem}
\begin{solbox}{A}
\textbf{\color{brand}Correct answer.} \textcolor{okgreen}{\textbf{(A)}} $0.995$ (coherent) and $0.036$ (stochastic)

\textbf{\color{brand}Concepts required.} Coherent vs stochastic accumulation.

\textbf{\color{brand}Relevant formulae.} Coherent $\sin^2(n\varepsilon/2)$; stochastic $\frac12[1 - (1-2p)^n]$.

\textbf{\color{brand}Step-by-step explanation.} Coherent: the rotations add up, $R_X(\varepsilon)^{60} = R_X(3)$: $P(1) = \sin^2(1.5) = 0.995$. Stochastic: $p = \sin^2(0.025) = 6.25\times10^{-4}$ per gate and
$P(1) = \frac12[1 - (1 - 2p)^{60}] = \frac12(1 - 0.9278) = 0.036 \approx np$. Coherent errors grow like $n^2\varepsilon^2$ (amplitudes add), stochastic ones like $np$.
Both values were computed with CUDA-Q for the lecture. Randomized compiling/Pauli twirling turns the former into the latter.

\textbf{\color{brand}Why the other options are incorrect.}
\begin{itemize}[leftmargin=2.4em, itemsep=1pt, topsep=2pt]
\item[\textbf{(B)}] $np = 0.0375$ is the stochastic estimate, not the coherent one.
\item[\textbf{(C)}] Stochastic errors grow linearly, not like coherent ones.
\item[\textbf{(D)}] Coherent errors do not add like probabilities.
\end{itemize}
\textbf{\color{brand}Key takeaway.} Coherent errors can be far worse than their per-gate fidelity suggests.
\end{solbox}
\qsep

\needspace{13\baselineskip}
\begin{qitem}{73}
\qmeta{Single-correct MCQ}{Pauli twirling}
Which statement about Pauli twirling (randomized compiling) is correct?
\begin{itemize}[leftmargin=2.6em, labelsep=0.4em, itemsep=2pt, topsep=3pt]
\item[\optlabel{A}] it removes all errors, so the circuit becomes noiseless
\item[\correctlabel{B}] averaging a gate over random Pauli frames that cancel ideally turns its coherent error into a stochastic Pauli channel with the same average fidelity
\item[\optlabel{C}] it requires a density-matrix simulator and cannot be done on hardware
\item[\optlabel{D}] it converts stochastic Pauli noise into coherent errors to make them easier to correct
\end{itemize}
\end{qitem}
\begin{solbox}{B}
\textbf{\color{brand}Correct answer.} \textcolor{okgreen}{\textbf{(B)}} averaging a gate over random Pauli frames that cancel ideally turns its coherent error into a stochastic Pauli channel with the same average fidelity

\textbf{\color{brand}Concepts required.} Pauli twirling / randomized compiling.

\textbf{\color{brand}Relevant formulae.} $\mathcal E_{\text{twirl}}(\rho) = \frac1{|\mathcal P|}\sum_PP^\dagger\mathcal E(P\rho P^\dagger)P$ is a Pauli channel.

\textbf{\color{brand}Step-by-step explanation.} Inserting random Paulis $P$ before a gate and the compensating $P'$ after it leaves the ideal circuit unchanged; averaged over many random compilations the
effective channel is diagonal in the Pauli basis --- a Pauli (stochastic) channel. This prevents the quadratic, worst-case accumulation of coherent errors and makes
noise models such as those of the notebook (Pauli channels) realistic. It does not reduce the average error itself.

\textbf{\color{brand}Why the other options are incorrect.}
\begin{itemize}[leftmargin=2.4em, itemsep=1pt, topsep=2pt]
\item[\textbf{(A)}] Twirling reshapes the error into a Pauli channel but keeps its average fidelity; the circuit is not noiseless.
\item[\textbf{(C)}] Randomized compiling is a compilation technique used on hardware; it needs no simulator.
\item[\textbf{(D)}] Twirling works in the opposite direction: it turns coherent errors into stochastic Pauli noise.
\end{itemize}
\textbf{\color{brand}Key takeaway.} Twirling makes noise stochastic, and Pauli-noise simulation realistic.
\end{solbox}
\qsep

\needspace{13\baselineskip}
\begin{qitem}{74}
\qmeta{Single-correct MCQ}{Notebook: \texttt{auto} method choice}
The noise notebook runs on a CPU (no GPU visible) with the settings\\
\hspace*{1em}\texttt{CIRCUIT = "ghz"},\ \ \texttt{NOISE\_MODEL = "amplitude\_damping"},\ \ \texttt{SIMULATION\_METHOD = "auto"}.\\
Which method does it choose?
\begin{itemize}[leftmargin=2.6em, labelsep=0.4em, itemsep=2pt, topsep=3pt]
\item[\optlabel{A}] \texttt{stabilizer} (target \texttt{stim})
\item[\optlabel{B}] \texttt{trajectory} (target \texttt{nvidia})
\item[\optlabel{C}] \texttt{statevector} (target \texttt{qpp-cpu})
\item[\correctlabel{D}] \texttt{density\_matrix} (target \texttt{density-matrix-cpu})
\end{itemize}
\end{qitem}
\begin{solbox}{D}
\textbf{\color{brand}Correct answer.} \textcolor{okgreen}{\textbf{(D)}} \texttt{density\_matrix} (target \texttt{density-matrix-cpu})

\textbf{\color{brand}Concepts required.} The notebook's \texttt{auto} method rule.

\textbf{\color{brand}Relevant formulae.} GPU $\to$ trajectory; CPU + GHZ + Pauli $\to$ stabilizer; else density matrix.

\textbf{\color{brand}Step-by-step explanation.} The \texttt{"auto"} rule: GPU visible $\to$ \texttt{trajectory}; otherwise \texttt{stabilizer} only for GHZ \emph{with Pauli noise}; otherwise \texttt{density\_matrix}.
Amplitude damping is not a Pauli channel, so \texttt{stim} would silently ignore it --- the notebook even raises an error if \texttt{stabilizer} is forced with it.
\texttt{statevector} is refused because noiseless simulators ignore noise models.

\textbf{\color{brand}Why the other options are incorrect.}
\begin{itemize}[leftmargin=2.4em, itemsep=1pt, topsep=2pt]
\item[\textbf{(A)}] \texttt{stim} cannot represent amplitude damping; the \texttt{auto} rule chooses it only for Pauli noise.
\item[\textbf{(B)}] \texttt{trajectory} is chosen only when a GPU is visible; this case is a CPU.
\item[\textbf{(C)}] The notebook refuses \texttt{statevector}: noiseless simulators ignore noise models.
\end{itemize}
\textbf{\color{brand}Key takeaway.} Method follows from device, circuit and noise type.
\end{solbox}
\qsep

\needspace{13\baselineskip}
\begin{qitem}{75}
\qmeta{Matching lists}{Mitigation, DD, QEC}
Match each technique (List-I) with its description (List-II).
\par\smallskip\noindent{\renewcommand{\arraystretch}{1.35}%
\begin{tabularx}{\linewidth}{@{}c>{\raggedright\arraybackslash}p{0.33\linewidth}@{\hspace{1.2em}}cL@{}}
\toprule
\multicolumn{2}{@{}l}{\textbf{List-I}} & \multicolumn{2}{l@{}}{\textbf{List-II}} \\ \midrule
\textbf{(P)} & zero-noise extrapolation & \textbf{(1)} & invert a calibrated confusion matrix applied to the measured distribution \\
\textbf{(Q)} & readout-error mitigation & \textbf{(2)} & run the circuit at amplified noise levels and extrapolate the observable to zero noise \\
\textbf{(R)} & dynamical decoupling & \textbf{(3)} & encode logical qubits redundantly, measure syndromes and correct; needs physical errors below a threshold \\
\textbf{(S)} & quantum error correction & \textbf{(4)} & apply pulse sequences during idle periods to refocus slow (quasi-static) dephasing \\
\bottomrule
\end{tabularx}}
\par\smallskip\noindent\begin{tabularx}{\linewidth}{@{}*{2}{L}@{}}
\correctlabel{A}\,P$\to$2, Q$\to$1, R$\to$4, S$\to$3 & \optlabel{B}\,P$\to$1, Q$\to$2, R$\to$4, S$\to$3 \\[3pt]
\optlabel{C}\,P$\to$2, Q$\to$1, R$\to$3, S$\to$4 & \optlabel{D}\,P$\to$2, Q$\to$4, R$\to$1, S$\to$3
\end{tabularx}
\end{qitem}
\begin{solbox}{A}
\textbf{\color{brand}Correct answer.} \textcolor{okgreen}{\textbf{(A)}} P$\to$2, Q$\to$1, R$\to$4, S$\to$3

\textbf{\color{brand}Concepts required.} Mitigation vs correction vs decoupling.

\textbf{\color{brand}Relevant formulae.} ZNE: extrapolate $E(\lambda)$; readout: $A^{-1}\mathbf q$; DD: refocus; QEC: encode + syndromes.

\textbf{\color{brand}Step-by-step explanation.} \begin{itemize}\itemsep0pt
\item ZNE: an \emph{estimator} for expectation values, post-processing (2);
\item readout mitigation: classical post-processing of counts (1);
\item DD: changes the physical evolution during idling (4);
\item QEC: protects the quantum state itself during the computation (3).
\end{itemize}
Only QEC corrects errors on the quantum state; ZNE and readout mitigation correct estimates; noise \emph{simulation} does neither. Hence P$\to$2, Q$\to$1, R$\to$4, S$\to$3.

\textbf{\color{brand}Why the other options are incorrect.}
\begin{itemize}[leftmargin=2.4em, itemsep=1pt, topsep=2pt]
\item[\textbf{(B)}] P$\to$1, Q$\to$2 swaps ZNE and readout mitigation.
\item[\textbf{(C)}] R$\to$3, S$\to$4 swaps dynamical decoupling (pulses during idling) and error correction (encoding and syndromes).
\item[\textbf{(D)}] Q$\to$4, R$\to$1 swaps readout mitigation (classical post-processing) and dynamical decoupling (physical pulses).
\end{itemize}
\textbf{\color{brand}Key takeaway.} Only QEC protects the state itself.
\end{solbox}
\qsep

\needspace{13\baselineskip}
\begin{qitem}{76}
\qmeta{Matching lists}{Error taxonomy}
Match each physical cause (List-I) with its category (List-II).
\par\smallskip\noindent{\renewcommand{\arraystretch}{1.35}%
\begin{tabularx}{\linewidth}{@{}c>{\raggedright\arraybackslash}p{0.33\linewidth}@{\hspace{1.2em}}cL@{}}
\toprule
\multicolumn{2}{@{}l}{\textbf{List-I}} & \multicolumn{2}{l@{}}{\textbf{List-II}} \\ \midrule
\textbf{(P)} & a pulse amplitude 1\,\% too large on every $X$ gate & \textbf{(1)} & decoherence (non-unital, $T_1$) \\
\textbf{(Q)} & a random Pauli error after each gate & \textbf{(2)} & classical measurement (readout) error \\
\textbf{(R)} & energy loss of $\ket1$ to the environment & \textbf{(3)} & coherent error (unitary, systematic) \\
\textbf{(S)} & the discriminator assigns some $\ket1$ shots to outcome 0 & \textbf{(4)} & stochastic (incoherent) gate noise \\
\bottomrule
\end{tabularx}}
\par\smallskip\noindent\begin{tabularx}{\linewidth}{@{}*{2}{L}@{}}
\optlabel{A}\,P$\to$4, Q$\to$3, R$\to$1, S$\to$2 & \correctlabel{B}\,P$\to$3, Q$\to$4, R$\to$1, S$\to$2 \\[3pt]
\optlabel{C}\,P$\to$3, Q$\to$4, R$\to$2, S$\to$1 & \optlabel{D}\,P$\to$3, Q$\to$1, R$\to$4, S$\to$2
\end{tabularx}
\end{qitem}
\begin{solbox}{B}
\textbf{\color{brand}Correct answer.} \textcolor{okgreen}{\textbf{(B)}} P$\to$3, Q$\to$4, R$\to$1, S$\to$2

\textbf{\color{brand}Concepts required.} Error taxonomy.

\textbf{\color{brand}Relevant formulae.} coherent $U_\varepsilon$; stochastic $\sum p_kP_k\cdot P_k$; decoherence $T_1, T_2$; readout $A$.

\textbf{\color{brand}Step-by-step explanation.} \begin{itemize}\itemsep0pt
\item over-rotation: the same unitary error every time (3);
\item random Pauli: a probabilistic mixture --- the notebook's model (4);
\item energy loss: amplitude damping, $T_1$ (1);
\item misassignment: acts on the classical outcome after the quantum measurement (2).
\end{itemize}
Hence P$\to$3, Q$\to$4, R$\to$1, S$\to$2.

\textbf{\color{brand}Why the other options are incorrect.}
\begin{itemize}[leftmargin=2.4em, itemsep=1pt, topsep=2pt]
\item[\textbf{(A)}] P$\to$4, Q$\to$3 swaps coherent and stochastic errors: the miscalibration is the same every time (coherent).
\item[\textbf{(C)}] R$\to$2, S$\to$1 swaps decoherence (quantum) and misassignment (classical readout).
\item[\textbf{(D)}] Q$\to$1, R$\to$4 swaps random Pauli errors (stochastic) and energy loss (amplitude damping, $T_1$).
\end{itemize}
\textbf{\color{brand}Key takeaway.} Different error types need different remedies.
\end{solbox}
\qsep

\needspace{13\baselineskip}
\begin{qitem}{77}
\qmeta{Assertion--Reason}{Cost and nature of ZNE}
\textbf{Assertion (A):} Zero-noise extrapolation increases the number of circuit executions needed to reach a given statistical precision.\\[1mm]
\textbf{Reason (R):} ZNE corrects errors on the quantum state during the computation, in the same way as quantum error correction.
\begin{itemize}[leftmargin=2.6em, labelsep=0.4em, itemsep=2pt, topsep=3pt]
\item[\optlabel{A}] Both Assertion (A) and Reason (R) are true, and R is the correct explanation of A
\item[\optlabel{B}] Both A and R are true, but R is \emph{not} the correct explanation of A
\item[\correctlabel{C}] A is true, but R is false
\item[\optlabel{D}] A is false, but R is true
\end{itemize}
\end{qitem}
\begin{solbox}{C}
\textbf{\color{brand}Correct answer.} \textcolor{okgreen}{\textbf{(C)}} A is true, but R is false

\textbf{\color{brand}Concepts required.} Cost and nature of ZNE.

\textbf{\color{brand}Relevant formulae.} $\operatorname{Var}[E_0] = \sum_k\gamma_k^2\operatorname{Var}[E(\lambda_k)]$.

\textbf{\color{brand}Step-by-step explanation.} \textbf{A is true}: several noise levels must be run, and extrapolation coefficients (e.g.\ $3, -3, 1$) amplify the statistical error, so more shots are needed for the same
error bar. \textbf{R is false}: ZNE never touches the quantum state; it post-processes noisy expectation values and returns a (biased) estimate. Only QEC protects the
state, and it can do so for arbitrary-length computations. Answer: (C).

\textbf{\color{brand}Why the other options are incorrect.}
\begin{itemize}[leftmargin=2.4em, itemsep=1pt, topsep=2pt]
\item[\textbf{(A)}] requires R to be true; ZNE does not act on the state.
\item[\textbf{(B)}] also requires R to be true (with R not explaining A); R is false because ZNE never acts on the state.
\item[\textbf{(D)}] requires A to be false, but extra noise levels and amplified variance always cost extra shots.
\end{itemize}
\textbf{\color{brand}Key takeaway.} Mitigation trades shots for bias reduction; it does not fix the state.
\end{solbox}
\qsep

\needspace{13\baselineskip}
\begin{qitem}{78}
\qmeta{Statement I / II}{Readout and hardware in the notebook}
\textbf{Statement I:} \texttt{Noise\_Simulation.ipynb} includes measurement (readout) errors in its noise model.\\[1mm]
\textbf{Statement II:} The recorded results of \texttt{Noise\_Simulation.ipynb} were obtained on a physical quantum processor.
\begin{itemize}[leftmargin=2.6em, labelsep=0.4em, itemsep=2pt, topsep=3pt]
\item[\optlabel{A}] Both Statement I and Statement II are true
\item[\correctlabel{B}] Both Statement I and Statement II are false
\item[\optlabel{C}] Statement I is true, but Statement II is false
\item[\optlabel{D}] Statement I is false, but Statement II is true
\end{itemize}
\end{qitem}
\begin{solbox}{B}
\textbf{\color{brand}Correct answer.} \textcolor{okgreen}{\textbf{(B)}} Both Statement I and Statement II are false

\textbf{\color{brand}Concepts required.} What the repository implements and where it ran.

\textbf{\color{brand}Relevant formulae.} Noise model: gates \texttt{h, x, ry, rz}, CNOT; target \texttt{stim}; no hardware.

\textbf{\color{brand}Step-by-step explanation.} \textbf{I is false}: \texttt{build\_noise\_model} attaches channels only to \texttt{h}, \texttt{x}, \texttt{ry}, \texttt{rz} and CNOTs; measurements are ideal.
\textbf{II is false}: the recorded run used the \texttt{stim} simulator on a CPU (no GPU visible); the notebook even refuses \texttt{BACKEND = "hardware"}, because noise
models are a simulator feature. Answer: (B).

\textbf{\color{brand}Why the other options are incorrect.}
\begin{itemize}[leftmargin=2.4em, itemsep=1pt, topsep=2pt]
\item[\textbf{(A)}] requires both statements to be true; neither is.
\item[\textbf{(C)}] requires Statement I to be true, but the noise model has no measurement channel.
\item[\textbf{(D)}] requires Statement II to be true, but the recorded run used the \texttt{stim} simulator.
\end{itemize}
\textbf{\color{brand}Key takeaway.} Simulated noise $\neq$ hardware results.
\end{solbox}
\qsep

\needspace{20\baselineskip}
\begin{passagebox}{Paragraph for Questions 79 and 80}
For the lecture (illustrative experiment, not part of the notebook), zero-noise extrapolation (ZNE) was applied to the parity
$\langle X^{\otimes4}\rangle$ of a 4-qubit GHZ state. The circuit was simulated exactly on \texttt{density-matrix-cpu} with the notebook's
depolarizing noise model at the scaled probability $p = 0.01\lambda$ (noise scale factor $\lambda$). The ideal value is $1$.
\begin{center}\small
\begin{tabular}{cccc}\toprule $\lambda$ & 1 & 2 & 3\\ \midrule $\langle X^{\otimes4}\rangle$ & 0.9055 & 0.8189 & 0.7396\\ \bottomrule\end{tabular}
\end{center}
\end{passagebox}
\needspace{13\baselineskip}
\begin{qitem}{79}
\qmeta{Paragraph based}{ZNE extrapolation values}
The linear extrapolation from $\lambda = 1, 2$ and the Richardson (quadratic) extrapolation from $\lambda = 1, 2, 3$ give
\par\smallskip\noindent\begin{tabularx}{\linewidth}{@{}*{2}{L}@{}}
\correctlabel{A}\,$0.9921$ and $0.9994$ & \optlabel{B}\,$0.9055$ and $0.9921$ \\[3pt]
\optlabel{C}\,$0.9921$ and $1.0021$ & \optlabel{D}\,$0.9872$ and $0.9994$
\end{tabularx}
\end{qitem}
\begin{solbox}{A}
\textbf{\color{brand}Correct answer.} \textcolor{okgreen}{\textbf{(A)}} $0.9921$ and $0.9994$

\textbf{\color{brand}Concepts required.} Applying linear and Richardson extrapolation to data.

\textbf{\color{brand}Relevant formulae.} $2E_1 - E_2$; $3E_1 - 3E_2 + E_3$.

\textbf{\color{brand}Step-by-step explanation.} Linear: $2E(1) - E(2) = 1.8110 - 0.8189 = 0.9921$. Richardson: $3E(1) - 3E(2) + E(3) = 2.7165 - 2.4567 + 0.7396 = 0.9994$.
Both reduce the error of the raw value $0.9055$ ($9.5\,\%$) to below $1\,\%$. $1.0021$ is the exponential fit, which overshoots the physical maximum.

\textbf{\color{brand}Why the other options are incorrect.}
\begin{itemize}[leftmargin=2.4em, itemsep=1pt, topsep=2pt]
\item[\textbf{(B)}] $0.9055$ is the raw value at $\lambda = 1$, not an extrapolation; $0.9921$ is the linear (not the Richardson) estimate.
\item[\textbf{(C)}] $1.0021$ is the exponential fit, which overshoots the physical maximum of a parity.
\item[\textbf{(D)}] $0.9872$ is the least-squares straight line through all three points, a different (and here worse) linear estimator than the two-point formula.
\end{itemize}
\textbf{\color{brand}Key takeaway.} Extrapolation removes most of the bias in this example.
\end{solbox}
\qsep

\needspace{13\baselineskip}
\begin{qitem}{80}
\qmeta{Paragraph based}{Interpreting ZNE}
Which conclusion about this ZNE experiment is correct?
\begin{itemize}[leftmargin=2.6em, labelsep=0.4em, itemsep=2pt, topsep=3pt]
\item[\optlabel{A}] on hardware the same result is obtained without scaling the noise, since $\lambda$ is a software parameter
\item[\optlabel{B}] ZNE has corrected the GHZ state itself, so a subsequent circuit could use the improved state
\item[\correctlabel{C}] the extrapolated values depend on the fitted model (an exponential fit gives $1.0021 > 1$, unphysical) and, with shot noise, would carry larger error bars than the raw value
\item[\optlabel{D}] because the exact density matrix was used, ZNE is no longer needed: the raw value $0.9055$ is already the ideal one
\end{itemize}
\end{qitem}
\begin{solbox}{C}
\textbf{\color{brand}Correct answer.} \textcolor{okgreen}{\textbf{(C)}} the extrapolated values depend on the fitted model (an exponential fit gives $1.0021 > 1$, unphysical) and, with shot noise, would carry larger error bars than the raw value

\textbf{\color{brand}Concepts required.} Interpreting ZNE results honestly.

\textbf{\color{brand}Relevant formulae.} Bias depends on the model; variance amplification $\sqrt{\sum\gamma_k^2}$.

\textbf{\color{brand}Step-by-step explanation.} ZNE estimates an expectation value; the state remains noisy. Different extrapolation models give $0.9921$, $0.9994$, $1.0021$ --- a model-dependent
bias, and $1.0021$ exceeds the allowed range of a parity. Here the inputs were exact; with $S$ shots per point, the Richardson combination would multiply the
statistical error by about $\sqrt{19} \approx 4.4$. On hardware the noise cannot be scaled by changing $p$; it is amplified physically, e.g.\ by gate folding or pulse
stretching.

\textbf{\color{brand}Why the other options are incorrect.}
\begin{itemize}[leftmargin=2.4em, itemsep=1pt, topsep=2pt]
\item[\textbf{(A)}] On hardware $p$ cannot be set in software; noise must be amplified physically (gate folding, pulse stretching).
\item[\textbf{(B)}] ZNE post-processes expectation values; the GHZ state itself stays noisy and cannot be reused in a better form.
\item[\textbf{(D)}] The exact density matrix is still the \emph{noisy} state; $0.9055 \neq 1$ shows the remaining bias.
\end{itemize}
\textbf{\color{brand}Key takeaway.} ZNE is a model-dependent estimator, not error correction.
\end{solbox}
\qsep

\end{document}